% Lutte contre l'effet de serre — SVT 1ère TI — Cours signé OnBuch+
% Programme officiel MINESEC. Compiler avec : tectonic NCT12-lutte-effet-serre-c-ti.tex
\newcommand{\DOCMATIERE}{SVT}
\newcommand{\DOCNIVEAU}{1ère TI}
\newcommand{\DOCMODULE}{Module 3 : L'Éducation à l'Environnement et au Développement Durable}
\newcommand{\DOCLECON}{NCT12}
\newcommand{\DOCTITRE}{Lutte contre l'effet de serre}
% OnBuch+ — préambule « cours signés » ENRICHI (compilé avec tectonic / XeLaTeX)
% Système visuel « Pop » d'OnBuch+ : fond crème (jamais blanc), bordures encre
% épaisses, ombres « dures » décalées, titres Archivo Black en capitales.
% Version MATHÉMATIQUES : graphes (pgfplots/tikz), boîtes pédagogiques
% (définition, propriété, méthode, attention, le savais-tu, exercice résolu,
% exercice, TP, bilan), page de garde et signature OnBuch+ ancrée en bas.
%
% Macros attendues dans main.tex AVANT \input{preamble} :
%   \DOCMATIERE  \DOCNIVEAU  \DOCMODULE  \DOCLECON (numéro)  \DOCTITRE
\documentclass[11pt]{article}

\usepackage{fontspec}
\usepackage[french]{babel}
\usepackage[a4paper,margin=2cm,top=3.4cm,bottom=2.8cm,headheight=1.4cm,footskip=1.3cm]{geometry}
\usepackage{amsmath,amssymb,amsfonts}
\usepackage{mathtools} % \xmapsto, \coloneqq... souvent utilisés par les modèles
\usepackage{cancel} % simplifications barrées (bilans de forces)
% \abs{z} (module, valeur absolue) et \norm{u} (norme), utilisés par les modèles
\providecommand{\abs}{}\let\abs\relax\DeclarePairedDelimiter{\abs}{\lvert}{\rvert}
\providecommand{\norm}{}\let\norm\relax\DeclarePairedDelimiter{\norm}{\lVert}{\rVert}
\usepackage[table]{xcolor} % [table] : fournit aussi \rowcolors (alternance de lignes)
\usepackage{pagecolor}
\usepackage{tikz}
\usetikzlibrary{arrows.meta,positioning,calc,shapes.geometric,shapes.misc,shapes.arrows,decorations.pathmorphing,decorations.pathreplacing,decorations.markings,patterns,shadows,mindmap,trees,fit,backgrounds,matrix,intersections,angles,quotes}
\usepackage{pgfplots}
\usepackage[version=4]{mhchem} % équations nucléaires \ce{^{235}_{92}U}
\usepackage[european,siunitx]{circuitikz} % schémas électriques
\pgfplotsset{compat=1.18}
\usepgfplotslibrary{fillbetween,groupplots} % aires entre courbes (calcul intégral)
\usepackage{enumitem}
\usepackage{fancyhdr}
% [most] charge toutes les bibliothèques tcolorbox (listings, minted, raster,
% documentation...) et gonfle inutilement la mémoire TeX sur un document long
% (~300 boîtes) : on ne charge que ce qui est réellement utilisé.
\usepackage[skins,breakable]{tcolorbox}
\usepackage{graphicx}
\usepackage{colortbl}
\usepackage{booktabs}
\usepackage{tabularx}
\usepackage{diagbox} % case d'en-tête diagonale des tableaux à double entrée
% Colonnes extensibles centrée / gauche / droite (C, L, R), souvent utilisées par les modèles
\newcolumntype{C}{>{\centering\arraybackslash}X}
\newcolumntype{L}{>{\raggedright\arraybackslash}X}
\newcolumntype{R}{>{\raggedleft\arraybackslash}X}
\usepackage{array}
\usepackage{multicol}
\usepackage{siunitx}
% parse-numbers=false : accepte aussi \SI{2,5 \times 10^{-3}}{...}
\sisetup{locale=FR,output-decimal-marker={,},group-minimum-digits=4,parse-numbers=false}
\DeclareSIUnit\ohm{\ensuremath{\symup{\Omega}}} % U+2126 absent de la police
\DeclareSIUnit\division{div} % réglages d'oscilloscope (ms/div, V/div)
\DeclareSIUnit\tr{tr} % tours (vitesse de rotation, ex. \SI{1500}{\tr\per\minute})
\DeclareSIUnit\ch{ch} % cheval-vapeur (puissance moteur, unité usuelle hors SI)
\DeclareSIUnit\franccfa{FCFA} % franc CFA (coûts, contexte camerounais)
\DeclareSIUnit\dioptre{\ensuremath{\delta}} % vergence d'une lentille (dioptrie)
\usepackage{qrcode}
% \degree : commande fréquemment utilisée par les modèles pour un angle ou une
% température en dehors de \SI{}{} (non fournie par les paquets chargés).
\providecommand{\degree}{^{\circ}}
% \qed (fin de démonstration), utilisé par les modèles sans amsthm
\providecommand{\qed}{\hfill$\square$}
% \barre : double barre des valeurs interdites dans un tableau de variations (tkz-tab)
\providecommand{\barre}{\ensuremath{\|}}
% environnement proof (amsthm non chargé) utilisé par les modèles
\ifdefined\proof\else\newenvironment{proof}[1][Démonstration]{\par\noindent\textit{#1.}\ }{\qed\par}\fi
\usepackage{lastpage}
\usepackage{hyperref}
\hypersetup{hidelinks}

% ============================================================
% Palette « Pop » OnBuch+ — mêmes hex que la classe P (plus_ui.dart)
% ============================================================
\definecolor{poporange}{HTML}{F59321}
\definecolor{popdark}{HTML}{E07A0C}
\definecolor{popcream}{HTML}{FFF6E8}
\definecolor{popink}{HTML}{1C1714}
\definecolor{poppurple}{HTML}{7A5AE0}
\definecolor{popgreen}{HTML}{2E9E6A}
\definecolor{poppink}{HTML}{E0457B}
\definecolor{popblue}{HTML}{2F7DE1}
\definecolor{popgold}{HTML}{C98A12}
\definecolor{popmuted}{HTML}{8A7F76}
\definecolor{popline}{HTML}{EADFD2}
\definecolor{popgray}{HTML}{8A7F76}
\colorlet{popgrey}{popgray}
% Alias tolérants (le modèle nomme parfois les couleurs différemment)
\colorlet{popwhite}{white}
\colorlet{popblack}{popink}
\colorlet{popred}{poppink}
\colorlet{popyellow}{popgold}
% Teintes claires pour fonds de tableaux / zones
\colorlet{popblueL}{popblue!10!white}
\colorlet{poporangeL}{poporange!14!white}
\colorlet{popgreenL}{popgreen!12!white}
\colorlet{poppurpleL}{poppurple!12!white}
\colorlet{poppinkL}{poppink!12!white}
\colorlet{popgoldL}{popgold!14!white}

\pagecolor{popcream}

% ============================================================
% Typographie
% ============================================================
\defaultfontfeatures{Ligatures=TeX}
\setmainfont{PlusJakartaSans-Medium.ttf}[Path=../../../fonts/,BoldFont=PlusJakartaSans-Bold.ttf]
% Maths en sans-serif (Fira Math) avec chiffres et lettres droites de Plus
% Jakarta, pour que les formules se fondent dans le texte.
\usepackage[math-style=ISO,bold-style=ISO]{unicode-math}
\setmathfont{FiraMath-Regular.otf}[Path=../../../fonts/]
\setmathfont{PlusJakartaSans-Medium.ttf}[Path=../../../fonts/,range={up/num,up/{Latin,latin},\mathup}]
\usepackage{icomma} % virgule décimale sans espace en mode math
% Caractères mathématiques Unicode absents des polices de texte (Plus Jakarta,
% Archivo Black) : rendus par leur équivalent mathématique, en texte comme en
% formule — sinon ils s'affichent en carré vide (ex. « Arithmétique dans ℤ »).
\usepackage{newunicodechar}
\newunicodechar{ℝ}{\ensuremath{\mathbb{R}}}
\newunicodechar{ℤ}{\ensuremath{\mathbb{Z}}}
\newunicodechar{ℂ}{\ensuremath{\mathbb{C}}}
\newunicodechar{ℕ}{\ensuremath{\mathbb{N}}}
\newunicodechar{ℚ}{\ensuremath{\mathbb{Q}}}
\newunicodechar{𝔻}{\ensuremath{\mathbb{D}}}
\newunicodechar{α}{\ensuremath{\alpha}}
\newunicodechar{β}{\ensuremath{\beta}}
\newunicodechar{γ}{\ensuremath{\gamma}}
\newunicodechar{δ}{\ensuremath{\delta}}
\newunicodechar{ε}{\ensuremath{\varepsilon}}
\newunicodechar{θ}{\ensuremath{\theta}}
\newunicodechar{λ}{\ensuremath{\lambda}}
\newunicodechar{μ}{\ensuremath{\mu}}
\newunicodechar{σ}{\ensuremath{\sigma}}
\newunicodechar{τ}{\ensuremath{\tau}}
\newunicodechar{φ}{\ensuremath{\varphi}}
\newunicodechar{ψ}{\ensuremath{\psi}}
\newunicodechar{ω}{\ensuremath{\omega}}
\newunicodechar{Σ}{\ensuremath{\Sigma}}
% majuscules produites par \MakeUppercase dans les titres (α-aminés -> Α)
\newunicodechar{Α}{\ensuremath{\alpha}}
\newunicodechar{Β}{\ensuremath{\beta}}
\newunicodechar{Γ}{\ensuremath{\Gamma}}
\newunicodechar{Δ}{\ensuremath{\Delta}}
\newunicodechar{Ε}{\ensuremath{\varepsilon}}
\newunicodechar{Θ}{\ensuremath{\Theta}}
\newunicodechar{Λ}{\ensuremath{\Lambda}}
\newunicodechar{Μ}{\ensuremath{\mu}}
\newunicodechar{Τ}{\ensuremath{\tau}}
\newunicodechar{Φ}{\ensuremath{\Phi}}
\newunicodechar{Ψ}{\ensuremath{\Psi}}
\newunicodechar{Ω}{\ensuremath{\Omega}}
\newunicodechar{↦}{\ensuremath{\mapsto}}
\newunicodechar{∈}{\ensuremath{\in}}
\newunicodechar{∉}{\ensuremath{\notin}}
\newunicodechar{∧}{\ensuremath{\wedge}}
\newunicodechar{⇔}{\ensuremath{\Leftrightarrow}}
\newunicodechar{⟺}{\ensuremath{\Longleftrightarrow}}
\newunicodechar{⇒}{\ensuremath{\Rightarrow}}
\newunicodechar{⟹}{\ensuremath{\Longrightarrow}}
\newunicodechar{∘}{\ensuremath{\circ}}
\newunicodechar{∪}{\ensuremath{\cup}}
\newunicodechar{∩}{\ensuremath{\cap}}
\newunicodechar{≡}{\ensuremath{\equiv}}
\newunicodechar{⊂}{\ensuremath{\subset}}
\newunicodechar{⊥}{\ensuremath{\perp}}
\newunicodechar{∃}{\ensuremath{\exists}}
\newunicodechar{∀}{\ensuremath{\forall}}
\newunicodechar{∛}{\ensuremath{\sqrt[3]{\,}}}
\newunicodechar{∜}{\ensuremath{\sqrt[4]{\,}}}
\newunicodechar{⃗}{\ensuremath{{}^{\to}}}
\newunicodechar{ⁿ}{\ifmmode^{n}\else\textsuperscript{\ensuremath{n}}\fi}
\newunicodechar{ˣ}{\ifmmode^{x}\else\textsuperscript{\ensuremath{x}}\fi}
\newunicodechar{ᵇ}{\ifmmode^{b}\else\textsuperscript{\ensuremath{b}}\fi}
\newunicodechar{ʳ}{\ifmmode^{r}\else\textsuperscript{\ensuremath{r}}\fi}
\newunicodechar{ᵘ}{\ifmmode^{u}\else\textsuperscript{\ensuremath{u}}\fi}
\newunicodechar{ʸ}{\ifmmode^{y}\else\textsuperscript{\ensuremath{y}}\fi}
\newunicodechar{ᵃ}{\ifmmode^{a}\else\textsuperscript{\ensuremath{a}}\fi}
\newunicodechar{ᵉ}{\ifmmode^{e}\else\textsuperscript{\ensuremath{e}}\fi}
\newunicodechar{ᵏ}{\ifmmode^{k}\else\textsuperscript{\ensuremath{k}}\fi}
\newunicodechar{ᵖ}{\ifmmode^{p}\else\textsuperscript{\ensuremath{p}}\fi}
\newunicodechar{ᴺ}{\ifmmode^{N}\else\textsuperscript{\ensuremath{N}}\fi}
\newunicodechar{ᶜ}{\ifmmode^{c}\else\textsuperscript{\ensuremath{c}}\fi}
\newunicodechar{ʰ}{\ifmmode^{h}\else\textsuperscript{\ensuremath{h}}\fi}
\newunicodechar{ᵠ}{\ifmmode^{q}\else\textsuperscript{\ensuremath{q}}\fi}
\newunicodechar{ₙ}{\ifmmode_{n}\else\textsubscript{\ensuremath{n}}\fi}
\newunicodechar{ₐ}{\ifmmode_{a}\else\textsubscript{\ensuremath{a}}\fi}
\newunicodechar{ᵢ}{\ifmmode_{i}\else\textsubscript{\ensuremath{i}}\fi}
\newunicodechar{ᵣ}{\ifmmode_{r}\else\textsubscript{\ensuremath{r}}\fi}
\newunicodechar{ₖ}{\ifmmode_{k}\else\textsubscript{\ensuremath{k}}\fi}
\newunicodechar{ᵦ}{\ifmmode_{\beta}\else\textsubscript{\ensuremath{\beta}}\fi}
\newunicodechar{ₓ}{\ifmmode_{x}\else\textsubscript{\ensuremath{x}}\fi}
\newfontfamily\popdisplay{ArchivoBlack-Regular.ttf}[Path=../../../fonts/]
\newfontfamily\poplogo{SpaceGrotesk-Bold.ttf}[Path=../../../fonts/]
\newfontfamily\popsemi{PlusJakartaSans-SemiBold.ttf}[Path=../../../fonts/]
\newfontfamily\popxbold{PlusJakartaSans-ExtraBold.ttf}[Path=../../../fonts/]
\newfontfamily\popxboldls{PlusJakartaSans-ExtraBold.ttf}[Path=../../../fonts/,LetterSpace=3.0]

\setlist[enumerate]{leftmargin=1.7em,itemsep=5pt,topsep=4pt}
\setlist[itemize]{leftmargin=1.7em,itemsep=4pt,topsep=4pt,label={\color{poporange}\textbullet}}
\setlength{\parindent}{0pt}
\setlength{\parskip}{5pt}
\renewcommand{\arraystretch}{1.25}

% pgfplots : style Pop par défaut
\pgfplotsset{
  every axis/.append style={axis line style={popink,line width=1pt},tick style={popink},
    grid style={popline,line width=0.5pt},label style={font=\small},tick label style={font=\footnotesize}},
  popcurve/.style={poporange,line width=1.8pt,no markers,smooth},
}
% Même style disponible dans un \draw TikZ simple (hors pgfplots)
\tikzset{popcurve/.style={poporange,line width=1.8pt}}
\tikzset{popgrid/.style={popline,very thin}}
\colorlet{popcurve}{poporange} % le nom du style est parfois utilisé comme couleur

% ============================================================
% Pill / squiggle
% ============================================================
\newcommand{\poppill}[3]{%
  \tikz[baseline=-0.55ex]\node[draw=popink,line width=1.2pt,rounded corners=2.6pt,%
    fill=#2,inner xsep=6.5pt,inner ysep=3pt]{%
    \popxboldls\fontsize{7}{7}\selectfont\color{#3}\MakeUppercase{#1}};%
}
\newcommand{\popsquiggle}[1]{%
  \tikz[baseline=-0.4ex]{
    \draw[#1,line width=2.6pt,line cap=round]
      plot [smooth] coordinates {(0,0.09) (0.34,0.01) (0.68,0.21) (1.02,0.01) (1.36,0.10)};
  }%
}
% Mot-clé surligné (marqueur orange sous le texte)
\newcommand{\cle}[1]{{\popxbold\color{popdark}#1}}

% ============================================================
% En-tête / pied
% ============================================================
\pagestyle{fancy}
\fancyhf{}
\renewcommand{\headrulewidth}{0pt}
\renewcommand{\footrulewidth}{0pt}
\lhead{\raisebox{-0.15ex}{\poplogo\fontsize{13}{13}\selectfont\color{popink}OnBuch\color{poporange}+}}
\rhead{\poppill{\DOCMATIERE\ \textbullet\ \DOCNIVEAU\ \textbullet\ Leçon \DOCLECON}{white}{popink}}
\lfoot{\popxbold\fontsize{7.6}{7.6}\selectfont\color{popmuted}\MakeUppercase{Cours signé OnBuch+}\ \textbullet\ {\popsemi\DOCTITRE}}
\rfoot{\tikz[baseline=-0.6ex]\node[rounded corners=8.5pt,fill=popink,minimum height=17pt,minimum width=17pt,inner xsep=5pt,inner sep=0]{\popxbold\fontsize{8}{8}\selectfont\color{popcream}\thepage\,/\,\pageref*{LastPage}};}

% ============================================================
% Titres
% ============================================================
\newcommand{\chaptitre}[2]{%
  \begin{tcolorbox}[enhanced,colback=poporange,colframe=popink,boxrule=2.6pt,arc=11pt,
      drop shadow=popink,left=15pt,right=15pt,top=12pt,bottom=11pt,before skip=3pt,after skip=18pt]
    \poppill{#2}{popink}{popcream}\par\vspace{9pt}
    {\raggedright\hyphenpenalty=10000\exhyphenpenalty=10000\popdisplay\fontsize{22}{24}\selectfont\color{popcream}\MakeUppercase{#1}\par}
  \end{tcolorbox}
}
% Section numérotée : pastille orange avec le numéro + titre Archivo
\newcounter{popsec}
\newcommand{\coursec}[1]{%
  \stepcounter{popsec}\setcounter{popsub}{0}%
  \par\vspace{16pt}\noindent
  \begin{minipage}{\linewidth}
  \tikz[baseline=-0.7ex]\node[draw=popink,line width=1.6pt,fill=poporange,rounded corners=4pt,minimum size=22pt,inner sep=0]{\popdisplay\fontsize{12}{12}\selectfont\color{popcream}\thepopsec};%
  \hspace{8pt}{\popdisplay\fontsize{15}{17}\selectfont\color{popink}\MakeUppercase{#1}}\par
  \vspace{4pt}\noindent{\color{poporange}\rule{36pt}{3.6pt}}
  \end{minipage}\par\nopagebreak\vspace{8pt}
}
\newcounter{popsub}
\newcommand{\courssub}[1]{%
  \stepcounter{popsub}%
  \par\vspace{9pt}\noindent{\popxbold\fontsize{11.5}{13}\selectfont\color{popink}{\color{poporange}\thepopsec.\thepopsub}\hspace{6pt}#1}\par\nopagebreak\vspace{4pt}
}

% ============================================================
% Boîtes pédagogiques — même recette : fond blanc, bordure épaisse colorée,
% ombre dure encre, badge plein. Titre optionnel [#1].
% ============================================================
\tcbset{popbase/.style={breakable,enhanced,colback=white,boxrule=2.3pt,arc=9pt,
  shadow={3pt}{-3pt}{0pt}{popink},left=14pt,right=14pt,top=11pt,bottom=11pt,before skip=11pt,after skip=15pt,
  fontupper=\popsemi\fontsize{10.6}{14.5}\selectfont\color{popink},
  fontlower=\popsemi\fontsize{10.6}{14.5}\selectfont\color{popink}}}
% Coche dessinée (le glyphe ✓ n'existe pas dans les polices du document)
\newcommand{\popcheck}[1]{\tikz[baseline=-0.5ex]\draw[#1,line width=1.6pt,line cap=round,line join=round](-0.13,0.0)--(-0.04,-0.09)--(0.13,0.1);}
\providecommand{\coche}{\popcheck{popgreen}}
\providecommand{\resultat}[1]{\par\smallskip\noindent\fcolorbox{popgreen}{popgreenL}{\parbox{\dimexpr\linewidth-2\fboxsep-2\fboxrule\relax}{\textbf{Résultat.} #1}}\par\smallskip}
\newcommand{\popboxhead}[3]{\poppill{#1}{#2}{white}\ifx\relax#3\relax\else\hspace{6pt}{\popxbold\fontsize{10.6}{12}\selectfont\color{popink}#3}\fi\par\vspace{7pt}}

\newtcolorbox{aretenir}[1][]{popbase,colframe=popgreen,before upper={\popboxhead{À retenir}{popgreen}{#1}}}
\newtcolorbox{exemplebox}[1][]{popbase,colframe=poporange,before upper={\popboxhead{Exemple}{poporange}{#1}}}
\newtcolorbox{definition}[1][]{popbase,colframe=popblue,before upper={\popboxhead{Définition}{popblue}{#1}}}
\newtcolorbox{propriete}[1][]{popbase,colframe=poppurple,before upper={\popboxhead{Propriété}{poppurple}{#1}}}
\newtcolorbox{methode}[1][]{popbase,colframe=popgold,before upper={\popboxhead{Méthode}{popgold}{#1}}}
\newtcolorbox{attention}[1][]{popbase,colframe=poppink,before upper={\popboxhead{Attention}{poppink}{#1}}}
\newtcolorbox{experience}[1][]{popbase,colframe=popblue,colback=popblueL,before upper={\popboxhead{Expérience}{popblue}{#1}}}
% « Le savais-tu ? » : carte violette pleine (contexte, culture, Cameroun)
\newtcolorbox{savaistu}[1][]{popbase,colback=poppurple,colframe=popink,
  fontupper=\popsemi\fontsize{10.4}{14.2}\selectfont\color{white},
  before upper={\poppill{Le savais-tu\,?}{white}{poppurple}\ifx\relax#1\relax\else\hspace{6pt}{\popxbold\color{white}#1}\fi\par\vspace{7pt}}}
% Exercice résolu : énoncé en haut, \tcblower puis la solution sur fond crème
\newcounter{popexo}\newcounter{popexr}
\newtcolorbox{exoresolu}[1][]{popbase,colframe=poporange,colbacklower=popcream,
  segmentation style={popink,line width=1.2pt,dashed},
  before upper={\stepcounter{popexr}\popboxhead{Exercice résolu \thepopexr}{poporange}{#1}},
  before lower={\poppill{Solution}{popink}{popcream}\par\vspace{6pt}}}
% Exercice (non résolu en place) — la correction est dans la partie Corrigés
\newtcolorbox{exercice}[1][]{popbase,colframe=popink,
  before upper={\stepcounter{popexo}\popboxhead{Exercice \thepopexo}{popink}{#1}}}
% Corrigé
\newtcolorbox{corrige}[1][]{popbase,colframe=popgreen,colback=popgreenL,
  before upper={\popboxhead{Corrigé}{popgreen}{#1}}}
% Objectifs / prérequis / bilan
\newtcolorbox{objectifs}{popbase,colframe=popink,colback=poporangeL,
  before upper={\popboxhead{Objectifs de la leçon}{popink}{}}}
\newtcolorbox{prerequis}{popbase,colframe=popink,colback=popblueL,
  before upper={\popboxhead{Prérequis}{popblue}{}}}
\newtcolorbox{bilan}{popbase,colframe=popink,colback=white,boxrule=2.8pt,
  before upper={\popboxhead{Fiche bilan}{poporange}{L'essentiel à connaître pour l'examen}}}
% Figure encadrée avec légende
\newtcolorbox{popfigure}[1][]{enhanced,breakable=false,colback=white,colframe=popline,boxrule=1.4pt,arc=8pt,
  left=10pt,right=10pt,top=10pt,bottom=8pt,before skip=10pt,after skip=12pt,halign=center,
  lower separated=false,halign lower=center,
  fontlower=\popsemi\fontsize{9}{11.5}\selectfont\color{popmuted},
  #1}
\newcommand{\legende}[1]{\par\vspace{6pt}{\popsemi\fontsize{9}{11.5}\selectfont\color{popmuted}#1}}

% ============================================================
% Signature OnBuch+
% ============================================================
\newcommand{\poplogoinline}[1]{{\poplogo\fontsize{#1}{#1}\selectfont\color{popink}OnBuch\color{poporange}+}}
\newcommand{\signatureonbuch}{%
  \begin{tcolorbox}[enhanced,colback=popink,colframe=popink,boxrule=2.2pt,arc=10pt,
      drop shadow=poporange,left=14pt,right=14pt,top=10pt,bottom=10pt,before skip=0pt,after skip=0pt]
    \tikz[baseline=-0.7ex]\node[circle,fill=popgreen,draw=popcream,line width=1.3pt,minimum size=22pt,inner sep=0]{\popcheck{white}};%
    \hspace{9pt}\begin{minipage}[c]{0.62\linewidth}
      {\popxbold\fontsize{12}{13}\selectfont\color{popcream}Signé\ \textbullet\ {\poplogo OnBuch\color{poporange}+}}\par\vspace{2pt}
      {\raggedright\popsemi\fontsize{8}{10.5}\selectfont\color{popline}\MakeUppercase{Équipe pédagogique OnBuch+}\\\MakeUppercase{Conforme au programme officiel MINESEC}\par}
    \end{minipage}\hfill
    {\poplogo\fontsize{22}{22}\selectfont\color{popcream}OnBuch\color{poporange}+}
  \end{tcolorbox}%
}

% ============================================================
% Page de garde
% #1 titre · #2 matière · #3 niveau · #4 module · #5 n° leçon · #6 durée ·
% #7 description / objectifs (texte ou itemize)
% ============================================================
\newcommand{\pagegarde}[7]{%
  \thispagestyle{empty}
  \newgeometry{a4paper,margin=1.7cm,top=1.6cm,bottom=1.4cm}%
  \begin{tikzpicture}[remember picture,overlay]
    \fill[poporange!18] ([xshift=-3.2cm,yshift=-3.6cm]current page.north east) circle (4.6cm);
    \fill[poppurple!14] ([xshift=2.4cm,yshift=7.6cm]current page.south west) circle (3.8cm);
  \end{tikzpicture}%
  {\poplogo\fontsize{30}{30}\selectfont\color{popink}OnBuch\color{poporange}+}\hfill
  \raisebox{0.6ex}{\poppill{Cours signé \textbullet\ Premium}{popink}{popcream}}\par
  \vspace{1pt}\popsquiggle{poppurple}\par
  \vspace{8pt}
  \begin{tcolorbox}[enhanced,colback=poporange,colframe=popink,boxrule=2.8pt,arc=13pt,
      drop shadow=popink,left=18pt,right=18pt,top=12pt,bottom=13pt,after skip=10pt]
    \poppill{#2 \textbullet\ #3}{popink}{popcream}\hspace{5pt}\poppill{#4}{white}{popink}\par\vspace{8pt}
    {\popdisplay\fontsize{14}{14}\selectfont\color{popink}LEÇON #5}\par\vspace{4pt}
    {\raggedright\hyphenpenalty=10000\exhyphenpenalty=10000\popdisplay\fontsize{25}{27}\selectfont\color{popcream}\MakeUppercase{#1}\par}
    \vspace{8pt}
    {\popxbold\fontsize{9.5}{9.5}\selectfont\color{popink}\MakeUppercase{Durée conseillée : #6}}
  \end{tcolorbox}
  \begin{tcolorbox}[enhanced,colback=white,colframe=popink,boxrule=2.2pt,arc=10pt,
      drop shadow=popink,left=16pt,right=16pt,top=10pt,bottom=10pt,after skip=8pt]
    \poppill{Dans cette leçon}{poppurple}{white}\par\vspace{5pt}
    \popsemi\fontsize{9.3}{12.6}\selectfont\color{popink}#7
  \end{tcolorbox}
  \noindent\begin{minipage}[t]{0.487\linewidth}
    \begin{tcolorbox}[enhanced,colback=popgreen,colframe=popink,boxrule=2.2pt,arc=10pt,
        drop shadow=popink,left=11pt,right=11pt,top=10pt,bottom=10pt,height=3.1cm,valign=center]
      \tcbox[colback=white,colframe=popink,boxrule=1.3pt,arc=5pt,boxsep=0pt,nobeforeafter,
        left=3pt,right=3pt,top=3pt,bottom=3pt,valign=center]{\qrcode[height=1.9cm]{https://play.google.com/store/apps/details?id=com.luvvix.onbuch&hl=fr}}%
      \hspace{8pt}\begin{minipage}[c]{0.5\linewidth}
        {\popdisplay\fontsize{11}{12.5}\selectfont\color{white}\MakeUppercase{Télécharge OnBuch}}\par\vspace{4pt}
        {\raggedright\popsemi\fontsize{8.4}{11}\selectfont\color{white}Gratuit \textbullet\ Google Play\\Tuteur IA, annales, concours, résultats\par}
      \end{minipage}
    \end{tcolorbox}
  \end{minipage}\hfill
  \begin{minipage}[t]{0.487\linewidth}
    \begin{tcolorbox}[enhanced,colback=poppurple,colframe=popink,boxrule=2.2pt,arc=10pt,
        drop shadow=popink,left=11pt,right=11pt,top=10pt,bottom=10pt,height=3.1cm,valign=center]
      \tikz[baseline=-0.7ex]\node[circle,fill=white,draw=popink,line width=1.3pt,minimum size=1.6cm,inner sep=0]{\popdisplay\fontsize{24}{24}\selectfont\color{poppurple}+};%
      \hspace{8pt}\begin{minipage}[c]{0.6\linewidth}
        {\popdisplay\fontsize{11}{12.5}\selectfont\color{white}\MakeUppercase{Passe à OnBuch+}}\par\vspace{4pt}
        {\raggedright\popsemi\fontsize{8.4}{11}\selectfont\color{white}Tous les cours signés, Tuteur IA illimité, fiches PDF — onglet OnBuch+ de l'app\par}
      \end{minipage}
    \end{tcolorbox}
  \end{minipage}
  \vspace{8pt}
  \popcontenu
  \vfill
  \signatureonbuch
  \restoregeometry
  \newpage
}

% Rangée « Ce cours contient » (page de garde)
\newcommand{\poptile}[3]{% #1 couleur #2 gros texte #3 légende
  \begin{tcolorbox}[enhanced,colback=#1,colframe=popink,boxrule=2pt,arc=9pt,shadow={2.5pt}{-2.5pt}{0pt}{popink},
      left=6pt,right=6pt,top=9pt,bottom=9pt,halign=center,valign=center,height=1.75cm,width=\linewidth,nobeforeafter]
    {\popdisplay\fontsize{12.5}{13}\selectfont\color{white}\MakeUppercase{#2}}\par\vspace{3pt}
    {\centering\popsemi\fontsize{7.8}{9.5}\selectfont\color{white}#3\par}
  \end{tcolorbox}}
\newcommand{\popcontenu}{%
  \par\noindent{\popxboldls\fontsize{7.5}{7.5}\selectfont\color{popmuted}CE COURS SIGNÉ CONTIENT}\par\vspace{7pt}
  \noindent\begin{minipage}[t]{0.235\linewidth}\poptile{popblue}{Cours}{complet, illustré, pas à pas}\end{minipage}\hfill
  \begin{minipage}[t]{0.235\linewidth}\poptile{poporange}{Méthodes}{et exercices résolus}\end{minipage}\hfill
  \begin{minipage}[t]{0.235\linewidth}\poptile{poppink}{Exercices}{type examen + corrigés}\end{minipage}\hfill
  \begin{minipage}[t]{0.235\linewidth}\poptile{popgreen}{Fiche bilan}{l'essentiel à retenir}\end{minipage}\par}

% Sommaire
\newtcolorbox{sommaire}{enhanced,colback=white,colframe=popink,boxrule=2.2pt,arc=10pt,
  drop shadow=popink,left=18pt,right=18pt,top=13pt,bottom=13pt,before skip=6pt,after skip=20pt,
  before upper={\poppill{Sommaire}{poppurple}{white}\par\vspace{9pt}\popsemi\fontsize{10.6}{14.5}\selectfont\color{popink}}}

% Signature de fin de cours — poussée tout en bas de la dernière page
\newcommand{\findecours}{%
  \par\vfill\nopagebreak
  \begin{center}\popsquiggle{poporange}\end{center}\vspace{4pt}
  \signatureonbuch
}
\AtBeginDocument{\def\llbracket{\mathopen{[\mkern-3.5mu[}}\def\rrbracket{\mathclose{]\mkern-3.5mu]}}} % intervalles d'entiers (glyphe ⟦ absent de la police)
% Glyphes absents de Fira Math : redéfinis à partir de glyphes disponibles
\AtBeginDocument{%
  \def\setminus{\mathbin{\backslash}}\let\smallsetminus\setminus
  \def\vdots{\mathord{\vbox{\baselineskip3.2pt\lineskiplimit0pt\kern4pt\hbox{.}\hbox{.}\hbox{.}}}}%
  \def\ddots{\mathinner{\mkern1mu\raise6.5pt\hbox{.}\mkern2mu\raise3.6pt\hbox{.}\mkern2mu\raise0.7pt\hbox{.}\mkern1mu}}%
  \def\triangle{\mathord{\scalebox{0.9}{$\symup{\Delta}$}}}%
  \def\nabla{\mathord{\raisebox{\depth}{\scalebox{1}[-1]{$\symup{\Delta}$}}}}%
  \def\checkmark{\popcheck{popdark}}%
  \def\star{\mathbin{\ast}}\def\bigstar{\mathord{\ast}}%
  \def\diamond{\mathbin{\scalebox{0.6}{$\Diamond$}}}%
  \def\bigcup{\mathop{\mathchoice{\raisebox{-0.3em}{\scalebox{1.7}{$\cup$}}}{\scalebox{1.25}{$\cup$}}{\cup}{\cup}}\displaylimits}%
  \def\bigcap{\mathop{\mathchoice{\raisebox{-0.3em}{\scalebox{1.7}{$\cap$}}}{\scalebox{1.25}{$\cap$}}{\cap}{\cap}}\displaylimits}%
  \def\lesseqqgtr{\mathrel{\lessgtr}}\def\bigoplus{\mathop{\oplus}\displaylimits}%
}
\providecommand{\ssi}{si et seulement si} % abréviation fréquente
\AtBeginDocument{\providecommand{\wideparen}{\overparen}} % arc de cercle

%% FIGFIX-BEGIN v4 — correctifs globaux des figures (docs/onbuch-generation/tools/figfix)
\usepackage{adjustbox}\usepackage{etoolbox}
% toute figure TikZ/pgfplots plus large que la ligne est réduite à la largeur disponible
\BeforeBeginEnvironment{tikzpicture}{\begin{adjustbox}{max width=\linewidth}}
\AfterEndEnvironment{tikzpicture}{\end{adjustbox}}
% légendes pgfplots lisibles par-dessus les courbes
\pgfplotsset{every axis/.append style={legend style={fill=white,fill opacity=0.92,text opacity=1,draw=black!15,inner sep=3pt}}}
% étiquettes d'arêtes (syntaxe edge["..."]) avec fond blanc
\tikzset{every edge quotes/.append style={fill=white,inner sep=1.2pt}}
%% FIGFIX-END
\begin{document}

\pagegarde{Lutte contre l'effet de serre}{SVT}{1ère TI}{Module 3 : L'Éducation à l'Environnement et au Développement Durable}{NCT12}{3 h}{Cette leçon explique le mécanisme naturel de l'effet de serre, montre comment les activités humaines l'intensifient, analyse les impacts du réchauffement climatique et propose des actions concrètes pour réduire les émissions de gaz à effet de serre.
\vspace{4pt}
\begin{multicols}{2}\small
\begin{itemize}
\item À la fin de cette leçon, je sais expliquer le mécanisme de l'effet de serre naturel à partir d'un schéma.
\item Je sais identifier les principaux gaz à effet de serre (CO₂, CH₄, N₂O, vapeur d'eau) et leur rôle.
\item Je sais relier une activité humaine (transport, agriculture, déforestation, industrie) à l'émission de gaz à effet de serre.
\item Je sais décrire les conséquences majeures du réchauffement climatique (montée des eaux, événements extrêmes, impacts sur l'agriculture et la santé).
\item Je sais proposer des gestes quotidiens et des politiques publiques pour réduire les émissions de GES.
\item Je sais calculer une empreinte carbone simplifiée pour une activité donnée.
\end{itemize}\end{multicols}}

\begin{sommaire}
\begin{multicols}{2}
\begin{enumerate}
\item 1. L'effet de serre naturel : principe et gaz impliqués
\item 2. Origines anthropiques du renforcement de l'effet de serre
\item 3. Conséquences du réchauffement climatique
\item 4. Moyens d'action : atténuation et adaptation
\item 5. Synthèse, modélisation simple et engagement citoyen
\item Activité d'intégration
\item Méthodes et exercices résolus
\item Exercices
\item Corrigés des exercices
\item Fiche bilan
\end{enumerate}
\end{multicols}
\end{sommaire}

\begin{objectifs}
\begin{itemize}
\item À la fin de cette leçon, je sais expliquer le mécanisme de l'effet de serre naturel à partir d'un schéma.
\item Je sais identifier les principaux gaz à effet de serre (CO₂, CH₄, N₂O, vapeur d'eau) et leur rôle.
\item Je sais relier une activité humaine (transport, agriculture, déforestation, industrie) à l'émission de gaz à effet de serre.
\item Je sais décrire les conséquences majeures du réchauffement climatique (montée des eaux, événements extrêmes, impacts sur l'agriculture et la santé).
\item Je sais proposer des gestes quotidiens et des politiques publiques pour réduire les émissions de GES.
\item Je sais calculer une empreinte carbone simplifiée pour une activité donnée.
\item Je sais argumenter, à l'aide de données chiffrées, l'urgence d'agir pour limiter le réchauffement à 1,5 °C.
\end{itemize}
\end{objectifs}

\begin{prerequis}
\begin{itemize}
\item Notion de rayonnement électromagnétique (visible, infrarouge) vue en physique en classe de Seconde.
\item Connaissance de base du cycle du carbone (photosynthèse, respiration) du programme de SVT de Seconde.
\item Notions de pourcentage et de calcul de moyennes (mathématiques, Première).
\item Repérage géographique du Cameroun et de ses principales zones climatiques (géographie, Première).
\item Sensibilisation aux enjeux du développement durable (éducation civique, Première).
\end{itemize}
\end{prerequis}

\newpage

\coursec{L'effet de serre naturel : principe et gaz impliqués}

\textbf{Situation d'accroche.} Il est 14 h dans la cour du lycée. Sous l'auvent, les élèves s'essuient le front : « Il fait de plus en plus chaud, même en saison sèche on n'avait pas autant chaud avant ! » s'exclame Nadège. Un autre élève ajoute : « Chez ma grand-mère au village, la saison des pluies arrivait en mars ; maintenant, elle arrive parfois en avril, et les semis de maïs sont retardés. » Le professeur de SVTEEHB, qui passait par là, s'arrête et leur pose une question : « Pourquoi la Terre se réchauffe-t-elle, et que pouvons-nous faire, ici à Yaoundé, pour limiter ce phénomène ? » Cette question lance notre enquête scientifique sur l'\cle{effet de serre} et ses conséquences.

\textbf{Problématique.} Comment l'atmosphère terrestre maintient-elle naturellement une température compatible avec la vie ? Quels gaz en sont responsables, et comment mesure-t-on leur présence ?

\courssub{Une observation de la vie quotidienne : la voiture garée au soleil}

Avant de parler de la planète entière, partons d'une expérience que tout le monde a vécue. Une voiture garée en plein soleil, vitres fermées, devient rapidement une étuve : la température intérieure peut dépasser \SI{50}{\celsius} alors qu'il fait \SI{32}{\celsius} dehors. Pourquoi ?

\begin{itemize}
\item Le rayonnement solaire (lumière visible) traverse facilement les vitres et chauffe les sièges et le tableau de bord ;
\item ces surfaces chauffées réémettent de la chaleur sous forme de rayonnement \cle{infrarouge} (invisible) ;
\item or le verre laisse passer la lumière visible mais \textbf{bloque en grande partie l'infrarouge} : la chaleur reste piégée à l'intérieur.
\end{itemize}

C'est exactement le principe d'une serre de jardinier — d'où le nom du phénomène. Sur Terre, ce ne sont pas des vitres qui jouent ce rôle, mais certains \textbf{gaz de l'atmosphère}. Transposons :

\begin{center}
\begin{tabularx}{\linewidth}{|l|X|X|}
\hline
\rowcolor{popblueL} \textbf{Élément} & \textbf{Voiture au soleil} & \textbf{Planète Terre} \\
\hline
Source d'énergie & Le Soleil & Le Soleil \\
\hline
Surface chauffée & Sièges, tableau de bord & Sol et océans \\
\hline
« Vitre » piégeant la chaleur & Les glaces de la voiture & Les gaz à effet de serre (GES) \\
\hline
Chaleur piégée & Infrarouge bloqué par le verre & Infrarouge absorbé par les GES \\
\hline
\end{tabularx}
\end{center}

\courssub{Définition et rôle vital de l'effet de serre naturel}

\begin{definition}[Effet de serre naturel]
L'\cle{effet de serre naturel} est le phénomène par lequel certains gaz de l'atmosphère, appelés \cle{gaz à effet de serre} (GES), absorbent puis réémettent le rayonnement infrarouge émis par la surface terrestre, renvoyant une partie de cette chaleur vers le sol. Ce processus maintient la température moyenne de la Terre à environ \SI{15}{\celsius}, rendant la planète habitable.
\end{definition}

\begin{aretenir}[Un phénomène vital, pas un problème en soi]
Sans effet de serre naturel, la température moyenne de la Terre serait d'environ \SI{-18}{\celsius} : les océans seraient gelés et la vie telle que nous la connaissons serait impossible. L'effet de serre est donc un \textbf{phénomène naturel et indispensable}. Le problème actuel n'est pas son existence, mais son \textbf{renforcement} par les activités humaines (étudié dans la section suivante).
\end{aretenir}

\begin{savaistu}[La Lune, témoin d'une planète sans atmosphère]
La Lune, située à la même distance du Soleil que la Terre mais dépourvue d'atmosphère, subit des écarts de température extrêmes : environ \SI{+120}{\celsius} le jour et \SI{-170}{\celsius} la nuit. Sans « couverture » de gaz à effet de serre, aucune régulation thermique n'est possible. Notre atmosphère agit comme une couverture qui lisse les températures entre le jour et la nuit, entre l'équateur et les pôles.
\end{savaistu}

\courssub{Le flux radiatif : le voyage de l'énergie solaire}

Pour expliquer rigoureusement le mécanisme, suivons le parcours de l'énergie depuis le Soleil jusqu'à l'espace. C'est le \cle{bilan radiatif} de la Terre.

\begin{enumerate}
\item \textbf{Rayonnement solaire incident.} Le Soleil émet un rayonnement riche en lumière visible et en proche infrarouge. Environ \SI{30}{\percent} de cette énergie est réfléchie vers l'espace par les nuages, les glaces et les surfaces claires (c'est l'\cle{albédo}) ; le reste traverse l'atmosphère.
\item \textbf{Absorption par la surface.} Le sol et les océans absorbent le rayonnement solaire et s'échauffent.
\item \textbf{Réémission infrarouge.} Tout corps chauffé émet un rayonnement dont la longueur d'onde dépend de sa température. La Terre, à environ \SI{15}{\celsius}, réémet l'énergie reçue sous forme d'\textbf{infrarouge thermique} (invisible à l'œil).
\item \textbf{Absorption et réémission par les GES.} Les molécules de GES (CO$_2$, CH$_4$, H$_2$O, N$_2$O...) absorbent une grande partie de cet infrarouge, puis le réémettent \textbf{dans toutes les directions}, y compris vers le sol. La surface reçoit donc de l'énergie à la fois du Soleil et de l'atmosphère : elle est plus chaude qu'elle ne le serait sans GES.
\item \textbf{Équilibre.} En régime naturel stable, l'énergie reçue du Soleil égale l'énergie renvoyée vers l'espace : la température moyenne reste constante.
\end{enumerate}

\begin{popfigure}
\begin{center}
\begin{tikzpicture}[scale=1.0, every node/.style={font=\small}]
% Soleil
\shade[inner color=popgold, outer color=poporange] (-4.6,4.6) circle (0.75);
\node[popdark] at (-4.6,5.7) {\textbf{Soleil}};
% Sol
\fill[popgreenL] (-6,-2.2) rectangle (6,-1.4);
\draw[popgreen!60!black] (-6,-1.4) -- (6,-1.4);
\node[popdark] at (4.6,-1.85) {Surface terrestre};
% Atmosphère (bande)
\fill[popblue!18] (-6,1.6) rectangle (6,2.6);
\draw[popblue, dashed] (-6,1.6) -- (6,1.6);
\draw[popblue, dashed] (-6,2.6) -- (6,2.6);
\node[popblue] at (4.9,2.1) {\textbf{Atmosphère (GES)}};
% Molécules GES
\foreach \x/\y in {-3.5/2.1,-1.2/2.35,1.0/1.9,3.0/2.3,-2.2/1.85,2.0/2.45}{
  \fill[popblue] (\x,\y) circle (0.09);}
\node[popblue, font=\footnotesize] at (-4.4,2.1) {CO$_2$, CH$_4$, H$_2$O, N$_2$O};
% Rayonnement solaire incident
\draw[-{Stealth[length=3mm]}, line width=1.4pt, poporange] (-3.9,3.9) -- (-2.2,1.7);
\draw[-{Stealth[length=3mm]}, line width=1.4pt, poporange] (-2.2,1.55) -- (-0.6,-1.35);
\node[poporange, align=center, font=\footnotesize] at (-3.4,0.3) {Rayonnement solaire\\(visible) : traverse\\l'atmosphère};
% Part réfléchie
\draw[-{Stealth[length=3mm]}, line width=1.2pt, popmuted] (-5.2,3.6) -- (-4.4,2.75);
\draw[-{Stealth[length=3mm]}, line width=1.2pt, popmuted] (-4.4,2.75) -- (-5.6,4.3);
\node[popmuted, font=\footnotesize, align=center] at (-5.5,3.35) {Part\\réfléchie};
% Infrarouge montant
\draw[-{Stealth[length=3mm]}, line width=1.4pt, poppink] (0.9,-1.35) -- (0.9,1.55);
\node[poppink, font=\footnotesize, align=center] at (0.9,0.15) {Infrarouge émis\\par la surface};
% Absorption + réémission vers le sol
\draw[-{Stealth[length=3mm]}, line width=1.4pt, poppink] (1.0,1.9) -- (1.0,-1.3);
\node[poppink, font=\footnotesize, align=center] at (2.6,0.35) {Chaleur réémise\\vers la surface\\par les GES};
% Infrarouge qui s'échappe
\draw[-{Stealth[length=3mm]}, line width=1.2pt, poppink] (3.6,-1.35) -- (3.6,4.4);
\node[poppink, font=\footnotesize, align=right] at (5.6,3.6) {Part qui\\s'échappe\\vers l'espace};
\end{tikzpicture}
\end{center}
\legende{Schéma du mécanisme de l'effet de serre naturel. Le rayonnement solaire (flèches orange) traverse l'atmosphère et chauffe la surface ; celle-ci réémet un rayonnement infrarouge (flèches roses) dont une partie est absorbée par les gaz à effet de serre (points bleus) puis réémise dans toutes les directions, notamment vers le sol, ce qui maintient la température moyenne autour de \SI{15}{\celsius}.}
\end{popfigure}

\begin{methode}[Lire un diagramme de bilan radiatif]
Face à un schéma de l'effet de serre au Probatoire, procède toujours dans cet ordre :
\begin{enumerate}
\item \textbf{Repère les flux entrants} : flèches venant du Soleil (rayonnement visible, courtes longueurs d'onde). Distingue la part réfléchie (albédo) de la part absorbée par la surface.
\item \textbf{Repère les flux sortants} : flèches partant de la surface vers le haut (infrarouge, grandes longueurs d'onde).
\item \textbf{Localise les GES} : cherche où l'infrarouge est intercepté (flèche qui « rebondit » ou se divise dans la couche atmosphérique).
\item \textbf{Conclus} : la part d'infrarouge renvoyée vers la surface par les GES est la manifestation de l'effet de serre ; si cette part augmente, l'effet de serre se renforce.
\end{enumerate}
\end{methode}

\courssub{Les principaux gaz à effet de serre}

Tous les gaz de l'atmosphère ne sont pas des GES. L'azote (N$_2$, 78 \%) et l'oxygène (O$_2$, 21 \%), qui forment l'essentiel de l'air, sont \textbf{transparents à l'infrarouge} : ce ne sont pas des gaz à effet de serre. Les GES sont des gaz traces, présents en très faible quantité, mais à l'effet considérable.

\begin{center}
\begin{tabularx}{\linewidth}{|l|X|l|X|}
\hline
\rowcolor{popblueL} \textbf{Gaz} & \textbf{Sources principales} & \textbf{PRG (100 ans)} & \textbf{Concentration typique} \\
\hline
Vapeur d'eau H$_2$O & Évaporation naturelle & non chiffré (rétroaction) & variable (jusqu'à \SI{4}{\percent}) \\
\hline
Dioxyde de carbone CO$_2$ & Combustion (charbon, pétrole, bois), déforestation, cimenteries & 1 (référence) & $\approx$ 420 ppm \\
\hline
Méthane CH$_4$ & Riziculture, élevage (ruminants), décharges, fuites de gaz & $\approx$ 28 & $\approx$ 1\,900 ppb \\
\hline
Protoxyde d'azote N$_2$O & Engrais azotés, industrie chimique & $\approx$ 265 & $\approx$ 335 ppb \\
\hline
\end{tabularx}
\end{center}

\begin{propriete}[Pouvoir de réchauffement global (PRG)]
Le \cle{pouvoir de réchauffement global} (PRG, en anglais GWP) est un indicateur qui compare l'effet de réchauffement d'une masse donnée d'un gaz à celui de la même masse de CO$_2$ sur une période de 100 ans. Par convention, le PRG du CO$_2$ vaut 1. Ainsi, 1 kg de méthane réchauffe environ 28 fois plus que 1 kg de CO$_2$ sur un siècle, et 1 kg de N$_2$O environ 265 fois plus.
\end{propriete}

\begin{attention}[GES ne veut pas dire polluant de l'air]
Ne confonds pas \textbf{gaz à effet de serre} et \textbf{polluants atmosphériques}. Les oxydes d'azote (NO$_x$), le dioxyde de soufre (SO$_2$) ou les particules fines dégradent la qualité de l'air et provoquent des maladies respiratoires, mais ne sont pas les principaux responsables du réchauffement. Inversement, le CO$_2$ n'est pas toxique aux concentrations actuelles : c'est un GES, pas un « poison » de l'air. Les deux notions répondent à des problèmes différents, même si certaines sources (comme les pots d'échappement) émettent les deux.
\end{attention}

\begin{savaistu}[La vapeur d'eau : le GES le plus abondant]
La vapeur d'eau est de loin le gaz à effet de serre le plus abondant de l'atmosphère. Pourtant, on ne la cite pas comme cause du réchauffement actuel : sa concentration dépend directement de la température (l'air chaud peut contenir plus de vapeur d'eau). Elle agit donc comme une \textbf{rétroaction positive} : quand le CO$_2$ réchauffe l'air, celui-ci retient plus de vapeur d'eau, ce qui amplifie encore le réchauffement. C'est un amplificateur, non un déclencheur.
\end{savaistu}

\courssub{Mesurer les GES : les unités ppm et ppb}

Les GES étant présents en traces, on exprime leur concentration avec des unités adaptées :

\begin{definition}[ppm et ppb]
La \cle{ppm} (partie par million) indique le nombre de molécules d'un gaz pour un million de molécules d'air : 420 ppm de CO$_2$ signifie 420 molécules de CO$_2$ pour 1\,000\,000 de molécules d'air, soit \SI{0,042}{\percent}. La \cle{ppb} (partie par milliard) est mille fois plus petite : 1\,900 ppb de CH$_4$ = 1,9 ppm.
\end{definition}

\begin{exemplebox}[L'évolution du CO$_2$ depuis l'ère industrielle]
Les carottes de glace prélevées en Antarctique permettent de mesurer le CO$_2$ piégé dans les bulles d'air anciennes. Résultat : la concentration était stable autour de \textbf{280 ppm} pendant des millénaires avant 1850 ; elle atteint aujourd'hui environ \textbf{420 ppm}, soit une hausse de \SI{50}{\percent} en moins de deux siècles. Cette augmentation rapide et sans précédent coïncide avec l'industrialisation.
\end{exemplebox}

\begin{popfigure}
\begin{center}
\begin{tikzpicture}
\begin{axis}[
  width=0.9\linewidth, height=5.6cm,
  xlabel={Année}, ylabel={CO$_2$ (ppm)},
  xmin=1850, xmax=2025, ymin=270, ymax=430,
  xtick={1850,1900,1950,2000,2025},
  ytick={280,320,360,400},
  grid=both, grid style={popline!50},
  axis line style={popdark},
  label style={popdark},
  legend style={font=\footnotesize, at={(0.03,0.97)}, anchor=north west}
]
\addplot[popcurve, domain=1850:2025, samples=120] {280 + 140*((x-1850)/175)^2.6};
\addlegendentry{Concentration en CO$_2$ (allure)}
\addplot[dashed, popmuted, domain=1850:2025] {280};
\addlegendentry{Niveau préindustriel (280 ppm)}
\end{axis}
\end{tikzpicture}
\end{center}
\legende{Allure de l'évolution de la concentration atmosphérique en CO$_2$ depuis 1850 : longue stabilité autour de 280 ppm, puis croissance accélérée jusqu'à environ 420 ppm aujourd'hui. Cette courbe illustre le renforcement récent de l'effet de serre.}
\end{popfigure}

\courssub{Effet de serre naturel et effet de serre renforcé : l'équilibre en question}

Résumons la distinction essentielle, cœur de cette leçon :

\begin{itemize}
\item \textbf{Effet de serre naturel} : avec des concentrations de GES stables (CO$_2$ à 280 ppm), l'énergie entrante égale l'énergie sortante. La température moyenne reste stable autour de \SI{15}{\celsius}. C'est un \textbf{équilibre} qui a permis le développement de la vie et des civilisations humaines.
\item \textbf{Effet de serre renforcé} : quand les activités humaines ajoutent des GES (CO$_2$ à 420 ppm, méthane en forte hausse), l'atmosphère absorbe et réémet davantage d'infrarouge vers la surface. L'énergie entrante devient \textbf{supérieure} à l'énergie sortante : la planète accumule de la chaleur, c'est le \textbf{déséquilibre} responsable du réchauffement climatique observé.
\end{itemize}

\begin{aretenir}[L'essentiel de la section]
\begin{itemize}
\item L'effet de serre est un phénomène \textbf{naturel et vital} : sans lui, la Terre serait à \SI{-18}{\celsius} en moyenne au lieu de \SI{15}{\celsius}.
\item Les GES (CO$_2$, CH$_4$, N$_2$O, vapeur d'eau) absorbent l'infrarouge émis par la surface et en réémettent une partie vers le sol.
\item Le \textbf{PRG} compare le pouvoir de réchauffement d'un gaz à celui du CO$_2$ (CO$_2$ = 1 ; CH$_4$ $\approx$ 28 ; N$_2$O $\approx$ 265).
\item Les concentrations se mesurent en \textbf{ppm} (CO$_2$ : 420 ppm contre 280 ppm avant l'industrialisation) ou en \textbf{ppb} (CH$_4$).
\item Le problème actuel est le \textbf{renforcement} de l'effet de serre, qui rompt l'équilibre radiatif.
\end{itemize}
\end{aretenir}

\begin{exoresolu}[Conversion ppm/ppb et calcul de hausse]
\textbf{Énoncé.} (1) La concentration en méthane est de 1\,900 ppb. Exprime-la en ppm. (2) Calcule le pourcentage d'augmentation du CO$_2$ entre l'ère préindustrielle (280 ppm) et aujourd'hui (420 ppm). (3) Un élève affirme : « Le méthane est négligeable car il est 220 fois moins abondant que le CO$_2$. » Corrige cette affirmation à l'aide du PRG.
\tcblower
\textbf{Solution détaillée.}
\begin{enumerate}
\item 1 ppm = 1\,000 ppb, donc $1\,900\ \text{ppb} = \dfrac{1\,900}{1\,000} = 1{,}9\ \text{ppm}$.
\item Augmentation : $\dfrac{420 - 280}{280} \times 100 = \dfrac{140}{280} \times 100 = \SI{50}{\percent}$. La concentration en CO$_2$ a augmenté de \SI{50}{\percent}.
\item L'abondance seule ne suffit pas à juger l'effet : il faut tenir compte du PRG. Avec un PRG de 28, chaque molécule de CH$_4$ réchauffe environ 28 fois plus qu'une molécule de CO$_2$ sur 100 ans. L'effet relatif du méthane est donc de l'ordre de $\dfrac{28}{220} \approx \SI{13}{\percent}$ de celui du CO$_2$ : loin d'être négligeable, le méthane est le deuxième contributeur au renforcement de l'effet de serre.
\end{enumerate}
\end{exoresolu}

\begin{exercice}[À toi de jouer]
(1) Explique, en quatre phrases maximum et à l'aide d'un schéma simple, pourquoi la température moyenne de la Terre est de \SI{15}{\celsius} et non de \SI{-18}{\celsius}. (2) Le protoxyde d'azote a un PRG de 265 : que signifie ce chiffre ? (3) Pourquoi dit-on que la vapeur d'eau est une rétroaction et non une cause initiale du réchauffement ?
\end{exercice}

\begin{corrige}[Exercice « À toi de jouer »]
(1) Le Soleil chauffe la surface terrestre, qui réémet de l'infrarouge. Les gaz à effet de serre absorbent une partie de cet infrarouge et le réémettent dans toutes les directions, dont vers le sol. Ce renvoi de chaleur élève la température moyenne de \SI{-18}{\celsius} (valeur sans atmosphère) à environ \SI{15}{\celsius}. Le schéma doit montrer : flèche solaire descendante, flèche infrarouge ascendante, flèche de réémission des GES vers la surface.
(2) Un PRG de 265 signifie qu'une masse de N$_2$O réchauffe l'atmosphère 265 fois plus que la même masse de CO$_2$ sur une période de 100 ans.
(3) La concentration en vapeur d'eau est contrôlée par la température : elle ne peut augmenter durablement que si l'air se réchauffe d'abord (à cause des autres GES). Elle amplifie donc un réchauffement déclenché par ailleurs, sans pouvoir l'initier.
\end{corrige}

\coursec{Origines anthropiques du renforcement de l'effet de serre}

Dans la section précédente, nous avons vu que l'\cle{effet de serre} est un phénomène \emph{naturel} et indispensable : sans lui, la température moyenne de la Terre serait d'environ $-18\ ^\circ\text{C}$ au lieu de $+15\ ^\circ\text{C}$. Les \cle{gaz à effet de serre} (GES) — vapeur d'eau, dioxyde de carbone \ce{CO2}, méthane \ce{CH4}, protoxyde d'azote \ce{N2O} — absorbent une partie du rayonnement infrarouge émis par le sol et le réémettent dans toutes les directions, dont vers la surface. Le problème actuel n'est donc pas l'effet de serre lui-même, mais son \textbf{renforcement} : depuis la révolution industrielle, les activités humaines injectent dans l'atmosphère des quantités croissantes de GES, bien au-delà de ce que les puits naturels peuvent absorber. Cette section répond à une question simple : \emph{d'où viennent ces émissions supplémentaires ?}

\courssub{Émissions anthropiques : définition et situation concrète}

\textbf{Observation de départ.} À Douala, un matin de semaine : des milliers de motos-taxis et de voitures embouteillées sur le pont de la Wouri, les cheminées des zones industrielles de Bonabéri, les groupes électrogènes qui ronronnent pendant les délestages, la fumée des feux de brousse et des charbonniers en périphérie. Toutes ces activités ont un point commun invisible : elles rejettent des gaz, en particulier du \ce{CO2}, qui s'accumulent dans l'atmosphère.

\begin{definition}[Émissions anthropiques]
On appelle \cle{émissions anthropiques} de gaz à effet de serre l'ensemble des rejets de GES dans l'atmosphère qui résultent des \textbf{activités humaines} (production et consommation d'énergie, agriculture, élevage, déforestation, industrie, gestion des déchets). Elles s'ajoutent aux \textbf{flux naturels} (respiration, décomposition, volcans, feux spontanés) et se distinguent d'eux par leur origine et par leur croissance rapide depuis 1750.
\end{definition}

\begin{attention}[Ne pas confondre flux naturels et émissions anthropiques]
Les flux naturels (respiration des êtres vivants, décomposition de la matière organique) étaient, avant l'ère industrielle, \emph{équilibrés} par les absorptions naturelles (photosynthèse, dissolution dans les océans) : la concentration atmosphérique de \ce{CO2} était stable autour de 280 ppm. Les émissions anthropiques, elles, viennent \textbf{s'ajouter} à ce cycle équilibré. C'est ce supplément — et non la respiration des humains ou des animaux — qui fait augmenter la concentration de GES. Dire « le \ce{CO2} est naturel donc il n'y a pas de problème » est un raisonnement faux : c'est le \emph{déséquilibre} du bilan qui compte.
\end{attention}

\begin{propriete}[Bilan carbone]
Le \cle{bilan carbone} de la planète (ou d'un territoire) est la différence entre les \textbf{émissions} de GES et les \textbf{absorptions} par les puits (forêts, sols, océans) :
\[ \text{Bilan} = \text{Émissions} - \text{Absorptions par les puits} \]
\begin{itemize}
\item Si le bilan est \textbf{positif} (émissions $>$ absorptions), la concentration atmosphérique de GES \textbf{augmente} : c'est la situation actuelle.
\item Si le bilan est \textbf{négatif}, la concentration diminue.
\item La \textbf{neutralité carbone} correspond à un bilan nul.
\end{itemize}
Actuellement, les puits naturels n'absorbent qu'environ la moitié des émissions anthropiques de \ce{CO2} ; le reste s'accumule dans l'atmosphère (fraction aéroportée $\approx$ 45--50 \%). 
\end{propriete}

\courssub{Les mécanismes d'émission : combustion et fermentation anaérobie}

\textbf{Deux grandes familles de mécanismes} expliquent l'essentiel des émissions anthropiques : la \emph{combustion} (qui produit du \ce{CO2}) et la \emph{fermentation anaérobie} (qui produit du \ce{CH4}).

\textbf{a) La combustion : source majeure de \ce{CO2}.} Brûler un combustible carboné (bois, charbon, pétrole, gaz naturel) en présence de dioxygène libère de l'énergie, du dioxyde de carbone et de l'eau. Pour un combustible purement carboné (charbon), l'équation bilan est :
\[ \ce{C + O2 -> CO2} \]
Pour les hydrocarbures, le principe est identique. Exemple de l'octane, constituant majeur de l'essence :
\[ \ce{2C8H18 + 25O2 -> 16CO2 + 18H2O} \]
Chaque litre d'essence brûlé dans un moteur rejette environ \num{2,31} kg de \ce{CO2}. C'est ce mécanisme qui fait du \textbf{transport}, de la \textbf{production d'électricité} (centrales thermiques au fioul ou au gaz, groupes électrogènes) et de l'\textbf{industrie} (cimenteries, fours) les premières sources mondiales de \ce{CO2}.

\begin{methode}[Calculer l'émission de \ce{CO2} d'un trajet en voiture]
Pour estimer la masse de \ce{CO2} émise par un trajet en véhicule à essence :
\begin{enumerate}
\item Relever la \textbf{distance} parcourue $d$ (en km).
\item Connaître la \textbf{consommation} du véhicule $c$ (en L/100 km).
\item Calculer le volume de carburant : $V = d \times \dfrac{c}{100}$.
\item Multiplier par le \textbf{facteur d'émission} : $m_{\ce{CO2}} = V \times \num{2,31}$ (en kg).
\end{enumerate}
Formule synthétique : $m_{\ce{CO2}} = d \times \dfrac{c}{100} \times \num{2,31}$.
\end{methode}

\begin{exoresolu}[Émissions d'un trajet Douala--Yaoundé]
\textbf{Énoncé.} Une voiture à essence consomme en moyenne 8 L/100 km. Elle effectue un aller-retour Douala--Yaoundé, soit $2 \times 230 = 460$ km. 1) Calculer la masse de \ce{CO2} émise. 2) Comparer avec l'absorption annuelle d'un arbre adulte, estimée à 25 kg de \ce{CO2}/an.
\tcblower
\textbf{Solution.} 1) Volume d'essence consommé :
\[ V = 460 \times \frac{8}{100} = \num{36,8}\ \text{L} \]
Masse de \ce{CO2} émise :
\[ m = \num{36,8} \times \num{2,31} \approx 85\ \text{kg de } \ce{CO2} \]
2) Il faudrait $\dfrac{85}{25} \approx 3,4$ arbres adultes pendant une année entière pour absorber le \ce{CO2} d'un seul aller-retour Douala--Yaoundé en voiture. Ce calcul simple montre pourquoi la multiplication des trajets motorisés alourdit rapidement le bilan carbone.
\end{exoresolu}

\textbf{b) La fermentation anaérobie : source majeure de \ce{CH4}.} En l'absence de dioxygène (milieu \emph{anaérobie}), des micro-organismes (archées méthanogènes) décomposent la matière organique en produisant du méthane. On parle de \cle{méthanisation} ou fermentation entérique. Trois sources anthropiques dominent :
\begin{itemize}
\item \textbf{Les ruminants} (bovins, ovins, caprins) : dans leur panse, les micro-organismes dégradent la cellulose des fourrages et libèrent du \ce{CH4}, évacué principalement par éructation. Un bovin émet de 50 à 100 kg de \ce{CH4} par an.
\item \textbf{La riziculture inondée} : l'eau qui recouvre les rizières bloque l'oxygène du sol ; la matière organique du sol est alors décomposée en anaérobiose, avec production de \ce{CH4} qui diffuse à travers les plants de riz vers l'atmosphère.
\item \textbf{Les décharges de déchets} : les ordures enfouies fermentent sans oxygène et dégagent un « biogaz » riche en méthane.
\end{itemize}

\begin{attention}[Le méthane : puissant mais de courte durée de vie]
Le \ce{CH4} a une durée de vie atmosphérique d'environ \textbf{12 ans} seulement (contre des siècles pour le \ce{CO2}), mais son \cle{pouvoir de réchauffement global} (PRG) est environ \textbf{28 à 30 fois} supérieur à celui du \ce{CO2} sur 100 ans (et plus de 80 fois sur 20 ans). Conséquence stratégique : réduire les émissions de méthane (meilleure gestion des décharges, alimentation des ruminants, riziculture améliorée) produit un effet \textbf{rapide} sur le rythme du réchauffement, alors que les réductions de \ce{CO2} agissent sur le long terme. Les deux sont nécessaires.
\end{attention}

\begin{savaistu}[Riz et méthane : la technique SRI]
La culture du riz en zone inondée permanente est une source importante de \ce{CH4}. La \textbf{SRI} (Système de Riziculture Intensif), qui alterne des phases d'inondation et de séchage des parcelles au lieu d'une submersion continue, réoxygène le sol et peut réduire les émissions de méthane de \textbf{30 à 50 \%}, tout en économisant l'eau et en augmentant souvent les rendements. Des projets pilotes de SRI sont expérimentés dans les plaines rizicoles du Cameroun (vallées du Logone, plaine des Mbo).
\end{savaistu}

\courssub{Les principales sources d'émissions par secteur}

\textbf{Au niveau mondial}, on classe les émissions anthropiques en grands secteurs (données GIEC / IEA 2022-2023, en \ce{CO2} équivalent) :
\begin{itemize}
\item \textbf{Énergie (électricité et chaleur)} : combustion de charbon, pétrole et gaz dans les centrales — premier secteur émetteur mondial.
\item \textbf{Industrie} : combustion pour les fours (cimenteries, sidérurgie) et réactions chimiques émettrices (la fabrication du ciment libère du \ce{CO2} par calcination du calcaire : $\ce{CaCO3 -> CaO + CO2}$).
\item \textbf{Transport} : carburants routiers, aériens, maritimes.
\item \textbf{Bâtiments} : chauffage, climatisation, cuisson (résidentiel et tertiaire).
\item \textbf{Agriculture et élevage} : méthane des ruminants et des rizières, \ce{N2O} des engrais azotés.
\item \textbf{Changement d'affectation des sols} : déforestation et dégradation des tourbières (émissions nettes).
\item \textbf{Déchets} : méthane des décharges et des eaux usées.
\end{itemize}

\begin{center}
\begin{popfigure}
\begin{tikzpicture}
\begin{axis}[
  ybar, bar width=0.55cm,
  width=\linewidth, height=7cm,
  ymin=0, ymax=18,
  ylabel={Émissions (Gt \ce{CO2} équivalent/an)},
  symbolic x coords={Énergie, Industrie, Transport, Bâtiments, Agriculture, Sols/déforestation, Déchets},
  xtick=data, x tick label style={rotate=25, anchor=east, font=\small},
  nodes near coords, nodes near coords style={font=\small},
  ymajorgrids, grid style={popline!40},
]
\addplot[fill=popblue!70, draw=popblue] coordinates {(Énergie,15.8) (Industrie,6.3) (Transport,8.4) (Bâtiments,3.3) (Agriculture,6.2) (Sols/déforestation,3.9) (Déchets,1.6)};
\end{axis}
\end{tikzpicture}
\legende{Ordres de grandeur des émissions mondiales de gaz à effet de serre par secteur (en gigatonnes de \ce{CO2} équivalent par an, données 2022-2023, arrondies). Le secteur de l'énergie (électricité et chaleur) domine, suivi du transport et de l'industrie ; l'agriculture et la déforestation contribuent aussi fortement, notamment en \ce{CH4} et \ce{N2O}.}
\end{popfigure}
\end{center}

\textbf{Au niveau du Cameroun}, le profil est différent de celui des pays industrialisés : les émissions totales restent modestes (de l'ordre de quelques dizaines de Mt de \ce{CO2} par an, soit moins de 0,1 \% des émissions mondiales), mais elles croissent rapidement avec l'urbanisation et la motorisation.

\begin{exemplebox}[Le secteur des transports au Cameroun]
Au Cameroun, le secteur des \textbf{transports} représente environ \textbf{30 \%} des émissions nationales de \ce{CO2}, soit de l'ordre de \textbf{12 Mt de \ce{CO2}/an}. Les causes : un parc automobile vieillissant (véhicules d'occasion importés, très consommateurs), l'essor des motos-taxis, l'engorgement de Douala et Yaoundé, et l'usage massif des groupes électrogènes lors des coupures d'électricité. Viennent ensuite la production d'énergie, l'industrie (cimenteries de Douala et Nomayos, brasseries) et le changement d'affectation des sols.
\end{exemplebox}

\courssub{La déforestation : une double peine pour le climat}

\textbf{Observation.} Dans la région du Sud-Est du Cameroun (départements de la Kadey, du Boumba-et-Ngoko), l'avancée des pistes forestières, l'exploitation minière artisanale de l'or et l'agriculture itinérante sur brûlis fragmentent la grande forêt du bassin du Congo — le deuxième massif forestier du monde après l'Amazonie.

\textbf{Hypothèse.} La déforestation aggrave l'effet de serre non pas d'une mais de \emph{deux} façons.

\textbf{Données et interprétation.}
\begin{enumerate}
\item \textbf{Émission directe} : abattre et brûler la forêt (ou laisser se décomposer la biomasse) relâche dans l'atmosphère le carbone stocké dans les troncs, les racines et le sol : $\ce{C + O2 -> CO2}$.
\item \textbf{Perte de séquestration} : une forêt vivante est un \cle{puits de carbone} : par photosynthèse ($\ce{6CO2 + 6H2O -> C6H12O6 + 6O2}$), elle prélève du \ce{CO2} atmosphérique et le fixe dans sa biomasse. Détruire la forêt, c'est supprimer ce « pompage » : la capacité d'absorption diminue alors même que les émissions augmentent.
\end{enumerate}

\textbf{Conclusion.} La déforestation agit comme une \textbf{double peine climatique} : elle ajoute des émissions \emph{et} retire des absorptions. Au Cameroun, le taux de déforestation est estimé à environ \num{0,5} à 1 \% par an dans les zones les plus touchées, soit des centaines de milliers d'hectares perdus par décennie. La forêt camerounaise (environ 20 millions d'hectares) stocke pourtant des milliards de tonnes de carbone : sa conservation est un enjeu climatique national et mondial.

\begin{center}
\begin{popfigure}
\begin{tikzpicture}[scale=2.5]
% Contour simplifié du Cameroun
\draw[thick, popdark, fill=popgreenL] (0,0) -- (0.4,1.2) -- (0.2,2.4) -- (0.8,3.4) -- (1.2,4.4) -- (2.0,5.2) -- (3.2,5.6) -- (4.4,5.2) -- (5.0,4.4) -- (4.8,3.2) -- (5.2,2.2) -- (4.6,1.0) -- (3.4,0.2) -- (2.0,-0.3) -- (0.8,-0.2) -- cycle;
% Zone de déforestation Sud-Est
\fill[poporange!50, draw=poporange, thick] (3.0,0.4) -- (4.6,1.0) -- (5.2,2.2) -- (4.4,2.6) -- (3.2,2.0) -- cycle;
% Villes
\fill[popink] (0.55,0.6) circle (0.07) node[right, font=\small]{Douala};
\fill[popink] (1.7,1.5) circle (0.07) node[right, font=\small]{Yaoundé};
\fill[popink] (1.15,4.5) circle (0.06) node[left, font=\small]{Garoua};
\fill[popink] (0.35,1.6) circle (0.06) node[left, font=\small]{Bafoussam};
% Légende interne
\draw[poporange, thick, fill=poporange!50] (0.6,-0.9) rectangle (1.0,-0.6);
\node[right, font=\small] at (1.1,-0.75) {Zone de déforestation active (Sud-Est)};
\fill[popink] (0.75,-1.4) circle (0.06);
\node[right, font=\small] at (1.1,-1.4) {Ville émettrice majeure (transport, industrie)};
\end{tikzpicture}
\legende{Carte simplifiée du Cameroun. Les grandes agglomérations (Douala, Yaoundé) concentrent les émissions liées au transport et à l'industrie ; le Sud-Est forestier subit une déforestation active (orpaillage, agriculture sur brûlis, exploitation du bois) qui dégrade le puits de carbone du bassin du Congo.}
\end{popfigure}
\end{center}

\courssub{L'évolution historique des concentrations : la preuve par les mesures}

Comment savons-nous que les émissions anthropiques s'accumulent réellement ? Par la \textbf{mesure directe}. Depuis 1958, l'observatoire du Mauna Loa (Hawaï) mesure en continu la concentration de \ce{CO2} atmosphérique : c'est la célèbre \cle{courbe de Keeling}. Les carottes de glace de l'Antarctique permettent de remonter bien plus loin : les bulles d'air piégées dans la glace conservent la composition de l'atmosphère d'autrefois.

\begin{center}
\begin{popfigure}
\begin{tikzpicture}
\begin{axis}[
  width=\linewidth, height=7.5cm,
  xlabel={Année}, ylabel={Concentration de \ce{CO2} (ppm)},
  xmin=1750, xmax=2025, ymin=270, ymax=430,
  xtick={1750,1800,1850,1900,1950,2000,2025},
  ymajorgrids, xmajorgrids, grid style={popline!40},
  legend style={at={(0.03,0.97)}, anchor=north west, font=\small},
]
% Courbe lissée passant par les points clés (carottes de glace + Mauna Loa)
\addplot[popcurve, color=popblue, very thick, smooth, tension=0.6] coordinates {
  (1750,277) (1800,280) (1850,285) (1900,295) (1950,310) (1958,315) (1970,325) (1980,338) (1990,354) (2000,369) (2010,390) (2020,412) (2023,421)
};
\addlegendentry{Concentration de \ce{CO2} (reconstitution)}
\draw[dashed, poporange, thick] (axis cs:1958,270) -- (axis cs:1958,425);
\node[poporange, font=\small, anchor=west, rotate=90] at (axis cs:1963,300) {Début des mesures de Keeling (1958)};
\addplot[only marks, mark=*, mark size=1.5pt, color=popink] coordinates {(1750,277) (1850,285) (1958,315) (1990,354) (2010,390) (2023,421)};
\end{axis}
\end{tikzpicture}
\legende{Évolution de la concentration atmosphérique de \ce{CO2} depuis 1750 (reconstitution d'après carottes de glace et courbe de Keeling). Stable autour de 280 ppm pendant des millénaires, elle dépasse 420 ppm en 2023 — un niveau jamais atteint depuis au moins 800 000 ans. Le méthane est passé dans le même temps d'environ 700 à plus de 1900 ppb.}
\end{popfigure}
\end{center}

\textbf{Interprétation.} Trois faits établissent l'origine humaine de cette hausse :
\begin{itemize}
\item La \textbf{rupture de pente} coïncide avec la révolution industrielle et s'accélère après 1950, en parallèle de la consommation mondiale d'énergies fossiles.
\item L'\textbf{analyse isotopique} du carbone atmosphérique montre que le \ce{CO2} ajouté provient bien de la combustion de matière organique fossile (charbon, pétrole, gaz) : il est appauvri en carbone-14 (radioactif, disparu dans les fossiles) et en carbone-13 (signature de la photosynthèse ancienne), c'est l'effet Suess.
\item La hausse mesurée ($\approx$ 2,4 ppm/an actuellement) correspond à environ la \textbf{moitié} des émissions anthropiques annuelles : l'autre moitié est absorbée par les océans et la végétation, conformément au bilan carbone.
\end{itemize}

\begin{aretenir}[L'essentiel de la section]
\begin{itemize}
\item Les \textbf{émissions anthropiques} de GES proviennent de quatre grandes familles de sources : la \textbf{combustion d'énergies fossiles} (énergie, transport, industrie : \ce{CO2}), l'\textbf{agriculture et l'élevage} (fermentation anaérobie : \ce{CH4} ; engrais : \ce{N2O}), la \textbf{déforestation} (émission + perte de puits de carbone) et les \textbf{déchets} (\ce{CH4} des décharges).
\item Le \textbf{bilan carbone} (émissions $-$ absorptions) est positif : les concentrations de GES augmentent, comme le montre la courbe de Keeling (280 ppm en 1750, plus de 420 ppm en 2023).
\item Le \ce{CH4} réchauffe beaucoup plus que le \ce{CO2} à masse égale (PRG $\approx$ 28--30 sur 100 ans), mais disparaît en $\approx$ 12 ans : le réduire agit vite.
\item Au Cameroun, les transports ($\approx$ 30 \% des émissions nationales, $\approx$ 12 Mt \ce{CO2}/an) et la déforestation du Sud-Est sont les deux principaux leviers d'action.
\end{itemize}
\end{aretenir}

\begin{exercice}[À toi de jouer]
1) Un moto-taximan parcourt en moyenne 120 km par jour avec une moto consommant 3 L/100 km. Calcule la masse de \ce{CO2} émise en une année de 300 jours de travail. 2) Explique en quelques lignes pourquoi la protection de la forêt du Sud-Est camerounais contribue à la lutte contre le réchauffement climatique, en mobilisant les notions d'émission et de puits de carbone.
\end{exercice}

\begin{corrige}[Exercice]
1) Distance annuelle : $120 \times 300 = \num{36000}$ km. Volume d'essence : $V = \num{36000} \times \dfrac{3}{100} = \num{1080}$ L. Masse de \ce{CO2} : $m = \num{1080} \times \num{2,31} \approx \num{2495}$ kg, soit environ \textbf{2,5 tonnes de \ce{CO2} par an} pour un seul moto-taxi. 2) La forêt du Sud-Est stocke d'immenses quantités de carbone dans sa biomasse et ses sols. La détruire libère ce carbone sous forme de \ce{CO2} (émission) et supprime en même temps la photosynthèse qui prélevait du \ce{CO2} atmosphérique (perte de puits). La protéger maintient donc le bilan carbone moins positif : c'est un double bénéfice climatique.
\end{corrige}

Maintenant que les \emph{origines} du renforcement de l'effet de serre sont établies et quantifiées, la section suivante examinera les \textbf{conséquences concrètes du réchauffement climatique}, à l'échelle du globe comme à celle du Cameroun.

\coursec{Conséquences du réchauffement climatique}

Dans la section précédente, nous avons établi que les activités humaines (combustion des énergies fossiles, déforestation, élevage, agriculture, industrie) rejettent massivement des gaz à effet de serre, renforçant l'effet de serre naturel. Reste une question essentielle : quelles sont les conséquences concrètes de ce déséquilibre énergétique pour la planète, et plus particulièrement pour le Cameroun ? C'est l'objet de cette section.

\courssub{Augmentation de la température moyenne mondiale}

\textbf{Situation concrète.} À Maroua, dans l'Extrême-Nord du Cameroun, les anciens se souviennent de saisons fraîches plus longues ; aujourd'hui, les journées de chaleur accablante se multiplient et les températures nocturnes restent élevées. Cette perception correspond-elle à une réalité mesurée ?

\begin{definition}[Réchauffement climatique]
Le \cle{réchauffement climatique} est l'augmentation durable de la température moyenne de l'atmosphère et des océans à la surface du globe, due à l'accumulation des gaz à effet de serre (GES) émis par les activités humaines. Il s'exprime en degrés Celsius par rapport à une période de référence (généralement 1850--1900, début de l'ère industrielle).
\end{definition}

\textbf{Données de référence.} Selon le sixième rapport du GIEC (Groupe d'experts intergouvernemental sur l'évolution du climat, 2021--2023), la température moyenne mondiale à la surface du globe a augmenté d'environ \cle{\SI{1,1}{\celsius} depuis 1850--1900}. Chacune des quatre dernières décennies a été plus chaude que toutes celles qui l'ont précédée. Les dix années les plus chaudes jamais enregistrées sont toutes postérieures à 2010.

\begin{popfigure}
\begin{center}
\begin{tikzpicture}
\begin{axis}[
  width=\linewidth, height=6.5cm,
  xlabel={Année}, ylabel={Anomalie de température (\si{\celsius})},
  xmin=1850, xmax=2025, ymin=-0.6, ymax=1.6,
  xtick={1850,1900,1950,2000,2020},
  ytick={-0.5,0,0.5,1,1.5},
  grid=both, grid style={popline!40},
  legend style={at={(0.02,0.98)},anchor=north west,font=\small},
  axis line style={popdark}
]
\addplot[popblue!25, draw=none, forget plot] coordinates {
(1850,-0.25) (1880,-0.15) (1900,-0.1) (1920,-0.25) (1940,0.05) (1960,-0.05) (1980,0.1) (2000,0.5) (2010,0.75) (2020,1.15) (2023,1.35)
} -- (axis cs:2023,1.05) -- (axis cs:2020,0.85) -- (axis cs:2010,0.45) -- (axis cs:2000,0.2) -- (axis cs:1980,-0.2) -- (axis cs:1960,-0.35) -- (axis cs:1940,-0.25) -- (axis cs:1920,-0.55) -- (axis cs:1900,-0.4) -- (axis cs:1880,-0.45) -- (axis cs:1850,-0.55) -- cycle;
\addplot[popcurve, color=popink, very thick] coordinates {
(1850,-0.4) (1880,-0.3) (1900,-0.25) (1920,-0.4) (1940,-0.1) (1960,-0.2) (1980,-0.05) (2000,0.35) (2010,0.6) (2020,1.0) (2023,1.2)
};
\legend{Bande d'incertitude, Anomalie de température moyenne mondiale}
\end{axis}
\end{tikzpicture}
\end{center}
\legende{Évolution de l'anomalie de température moyenne mondiale de 1850 à 2023 (référence : moyenne 1850--1900). La bande grise représente l'incertitude des mesures, plus grande au XIX\textsuperscript{e} siècle (peu de stations météorologiques). La tendance est clairement ascendante, avec une accélération depuis les années 1980.}
\end{popfigure}

\textbf{Projections pour 2100.} Le GIEC a construit cinq scénarios d'émissions, appelés \cle{SSP} (\emph{Shared Socioeconomic Pathways}), qui décrivent des futurs possibles selon les politiques menées :

\begin{center}
\begin{tabularx}{\linewidth}{|l|X|c|}
\hline
\rowcolor{popblueL} \textbf{Scénario} & \textbf{Hypothèse sur les émissions} & \textbf{Réchauffement en 2100} \\
\hline
SSP1--2.6 & Réduction forte et rapide (neutralité carbone vers 2075) & environ \SI{1,8}{\celsius} \\
\hline
SSP2--4.5 & Émissions stabilisées puis en légère baisse & environ \SI{2,7}{\celsius} \\
\hline
SSP3--7.0 & Émissions en forte hausse (rivalités, fossiles) & environ \SI{3,6}{\celsius} \\
\hline
SSP5--8.5 & Croissance fondée sur les énergies fossiles & environ \SI{4,4}{\celsius} \\
\hline
\end{tabularx}
\end{center}

\begin{methode}[Analyser un scénario du GIEC]
Pour exploiter un tableau ou un graphique de projection climatique :
\begin{enumerate}
\item Identifier la \textbf{variable projetée} (anomalie de température $\Delta T$, élévation du niveau marin $\Delta h$) et l'\textbf{horizon} (souvent 2100).
\item Repérer le \textbf{scénario d'émissions} associé à chaque courbe ou ligne du tableau (SSP1--2.6 = optimiste, SSP5--8.5 = pessimiste).
\item Lire la \textbf{fourchette de valeurs} (meilleure estimation + intervalle de confiance) : une projection n'est jamais une valeur unique.
\item Comparer les scénarios pour dégager le \textbf{rôle des choix humains} : l'écart entre SSP1--2.6 et SSP5--8.5 (environ \SI{2,5}{\celsius}) mesure l'impact de nos politiques actuelles.
\item Conclure en reliant la projection aux conséquences locales (agriculture, santé, littoral).
\end{enumerate}
\end{methode}

\begin{propriete}[Rétroactions climatiques]
Une \cle{rétroaction climatique} est un processus par lequel une perturbation initiale du climat en provoque une seconde qui modifie à son tour la première. Les rétroactions \textbf{amplificatrices} (positives) aggravent le réchauffement. Exemple fondamental : la fonte de la glace diminue l'\cle{albédo} (pouvoir réfléchissant de la surface terrestre) ; la surface océanique sombre ainsi exposée absorbe davantage de rayonnement solaire, ce qui réchauffe encore la région et accélère la fonte. Autre exemple : le dégel du \cle{pergélisol} (sol gelé en permanence) libère du méthane \ce{CH4}, un GES puissant, qui renforce l'effet de serre.
\end{propriete}

\begin{attention}[Météo et climat : ne pas confondre]
On ne peut pas attribuer un événement météorologique isolé (une inondation précise, une saison sèche particulière) au seul changement climatique. La formulation rigoureuse est : le réchauffement \textbf{augmente la probabilité} ou \textbf{l'intensité} de certains événements extrêmes. De même, un hiver froid local ne « réfute » pas le réchauffement global : la météo décrit l'état instantané de l'atmosphère, le climat décrit des moyennes sur au moins 30 ans.
\end{attention}

\courssub{Fonte des glaces et élévation du niveau de la mer}

\textbf{Observation.} Les images satellites montrent que la \cle{banquise} arctique perd environ \SI{13}{\percent} de sa superficie estivale par décennie depuis 1979. Les \cle{calottes glaciaires} du Groenland et de l'Antarctique perdent chaque année des centaines de milliards de tonnes de glace.

\textbf{Mécanisme.} Deux phénomènes expliquent l'élévation du niveau des océans :
\begin{itemize}
\item la \cle{fonte des glaces continentales} (calottes du Groenland et de l'Antarctique, glaciers de montagne) qui ajoute de l'eau aux océans ;
\item la \cle{dilatation thermique} : l'eau de mer se réchauffe et augmente de volume.
\end{itemize}
La fonte de la banquise (glace flottante) ne fait pas monter le niveau de la mer, de même que la fonte d'un glaçon dans un verre plein ne fait pas déborder le verre ; en revanche elle diminue l'albédo et amplifie le réchauffement (rétroaction amplificatrice).

\textbf{Donnée chiffrée.} Le niveau moyen mondial des océans s'élève actuellement d'environ \cle{\SI{3,3}{mm.an^{-1}}} (mesures satellitaires altimétriques), soit une vingtaine de centimètres depuis 1900, avec une accélération nette depuis les années 1990. Selon les scénarios du GIEC, l'élévation atteindrait \num{0,3} à \SI{1}{\meter} d'ici 2100.

\begin{exemplebox}[Le littoral camerounais menacé]
Si le niveau de la mer s'élève de \SI{0,5}{\meter} d'ici 2100 (hypothèse médiane des scénarios intermédiaires), plus de \textbf{\num{200000} habitants du littoral camerounais} seraient directement menacés : quartiers bas de Douala (Deïdo, Bonabéri), ville de Kribi, villages de pêcheurs de Limbé et de la péninsule de Bakassi. S'y ajoutent la salinisation des nappes phréatiques côtières (eau potable contaminée), l'érosion des plages et la destruction des mangroves, qui protègent naturellement les côtes.
\end{exemplebox}

\courssub{Multiplication des événements extrêmes}

Le réchauffement modifie le cycle de l'eau et la circulation atmosphérique : une atmosphère plus chaude contient plus de vapeur d'eau (environ \SI{7}{\percent} de plus par degré), ce qui alimente des précipitations plus violentes, tandis que l'évaporation accrue aggrave les sécheresses dans les zones déjà arides.

\begin{itemize}
\item \textbf{Vagues de chaleur} : plus fréquentes, plus longues et plus intenses ; elles touchent désormais régulièrement le Sahel et les grandes villes (stress thermique urbain à Yaoundé et Douala).
\item \textbf{Inondations} : pluies extrêmes concentrées sur des épisodes courts ; les inondations récurrentes de Yaoundé (quartiers de la plaine du Mfoundi) et de Douala illustrent la vulnérabilité urbaine.
\item \textbf{Sécheresses} : allongement de la saison sèche dans l'Extrême-Nord ; avancée de la désertification et recul du lac Tchad, qui a perdu plus de \SI{90}{\percent} de sa superficie depuis les années 1960 (effets combinés du climat et des prélèvements humains).
\item \textbf{Cyclones plus intenses} : l'océan plus chaud fournit davantage d'énergie aux tempêtes tropicales ; leur intensité maximale et leurs pluies augmentent, même si leur nombre total ne croît pas nécessairement.
\end{itemize}

\courssub{Impacts sur l'agriculture camerounaise}

\textbf{Situation concrète.} Un planteur de cacao de la région du Centre observe que la saison des pluies, autrefois régulière, devient imprévisible : la floraison des cacaoyers est perturbée, les pourritures brunes des cabosses se propagent lors des pluies diluviennes, et les jeunes plants souffrent des épisodes secs prolongés.

\textbf{Mécanismes.} Les cultures dépendent étroitement de la température et de la pluviométrie :
\begin{itemize}
\item le \textbf{maïs}, culture vivrière majeure, voit ses rendements chuter lorsque la température dépasse son optimum pendant la floraison ou lorsque la pluie manque au remplissage des grains ; les modèles prévoient des baisses de rendement de l'ordre de \SIrange{10}{20}{\percent} en Afrique centrale à l'horizon 2050 dans les scénarios pessimistes ;
\item le \textbf{cacao}, qui exige une humidité élevée et stable, devient vulnérable : les zones favorables se déplacent en altitude ou vers le sud, et certaines zones actuelles (savane humide) pourraient devenir impropres à sa culture ;
\item les \textbf{calendriers culturaux} sont bouleversés : les semis traditionnellement calés sur l'arrivée des pluies deviennent aléatoires, ce qui accroît l'\cle{insécurité alimentaire}.
\end{itemize}

\begin{popfigure}
\begin{center}
\begin{tikzpicture}[scale=1.0, font=\small]
% Contour simplifié du Cameroun
\draw[thick, popdark, fill=popcream] 
(2.0,0.0) -- (3.2,0.4) -- (4.0,1.2) -- (4.6,2.0) -- (5.2,2.4) -- (5.0,3.4) -- (4.4,4.2) -- (3.6,4.6) -- (3.4,5.4) -- (3.8,6.2) -- (3.4,7.0) -- (2.8,7.6) -- (2.6,8.4) -- (2.2,9.0) -- (1.8,8.6) -- (1.4,7.8) -- (1.6,7.0) -- (1.2,6.2) -- (1.4,5.2) -- (1.0,4.4) -- (1.2,3.4) -- (0.8,2.6) -- (1.0,1.6) -- (0.6,0.8) -- (1.2,0.2) -- cycle;
% Villes
\fill[popink] (1.0,1.2) circle (0.06); \node[right, popink] at (1.05,1.2) {Douala};
\fill[popink] (2.2,1.8) circle (0.06); \node[right, popink] at (2.25,1.8) {Yaoundé};
\fill[popink] (2.4,7.4) circle (0.06); \node[right, popink] at (2.45,7.4) {Maroua};
% Zone littorale : montée des eaux
\draw[-{Stealth[length=3mm]}, very thick, popblue] (-0.8,0.6) -- (0.5,0.9);
\draw[-{Stealth[length=3mm]}, very thick, popblue] (-0.8,1.4) -- (0.6,1.5);
\node[popblue, align=center, anchor=east] at (-0.9,1.0) {Montée des eaux\\ littoral menacé};
% Zone Nord : sécheresse
\draw[-{Stealth[length=3mm]}, very thick, poporange] (2.2,10.2) -- (2.2,9.2);
\draw[-{Stealth[length=3mm]}, very thick, poporange] (3.2,10.0) -- (2.9,9.1);
\node[poporange, align=center] at (2.7,10.6) {Sécheresses et\\ désertification (Nord)};
% Zone cacao : déclin
\draw[-{Stealth[length=3mm]}, very thick, poppurple] (5.6,1.0) -- (4.2,1.6);
\node[poppurple, align=center, anchor=west] at (5.7,1.0) {Zones cacaoyères\\ en déclin};
% Zone paludisme
\draw[-{Stealth[length=3mm]}, very thick, popgreen] (5.8,3.6) -- (4.6,3.2);
\node[popgreen, align=center, anchor=west] at (5.9,3.6) {Extension du\\ paludisme};
\end{tikzpicture}
\end{center}
\legende{Carte simplifiée du Cameroun montrant les principales zones de vulnérabilité au réchauffement climatique : littoral exposé à la montée des eaux (Douala, Kribi, Limbé), Nord exposé aux sécheresses et à la désertification, zones de culture du cacao (Centre, Sud, Sud-Ouest) menacées de déclin, et extension possible des zones de transmission du paludisme.}
\end{popfigure}

\courssub{Risques sanitaires}

Le réchauffement climatique est aussi un problème de santé publique :

\begin{itemize}
\item \textbf{Expansion du paludisme} : le moustique vecteur, \emph{Anopheles}, se développe d'autant mieux que les températures sont douces et l'humidité présente. Le réchauffement permet au parasite \emph{Plasmodium} de compléter son cycle plus rapidement et étend les zones de transmission vers les hautes terres (monts Bamboutos, Adamaoua) autrefois épargnées.
\item \textbf{Maladies diarrhéiques} : les inondations contaminent les puits et les points d'eau par les matières fécales, favorisant les épidémies de choléra (foyers récurrents dans l'Extrême-Nord et à Douala) et de typhoïde.
\item \textbf{Stress thermique} : l'exposition prolongée à des chaleurs extrêmes provoque déshydratation, coups de chaleur et aggravation des maladies cardiovasculaires, surtout chez les personnes âgées, les enfants et les travailleurs exposés (agriculteurs, ouvriers).
\item \textbf{Malnutrition} : la baisse des rendements agricoles et la hausse des prix alimentaires aggravent la sous-nutrition des populations vulnérables.
\end{itemize}

\courssub{Perte de biodiversité}

\textbf{Observation.} Les scientifiques constatent que de nombreuses espèces se déplacent vers les pôles ou en altitude, en moyenne de plusieurs kilomètres par décennie, pour suivre les températures qui leur conviennent.

\begin{itemize}
\item \textbf{Déplacement et extinction d'espèces} : les espèces incapables de migrer assez vite ou de s'adapter s'éteignent localement. Les forêts camerounaises, deuxième massif forestier d'Afrique, abritent des espèces endémiques (gorilles de Cross River, perroquets gris du Gabon) dont les habitats se fragmentent.
\item \textbf{Blanchiment des coraux} : lorsque la température de l'eau dépasse durablement un seuil, les coraux expulsent les algues symbiotiques (\cle{zooxanthelles}) qui les nourrissent et les colorent ; ils blanchissent puis meurent si le stress persiste, entraînant l'effondrement d'écosystèmes entiers.
\item \textbf{Disparition des habitats montagnards} : sur le Mont Cameroun (\SI{4095}{\meter}) et les monts Bamboutos, la ceinture de forêt de montagne et les prairies d'altitude remontent et rétrécissent ; les espèces confinées au sommet n'ont « nulle part où monter » et disparaissent.
\end{itemize}

\begin{savaistu}[Les mangroves, championnes du stockage du carbone]
Les mangroves du Cameroun (estuaire du Wouri, Rio del Rey, Campo) stockent jusqu'à \textbf{4 fois plus de carbone par hectare} que les forêts terrestres, grâce à l'accumulation de matière organique dans leurs sols gorgés d'eau. Les protéger est donc un double levier : \textbf{atténuation} (conserver ce réservoir de carbone) et \textbf{adaptation} (elles freinent l'érosion et protègent les côtes de la montée des eaux et des tempêtes). Leur destruction pour le bois de fumage du poisson ou l'urbanisation aggrave à la fois les émissions et la vulnérabilité du littoral.
\end{savaistu}

\begin{exoresolu}[Exploiter des données climatiques]
\textbf{Énoncé.} Le niveau moyen des océans s'élève actuellement de \SI{3,3}{mm.an^{-1}}. a) En supposant ce rythme constant, calculer l'élévation entre 2024 et 2100. b) Comparer ce résultat à la projection de \SI{0,5}{\meter} du scénario intermédiaire et expliquer l'écart. c) Citer deux conséquences pour le littoral camerounais.
\tcblower
\textbf{Solution.}
a) Durée : $2100 - 2024 = 76$ ans. Élévation :
\[
\Delta h = \SI{3,3}{mm.an^{-1}} \times 76 \ \text{ans} = \SI{250,8}{\milli\meter} \approx \SI{0,25}{\meter}.
\]
b) La valeur obtenue (\SI{0,25}{\meter}) est deux fois inférieure à la projection de \SI{0,5}{\meter}. L'écart s'explique par le fait que le rythme n'est \textbf{pas constant} : la fonte des calottes glaciaires et la dilatation thermique s'accélèrent avec le réchauffement (rétroactions amplificatrices). Le calcul à rythme constant sous-estime donc la réalité.
c) Conséquences pour le littoral camerounais : submersion des quartiers bas de Douala et des villages de pêcheurs, salinisation des nappes d'eau douce côtières, érosion des plages de Kribi et Limbé, destruction des mangroves protectrices, déplacement de plus de \num{200000} personnes.
\end{exoresolu}

\begin{aretenir}[L'essentiel de la section]
\begin{itemize}
\item La température moyenne mondiale a augmenté d'environ \textbf{\SI{1,1}{\celsius} depuis 1850} ; selon les scénarios du GIEC, elle atteindrait \SI{1,8}{\celsius} (SSP1--2.6) à \SI{4,4}{\celsius} (SSP5--8.5) en 2100.
\item Les \textbf{rétroactions amplificatrices} (fonte des glaces $\rightarrow$ albédo en baisse $\rightarrow$ absorption solaire en hausse) accélèrent le réchauffement.
\item Le niveau des océans monte d'environ \textbf{\SI{3,3}{mm.an^{-1}}} (fonte des glaces continentales + dilatation thermique), menaçant le littoral camerounais.
\item Les \textbf{événements extrêmes} (vagues de chaleur, inondations, sécheresses, cyclones) deviennent plus probables ou plus intenses — sans qu'on puisse attribuer un événement isolé au seul climat.
\item Conséquences majeures au Cameroun : baisse des rendements du maïs et du cacao, calendriers culturaux perturbés, extension du paludisme et des maladies diarrhéiques, stress thermique, déplacement d'espèces et disparition d'habitats.
\end{itemize}
\end{aretenir}

\begin{exercice}[Pour t'entraîner]
1. Définir le réchauffement climatique et citer la valeur de l'augmentation observée depuis 1850.
2. Expliquer, à l'aide de la notion d'albédo, pourquoi la fonte de la banquise arctique amplifie le réchauffement alors qu'elle ne fait pas monter le niveau de la mer.
3. Un élève affirme : « L'inondation de Yaoundé la semaine dernière prouve le réchauffement climatique. » Corriger cette affirmation en une phrase rigoureuse.
4. Expliquer pourquoi le paludisme pourrait s'étendre vers les hautes terres de l'Ouest camerounais.
\end{exercice}

\begin{corrige}[Exercice « Pour t'entraîner »]
1. Augmentation durable de la température moyenne de l'atmosphère et des océans due à l'accumulation des GES d'origine humaine ; environ \SI{1,1}{\celsius} depuis 1850--1900.
2. La banquise flotte déjà sur l'océan : sa fonte ne modifie pas le volume d'eau (comme un glaçon dans un verre). En revanche, la glace blanche réfléchit fortement le rayonnement solaire (albédo élevé) ; remplacée par une eau sombre (albédo faible), la surface absorbe davantage d'énergie, se réchauffe et accélère la fonte : c'est une rétroaction amplificatrice.
3. Formulation correcte : « Le réchauffement climatique augmente la probabilité et l'intensité des pluies extrêmes à l'origine de telles inondations » — un événement isolé ne peut être attribué directement au changement climatique.
4. Le réchauffement élève les températures en altitude, créant des conditions thermiques favorables au développement des moustiques \emph{Anopheles} et à l'achèvement du cycle du parasite \emph{Plasmodium} dans des zones de montagne autrefois trop froides.
\end{corrige}

\coursec{Moyens d'action : atténuation et adaptation}

Nous venons de voir que le renforcement de l'effet de serre entraîne des conséquences graves : élévation des températures, montée du niveau des mers, multiplication des sécheresses dans l'Adamaoua et l'Extrême-Nord, pluies extrêmes et inondations à Douala, recul des cultures et paludisme qui gagne les hauts plateaux de l'Ouest. Face à un tel diagnostic, une question s'impose naturellement : \emph{que pouvons-nous faire ?} C'est toute l'ambition de cette section : montrer qu'il existe deux grandes familles de réponses — \cle{atténuer} le problème à sa source et \cle{s'adapter} à ses effets déjà engagés — et que ces réponses se déclinent à toutes les échelles, du geste quotidien d'un élève de Première jusqu'aux accords internationaux.

\courssub{Deux stratégies complémentaires : atténuation et adaptation}

\textbf{Une situation pour comprendre.} Imagine ta maison pendant la saison des pluies : le toit fuit et l'eau s'infiltre. Deux réactions sont possibles. La première consiste à \emph{réparer le toit} : on traite la cause de la fuite. La seconde consiste à \emph{placer des bassines, surélever les meubles et creuser un canal d'évacuation} : on ne supprime pas la fuite, mais on réduit les dégâts qu'elle provoque. La lutte contre le réchauffement climatique repose exactement sur ces deux logiques complémentaires.

\begin{definition}[Atténuation et adaptation]
L'\cle{atténuation} désigne l'ensemble des actions humaines visant à \textbf{réduire les sources} d'émissions de gaz à effet de serre (GES) ou à \textbf{augmenter les puits} qui les absorbent (forêts, sols, océans). Elle agit sur les \emph{causes} du réchauffement.

L'\cle{adaptation} désigne l'ensemble des ajustements des systèmes naturels et humains visant à \textbf{réduire la vulnérabilité} face aux effets du changement climatique déjà en cours ou prévisibles : variétés de cultures résistantes à la sécheresse, digues contre la montée des eaux, systèmes d'alerte précoce, réorganisation des calendriers agricoles.
\end{definition}

\begin{attention}[Ne pas opposer les deux stratégies]
Une erreur fréquente au Probatoire consiste à présenter l'adaptation comme une alternative à l'atténuation. C'est faux : même si toutes les émissions s'arrêtaient demain, l'inertie du système climatique maintiendrait le réchauffement pendant des décennies. L'adaptation est donc \emph{indispensable}, mais sans atténuation, l'ampleur du réchauffement dépasserait nos capacités d'adaptation. Les deux stratégies sont \textbf{complémentaires et simultanées}.
\end{attention}

\courssub{Les gestes individuels : agir à son échelle}

\textbf{Observation.} Une famille de Yaoundé utilise quotidiennement une voiture à essence, consomme de la viande à chaque repas, laisse les ampoules allumées et jette ses déchets organiques à la poubelle. Chacune de ces habitudes émet des GES : le carburant brûlé dégage du \ce{CO2}, l'élevage produit du méthane (\ce{CH4}), l'électricité consommée provient en partie de centrales thermiques, et les déchets organiques enfouis fermentent en dégageant du méthane. Changer ces habitudes, c'est atténuer.

\begin{exemplebox}[Cinq familles de gestes efficaces]
\begin{itemize}
\item \textbf{Mobilité douce} : marcher, faire du vélo, utiliser les transports en commun ou le covoiturage plutôt que la voiture individuelle. Le transport est l'un des postes d'émission qui croît le plus vite dans les villes camerounaises.
\item \textbf{Isolation et sobriété du logement} : toitures claires, ventilation naturelle, ombrage végétal, ampoules LED, extinction des appareils en veille. Moins de climatisation et moins d'électricité consommée, c'est moins de \ce{CO2} émis.
\item \textbf{Alimentation moins carnée et locale} : réduire la fréquence des repas de viande de bœuf (l'élevage bovin émet beaucoup de méthane) et privilégier les produits locaux de saison (manioc, plantain, légumineuses, poisson local) limite les émissions liées à l'élevage et au transport des denrées.
\item \textbf{Tri et compostage} : composter les épluchures et résidus de cuisine évite leur fermentation anaérobie en décharge (source de \ce{CH4}) et produit un engrais naturel pour le jardin ou la ferme.
\item \textbf{Économie d'énergie} : cuisinières améliorées (foyers à bois économes) qui réduisent la consommation de bois de chauffe, donc la déforestation et les émissions.
\end{itemize}
\end{exemplebox}

\begin{exemplebox}[Exemple chiffré : le vélo contre la voiture]
Une voiture essence émet en moyenne environ 120 g de \ce{CO2} par kilomètre parcouru (émissions directes à l'échappement). Remplacer chaque jour 10 km de trajet en voiture par 10 km de vélo évite :
\[ 10\ \text{km/j} \times 0{,}120\ \text{kg/km} \times 365\ \text{j} \approx 438\ \text{kg/an} \]
Si l'on intègre les émissions indirectes (extraction, raffinage, transport du carburant, fabrication du véhicule), le facteur d'émission monte à environ 300 g/km, soit \textbf{environ 1,1 tonne de \ce{CO2} évitée par an et par personne} pour un usage quotidien soutenu. À l'échelle d'une ville entière, l'effet cumulé devient considérable.
\end{exemplebox}

\begin{methode}[Élaborer un plan d'action personnel « 5 gestes pour le climat »]
\begin{enumerate}
\item \textbf{Identifier} : liste tes principales sources d'émissions (transport, alimentation, énergie, déchets) pendant une semaine d'observation.
\item \textbf{Quantifier} : estime pour chaque geste envisagé les kilogrammes de \ce{CO2} évités par an, à l'aide de facteurs d'émission simples (ex. 0,12 kg/km en voiture essence, environ 0,5 kg par repas végétarien substitué à un repas de bœuf).
\item \textbf{Planifier} : choisis 5 gestes réalistes, fixe une date de début et un objectif chiffré (ex. « 500 kg de \ce{CO2} évités en un an »).
\item \textbf{Suivre} : tiens un tableau mensuel de tes actions et ajuste ton plan. Partage tes résultats en classe ou en famille : l'exemplarité est un levier collectif.
\end{enumerate}
\end{methode}

\begin{attention}[Le « greenwashing »]
Certaines entreprises affichent des produits « verts », « naturels » ou « climatiquement neutres » sans justification réelle : c'est le \cle{greenwashing} (écoblanchiment). Piège classique : un emballage vert décoré de feuilles ne prouve rien. Réflexe à acquérir : vérifier l'existence de \textbf{labels reconnus} (labels écologiques officiels, certification forestière FSC ou PEFC pour le bois), la \textbf{traçabilité} du produit et des données chiffrées vérifiables plutôt que des slogans vagues.
\end{attention}

\courssub{Les politiques publiques : agir à l'échelle d'un pays}

Les gestes individuels sont nécessaires mais insuffisants : les choix de consommation dépendent des infrastructures, des prix et des lois. D'où le rôle central des politiques publiques.

\begin{exemplebox}[Principaux instruments de politique climatique]
\begin{itemize}
\item \textbf{Taxe carbone} : elle fait payer chaque tonne de \ce{CO2} émise, ce qui renchérit les énergies fossiles et oriente consommateurs et entreprises vers des alternatives propres. Limite connue : si elle est mal conçue, elle pèse sur les ménages modestes ; elle doit donc être accompagnée de compensations.
\item \textbf{Subventions aux énergies renouvelables} : soutien financier au solaire photovoltaïque, à l'hydroélectricité (atout majeur du Cameroun avec les barrages de la Sanaga), à la biomasse durable, afin de rendre ces filières compétitives.
\item \textbf{Normes d'efficacité énergétique} : obligations de performance pour les bâtiments neufs, les appareils électriques et les véhicules ; interdiction progressive des équipements les plus énergivores.
\item \textbf{Reboisement et protection des forêts} : programmes de plantation, lutte contre l'exploitation forestière illégale et la déforestation liée à l'extension agricole.
\item \textbf{Plan National d'Adaptation (PNA)} : le Cameroun a élaboré un PNA qui identifie les secteurs vulnérables (agriculture, eau, santé, zones côtières) et programme des mesures concrètes : variétés de maïs et de sorgho tolérantes à la sécheresse dans le Grand Nord, protection des mangroves de l'estuaire du Wouri contre l'érosion et la montée des eaux, renforcement de la surveillance épidémiologique du paludisme.
\end{itemize}
\end{exemplebox}

\begin{popfigure}
\begin{center}
\begin{tikzpicture}[scale=1]
% Pyramide en trois étages
\fill[popgreenL] (-4,0) -- (4,0) -- (2.6,1.8) -- (-2.6,1.8) -- cycle;
\fill[popblueL] (-2.6,1.8) -- (2.6,1.8) -- (1.3,3.6) -- (-1.3,3.6) -- cycle;
\fill[poporangeL] (-1.3,3.6) -- (1.3,3.6) -- (0,5.4) -- cycle;
\draw[popdark,thick] (-4,0) -- (4,0) -- (0,5.4) -- cycle;
\draw[popdark] (-2.6,1.8) -- (2.6,1.8);
\draw[popdark] (-1.3,3.6) -- (1.3,3.6);
\node[popdark,font=\small\bfseries,align=center] at (0,0.9) {Gestes quotidiens\\ (individus, familles, élèves)};
\node[popdark,font=\small\bfseries,align=center] at (0,2.7) {Politiques locales et nationales\\ (communes, État, PNA, normes)};
\node[popdark,font=\small\bfseries,align=center,text width=6cm] at (0,4.35) {Accords internationaux\\ (Accord de Paris, CDN)};
% Flèches d'interaction
\draw[-{Stealth},popmuted,thick] (4.3,0.6) -- (4.3,4.8) node[midway,right,font=\small,align=left,text width=3.6cm] {les échelles se renforcent mutuellement};
\end{tikzpicture}
\end{center}
\legende{Pyramide des actions climatiques : la base repose sur les gestes quotidiens de chacun, soutenus par les politiques locales et nationales, elles-mêmes encadrées par les accords internationaux. Aucun niveau ne peut se substituer aux autres.}
\end{popfigure}

\courssub{Les puits de carbone : forêts et sols, alliés du climat}

\textbf{Observation et hypothèse.} Le Cameroun abrite une partie du bassin du Congo, deuxième massif forestier tropical du monde. On observe qu'une forêt dense en croissance « pompe » du \ce{CO2} atmosphérique. Hypothèse : la végétation stocke le carbone de l'air dans sa biomasse et dans le sol.

\textbf{Données et interprétation.} La photosynthèse fixe le carbone atmosphérique dans la matière organique :
\[ \ce{6CO2 + 6H2O ->[h\nu] C6H12O6 + 6O2} \]
Les mesures montrent qu'un hectare de forêt tropicale humide peut stocker de l'ordre de 150 à 250 tonnes de carbone, dont une part importante dans le sol. Conclusion : une forêt est un \cle{puits de carbone} ; la détruire par la coupe ou le brûlis transforme ce puits en \emph{source} (le carbone stocké repart dans l'atmosphère sous forme de \ce{CO2}).

\begin{exemplebox}[Gérer et restaurer les puits de carbone]
\begin{itemize}
\item \textbf{Gestion forestière durable} : exploitation sélective avec reboisement, certification du bois, aires protégées (parcs nationaux de Lobéké, de la Bénoué, de Waza).
\item \textbf{Agroforesterie} : associer arbres et cultures sur une même parcelle (cacaoyers sous ombrage, fruitiers dans les champs de vivriers) : on stocke du carbone tout en maintenant la production agricole et en protégeant les sols de l'érosion.
\item \textbf{Restauration des terres dégradées} : reboisement des savanes surexploitées, fixation des dunes, techniques anti-érosives (cordons pierreux, bandes végétales) qui reconstituent la matière organique des sols — or les sols constituent, avec les océans, l'un des plus grands réservoirs de carbone de la planète.
\end{itemize}
\end{exemplebox}

\begin{savaistu}[Le programme « 1 milliard d'arbres »]
En 2022, le Cameroun a lancé un ambitieux programme national de reboisement visant à planter \textbf{un milliard d'arbres} et à restaurer environ \textbf{12 millions d'hectares} de terres dégradées à l'horizon 2030. Ce programme s'inscrit dans l'initiative panafricaine AFR100 et contribue directement aux engagements climatiques du pays : chaque hectare restauré redevient un puits de carbone tout en protégeant les sols, l'eau et la biodiversité.
\end{savaistu}

\courssub{Captage et stockage du carbone (CSC) : promesses et limites}

\begin{definition}[Captage et stockage du carbone]
Le \cle{CSC} (en anglais CCS) est une technologie qui consiste à \textbf{capter} le \ce{CO2} à la sortie des grandes installations industrielles (cimenteries, centrales thermiques), à le \textbf{comprimer}, à le \textbf{transporter}, puis à l'\textbf{injecter} dans des formations géologiques profondes et étanches (anciens gisements de gaz ou de pétrole, aquifères salins profonds) où il resterait piégé pendant des millénaires.
\end{definition}

\textbf{État d'avancement et limites.} À l'échelle mondiale, quelques dizaines d'installations fonctionnent, mais la capacité totale captée reste très inférieure à 1 \% des émissions mondiales. Les limites sont réelles :
\begin{itemize}
\item \textbf{coût élevé} (captage et compression consomment eux-mêmes de l'énergie, ce qui réduit le bénéfice net) ;
\item \textbf{risques de fuites} à long terme et nécessité d'une surveillance géologique permanente ;
\item \textbf{risque moral} : le CSC peut servir d'excuse pour continuer à brûler des énergies fossiles au lieu d'en sortir.
\end{itemize}
Le CSC est donc un \emph{complément} possible pour les industries difficiles à décarboner, jamais un substitut à la réduction des émissions.

\courssub{Les engagements internationaux : de l'Accord de Paris aux CDN du Cameroun}

\textbf{Pourquoi un cadre international ?} L'atmosphère ne connaît pas les frontières : le \ce{CO2} émis en Asie réchauffe l'Afrique, et inversement. Aucun pays ne peut donc résoudre seul le problème. C'est le sens de la Convention-cadre des Nations unies sur les changements climatiques (CCNUCC, 1992) puis de l'\textbf{Accord de Paris} (2015), qui fixe l'objectif de contenir le réchauffement « nettement en dessous de 2 °C » par rapport à l'ère préindustrielle, en poursuivant les efforts vers 1,5 °C. Chaque pays traduit cet objectif par sa \cle{CDN} (Contribution Déterminée au niveau National), révisée tous les cinq ans.

\begin{propriete}[Principe de responsabilité commune mais différenciée (CBDR)]
Tous les États ont une responsabilité \emph{commune} dans la lutte contre le changement climatique, mais \emph{différenciée} selon leur contribution historique aux émissions et leur niveau de développement : les pays développés, responsables de l'essentiel des émissions accumulées depuis la révolution industrielle, doivent agir en premier, réduire davantage et \textbf{soutenir financièrement et technologiquement} les pays en développement (objectif international de 100 milliards de dollars par an de financement climat). Ce principe fonde la justice climatique : l'Afrique, qui émet moins de 4 \% des GES mondiaux, subit pourtant des impacts disproportionnés.
\end{propriete}

\begin{exemplebox}[La CDN du Cameroun]
Dans sa Contribution Déterminée au niveau National actualisée, le Cameroun s'est engagé à \textbf{réduire ses émissions de GES de 35 \% d'ici 2030 par rapport au scénario de référence} (« business as usual », c'est-à-dire la trajectoire sans politique climatique). Une partie de cet engagement est conditionnée à l'appui financier et technique international, conformément au principe CBDR. Les leviers prioritaires : énergies renouvelables (hydroélectricité, solaire), gestion durable des forêts, agriculture durable et efficacité énergétique.
\end{exemplebox}

\begin{popfigure}
\begin{center}
\begin{tikzpicture}
\begin{axis}[
width=0.95\linewidth, height=7cm,
xlabel={Ann\'ee}, ylabel={\'Emissions de GES (Mt CO$_2$ \'eq.)},
xmin=2015, xmax=2035, ymin=0, ymax=160,
xtick={2015,2020,2025,2030,2035},
ytick={0,40,80,120,160},
legend style={at={(0.03,0.97)},anchor=north west,font=\small},
grid=major, grid style={popline!40},
axis line style={popdark},
]
% Scénario business as usual : 80 Mt en 2015, +3%/an environ
\addplot[popcurve,poporange,very thick,domain=2015:2030,samples=60]{80*1.03^(x-2015)};
\addlegendentry{Sc\'enario de r\'ef\'erence (\og business as usual \fg)}
% Scénario CDN : -35% en 2030 par rapport au BAU
\addplot[popcurve,popgreen,very thick,domain=2015:2030,samples=60]{80*1.03^(x-2015)*(1-0.35*(x-2015)/15)};
\addlegendentry{Sc\'enario \og CDN respect\'ee \fg\ ($-35\,\%$ en 2030)}
% Pointillé objectif 2030
\addplot[popblue,dashed,thick] coordinates {(2030,0) (2030,124)};
\node[popblue,font=\small,anchor=south west] at (axis cs:2030.2,128) {Objectif 2030};
\end{axis}
\end{tikzpicture}
\end{center}
\legende{Trajectoires d'émissions de gaz à effet de serre du Cameroun (valeurs indicatives, en millions de tonnes d'équivalent \ce{CO2}). En orange : poursuite des tendances actuelles. En vert : trajectoire compatible avec la CDN, soit 35 \% d'émissions en moins que le scénario de référence en 2030. L'écart entre les deux courbes représente l'effort d'atténuation.}
\end{popfigure}

\begin{exoresolu}[Atténuation ou adaptation ?]
\textbf{Énoncé.} Pour chacune des actions suivantes menées au Cameroun, indique s'il s'agit d'une mesure d'atténuation ou d'adaptation, en justifiant : (a) plantation de mangroves pour fixer le carbone et protéger les côtes ; (b) distribution de semences de sorgho tolérant à la sécheresse dans l'Extrême-Nord ; (c) construction d'une centrale solaire à Guider ; (d) mise en place d'un système d'alerte précoce aux crues à Douala ; (e) taxe sur les véhicules très polluants.
\tcblower
\textbf{Solution détaillée.}
\begin{itemize}
\item (a) \textbf{Les deux} : la mangrove stocke du carbone (atténuation, augmentation d'un puits) \emph{et} protège les côtes de l'érosion et des ondes de tempête (adaptation). C'est un exemple de mesure à double bénéfice.
\item (b) \textbf{Adaptation} : on ne réduit pas les émissions, on réduit la vulnérabilité de l'agriculture à la sécheresse.
\item (c) \textbf{Atténuation} : l'énergie solaire remplace une production thermique fossile, donc réduit les sources d'émissions.
\item (d) \textbf{Adaptation} : l'alerte précoce réduit les dégâts humains et matériels des inondations sans agir sur les émissions.
\item (e) \textbf{Atténuation} : la taxe dissuade l'usage de véhicules émetteurs de \ce{CO2}, donc réduit les sources.
\end{itemize}
\textbf{Méthode à retenir} : demande-toi toujours « cette action agit-elle sur les \emph{émissions} (atténuation) ou sur les \emph{dégâts} (adaptation) ? »
\end{exoresolu}

\begin{exercice}[Ton plan climat personnel et national]
1. Cite trois gestes quotidiens d'atténuation et estime, pour l'un d'eux, les kilogrammes de \ce{CO2} évités par an (montre ton calcul).
2. Explique pourquoi le reboisement est à la fois une mesure d'atténuation et, parfois, d'adaptation.
3. Rédige un court paragraphe (10 lignes maximum) expliquant à un camarade ce qu'est la CDN du Cameroun et pourquoi le principe CBDR justifie que les pays riches soutiennent cet effort.
\end{exercice}

\begin{corrige}[Exercice « Ton plan climat personnel et national »]
1. Exemples attendus : vélo ou marche au lieu de la voiture, ampoules LED et extinction des veilles, compostage des déchets organiques, repas moins carnés, cuisinière améliorée. Calcul type : 10 km/j évités en voiture $\times$ 0,12 kg/km $\times$ 300 j = 360 kg de \ce{CO2} évités par an. Toute estimation cohérente avec un facteur d'émission explicite est acceptée.
2. Le reboisement est une atténuation car les arbres fixent le \ce{CO2} atmosphérique par photosynthèse (augmentation d'un puits). Il peut aussi être une adaptation : les arbres limitent l'érosion, régulent le régime des eaux, ombragent les cultures et brisent le vent, ce qui réduit la vulnérabilité des terres et des populations aux sécheresses et aux pluies extrêmes.
3. Éléments attendus : la CDN est l'engagement national du Cameroun dans le cadre de l'Accord de Paris ; elle vise une réduction de 35 \% des émissions en 2030 par rapport au scénario de référence ; le principe CBDR rappelle que les pays développés, responsables historiques de l'essentiel des émissions accumulées et disposant de plus de moyens, doivent agir en premier et financer les efforts des pays en développement comme le Cameroun, dont la part dans les émissions mondiales est très faible alors que les impacts subis sont lourds.
\end{corrige}

\begin{aretenir}[L'essentiel de la section]
\begin{itemize}
\item \textbf{Atténuation} = réduire les sources et augmenter les puits de GES ; \textbf{adaptation} = réduire la vulnérabilité aux effets du réchauffement. Les deux sont complémentaires.
\item L'action climatique s'organise en pyramide : gestes individuels (mobilité douce, sobriété énergétique, alimentation, compostage), politiques publiques (taxe carbone, normes, reboisement, PNA), accords internationaux (Accord de Paris, CDN).
\item Forêts et sols sont de puissants puits de carbone : leur gestion durable, l'agroforesterie et la restauration des terres dégradées sont des leviers majeurs pour le Cameroun (programme « 1 milliard d'arbres »).
\item Le CSC est une technologie complémentaire coûteuse et limitée, non une solution miracle ; se méfier du greenwashing.
\item Le Cameroun s'est engagé à réduire ses émissions de 35 \% d'ici 2030 par rapport au scénario de référence, dans le cadre du principe de responsabilité commune mais différenciée.
\end{itemize}
\end{aretenir}

\coursec{Synthèse, modélisation simple et engagement citoyen}

Nous avons vu dans la section précédente que la lutte contre le changement climatique repose sur deux piliers : l'\cle{atténuation} (réduire les émissions de gaz à effet de serre) et l'\cle{adaptation} (se préparer aux impacts inévitables). Mais une question demeure : comment passer des grands discours internationaux — Accords de Paris, contributions nationales du Cameroun — à des actions concrètes, mesurables, à notre échelle ? C'est tout l'objet de cette dernière section : apprendre à \textbf{mesurer} pour mieux \textbf{agir}, en construisant ensemble un modèle simple du bilan carbone de la classe, puis en nous engageant collectivement. Car on ne peut réduire efficacement que ce que l'on sait compter.

\courssub{Récapitulatif : la chaîne complète, de l'activité humaine à l'action}

Faisons d'abord une pause pour assembler toutes les pièces du puzzle étudié dans cette leçon. Le raisonnement complet se déroule comme une chaîne de causes et de conséquences :

\begin{enumerate}
\item Les \textbf{activités humaines} (transport, production d'électricité, agriculture, élevage, déforestation, industrie, gestion des déchets) émettent des \textbf{gaz à effet de serre} : \ce{CO2}, méthane \ce{CH4}, protoxyde d'azote \ce{N2O}.
\item Ces émissions augmentent la \textbf{concentration atmosphérique} en GES (le \ce{CO2} est passé d'environ 280 ppm avant l'ère industrielle à plus de 420 ppm aujourd'hui).
\item L'augmentation de concentration renforce le \textbf{forçage radiatif} : l'atmosphère absorbe davantage le rayonnement infrarouge émis par la Terre et le réémet vers la surface.
\item Il en résulte un \textbf{réchauffement} de la basse atmosphère et des océans (environ $+1,1$~°C depuis l'ère préindustrielle).
\item Ce réchauffement engendre des \textbf{impacts} : élévation du niveau des mers, modification des régimes de pluie (crucial pour l'agriculture camerounaise), multiplication des sécheresses et des inondations, déplacement des zones de paludisme, fonte des derniers glaciers du Mont Cameroun et des montagnes tropicales.
\item Face à ces impacts, des \textbf{actions} sont possibles : atténuer (réduire les émissions) et s'adapter (réduire la vulnérabilité).
\end{enumerate}

\begin{aretenir}[La chaîne causale à retenir]
\textbf{Activités} $\rightarrow$ \textbf{Émissions de GES} $\rightarrow$ \textbf{Concentration atmosphérique} $\rightarrow$ \textbf{Forçage radiatif} $\rightarrow$ \textbf{Réchauffement} $\rightarrow$ \textbf{Impacts} $\rightarrow$ \textbf{Actions (atténuation + adaptation)}. Rompre la chaîne le plus en amont possible (au niveau des émissions) est la stratégie la plus efficace.
\end{aretenir}

\courssub{Mesurer pour agir : l'empreinte carbone}

\textbf{Situation concrète.} Le proviseur d'un lycée de Yaoundé souhaite « verdir » son établissement. Faut-il d'abord installer des lampes LED, interdire les sachets plastiques, ou réorganiser le transport des élèves ? Sans chiffres, impossible de trancher : chacun défend son idée favorite. Il faut d'abord \emph{mesurer} d'où viennent les émissions.

\begin{definition}[Empreinte carbone]
L'\cle{empreinte carbone} d'une entité (individu, classe, école, pays) est la \textbf{quantité totale de gaz à effet de serre émise directement ou indirectement} par ses activités, exprimée en \textbf{équivalent \ce{CO2}} (noté \ce{CO2}éq). L'équivalent \ce{CO2} convertit tous les GES en une unité commune grâce à leur pouvoir de réchauffement global : 1 kg de \ce{CH4} compte pour environ 28 kg \ce{CO2}éq sur 100 ans.
\end{definition}

\begin{methode}[Calcul d'une empreinte carbone simplifiée]
\begin{enumerate}
\item \textbf{Découper} l'activité en postes : transport, énergie (électricité, bois, gaz), alimentation, déchets.
\item \textbf{Collecter} les données d'activité pour chaque poste : kilomètres parcourus, kWh consommés, repas servis, kilogrammes de déchets.
\item \textbf{Multiplier} chaque donnée d'activité par son \textbf{facteur d'émission} (en kg \ce{CO2}éq par unité), issu de bases de référence (ADEME, GIEC/IPCC).
\item \textbf{Additionner} : $\text{Empreinte} = \sum_i (\text{activité}_i \times \text{facteur d'émission}_i)$.
\item \textbf{Analyser} les résultats par poste pour identifier les priorités.
\end{enumerate}
\end{methode}

\begin{attention}[Les facteurs d'émission dépendent du contexte local]
Un même kWh d'électricité n'a pas la même empreinte partout : tout dépend du \textbf{mix électrique}. Au Cameroun, environ 70~\% de l'électricité provient de l'hydroélectricité (barrages de Songloulou, Édéa, Lagdo), source à \textbf{faible intensité carbone} ; le facteur d'émission y est bien plus faible qu'en Europe ou dans les pays producteurs d'électricité au charbon. En revanche, les groupes électrogènes diesel, très utilisés lors des délestages, ont une intensité carbone élevée. Utiliser un facteur d'émission étranger sans vérification est une erreur fréquente.
\end{attention}

\courssub{Atelier : le bilan carbone de la classe}

Construisons maintenant le modèle simplifié. La classe collecte ses données pendant une semaine (carnet de bord des trajets, relevé du compteur de la salle, enquête cantine), puis extrapole à l'année.

\begin{popfigure}
\begin{center}
\begin{tikzpicture}
\node[fill=popblueL, rounded corners=2pt, font=\small\bfseries, minimum width=\linewidth, inner sep=4pt] (t0) at (0,0) {Bilan carbone simplifié d'une classe de 40 élèves (hypothèses annuelles)};
\node[anchor=north west, font=\small, text width=0.94\linewidth, inner sep=2pt] at (-0.49\linewidth,-0.5) {%
\textbf{Poste Transport} : $E_t = N \times d \times n \times f_t$ \quad ($N$ élèves motorisés, $d$ distance aller-retour, $n$ jours, $f_t$ facteur)\\[2pt]
Ex. : $20 \times 8~\text{km} \times 180 \times 0{,}12~\text{kg \ce{CO2}éq/km} \approx 3{,}5$ t \ce{CO2}éq\\[4pt]
\textbf{Poste Énergie} : $E_e = C \times f_e$ \quad ($C$ consommation en kWh, $f_e$ facteur du mix)\\[2pt]
Ex. : $2\,500~\text{kWh} \times 0{,}05~\text{kg \ce{CO2}éq/kWh (hydro)} \approx 0{,}13$ t \ce{CO2}éq\\[4pt]
\textbf{Poste Alimentation} : $E_a = R \times f_a$ \quad ($R$ repas, $f_a$ selon menu : viande $\approx 2$ kg, végétarien $\approx 0{,}5$ kg)\\[2pt]
Ex. : $6\,000$ repas mixtes $\times 1{,}2~\text{kg} \approx 7{,}2$ t \ce{CO2}éq (cantine incluse)\\[4pt]
\textbf{Poste Déchets} : $E_d = M \times f_d$ \quad ($M$ masse de déchets, $f_d \approx 0{,}5$ kg \ce{CO2}éq/kg en décharge)\\[2pt]
Ex. : $800~\text{kg} \times 0{,}5 \approx 0{,}4$ t \ce{CO2}éq\\[4pt]
\textbf{Total} : $E = E_t + E_e + E_a + E_d$};
\end{tikzpicture}
\end{center}
\legende{Tableau récapitulatif des postes d'émission de la classe, avec la formule de calcul associée à chaque poste (donnée d'activité $\times$ facteur d'émission) et un exemple chiffré. Les valeurs sont des ordres de grandeur pédagogiques.}
\end{popfigure}

\begin{propriete}[Principe de proportionnalité (analyse de Pareto)]
La réduction la plus efficace cible \textbf{le poste d'émission le plus important} : réduire de 10~\% un poste qui pèse 7 t \ce{CO2}éq économise 0,7 t, tandis que réduire de 50~\% un poste de 0,13 t n'économise que 0,065 t. On agit donc d'abord là où le « gisement » est le plus gros (souvent : alimentation et transport), sans négliger les petits gestes qui entretiennent la dynamique collective.
\end{propriete}

\begin{exemplebox}[Exemple chiffré : l'effet du covoiturage]
Une classe de 40 élèves émet au total \textbf{2 t \ce{CO2}éq/an} (modèle simplifié centré sur les déplacements quotidiens). Si le covoiturage est adopté pour \textbf{50~\% des trajets motorisés}, les kilomètres parcourus en voiture individuelle diminuent fortement : la réduction atteint environ \textbf{0,3 t \ce{CO2}éq/an}, soit \textbf{15~\%} de l'empreinte totale de la classe — avec une seule mesure, sans aucun investissement matériel.
\end{exemplebox}

\courssub{De la mesure à l'engagement : la Charte climat de la classe}

Une fois les leviers identifiés (covoiturage ou transport en commun, extinction des lampes et ventilateurs inutiles, menus végétariens hebdomadaires à la cantine, tri des déchets plastiques), la classe rédige collectivement une \textbf{Charte climat}. Une bonne charte comporte :

\begin{itemize}
\item des \textbf{objectifs chiffrés et datés} : par exemple « $-10~\%$ d'émissions en 6 mois », « 2 jours végétariens par mois à la cantine », « zéro appareil laissé en veille » ;
\item des \textbf{responsables} désignés par poste (référent transport, référent énergie, référent cantine) ;
\item un \textbf{tableau de bord} : indicateurs simples (km covoiturés, kWh relevés, repas végétariens servis) mesurés chaque mois ;
\item une \textbf{présentation orale} des résultats devant la classe ou l'établissement, qui renforce l'engagement et essaime les bonnes pratiques.
\end{itemize}

\begin{popfigure}
\begin{center}
\begin{tikzpicture}
\begin{axis}[
width=0.85\linewidth, height=6cm,
xlabel={Mois}, ylabel={Émissions (kg \ce{CO2}éq)},
xmin=0.5, xmax=12.5, ymin=0, ymax=200,
xtick={1,2,3,4,5,6,7,8,9,10,11,12},
legend style={at={(0.02,0.98)}, anchor=north west, font=\small},
grid=major, grid style={popline!40},
]
\addplot[popcurve, popmuted, mark=*, dashed] coordinates {(1,167) (2,167) (3,167) (4,167) (5,167) (6,167) (7,167) (8,167) (9,167) (10,167) (11,167) (12,167)};
\addplot[popcurve, popgreen, mark=*] coordinates {(1,167) (2,165) (3,160) (4,155) (5,152) (6,150) (7,150) (8,148) (9,146) (10,145) (11,144) (12,143)};
\legend{Scénario sans action (2 t/an), Scénario avec charte climat}
\end{axis}
\end{tikzpicture}
\end{center}
\legende{Suivi mensuel simulé des émissions de la classe. En pointillés : le scénario « sans action » (constante à 167 kg \ce{CO2}éq/mois, soit 2 t/an). En vert : après mise en œuvre de la charte (covoiturage, extinction des veilleuses, menus végétariens), les émissions décroissent progressivement pour atteindre environ 143 kg/mois, soit une baisse d'environ 14~\% sur l'année.}
\end{popfigure}

\begin{exoresolu}[Objectif de la charte est-il atteint ?]
Énoncé. La classe s'est fixé « $-10~\%$ d'émissions en 6 mois ». Au mois 1, les émissions mensuelles sont de 167 kg \ce{CO2}éq ; au mois 6, elles sont de 150 kg \ce{CO2}éq. L'objectif est-il atteint ?
\tcblower
Solution. La réduction relative se calcule par :
\[ \text{Réduction} = \frac{167 - 150}{167} \times 100 = \frac{17}{167} \times 100 \approx 10{,}2~\%. \]
La réduction (environ 10,2~\%) dépasse légèrement l'objectif de 10~\% : \textbf{l'objectif est atteint}. La classe peut alors se fixer un nouveau palier, par exemple $-15~\%$ à 12 mois, en activant un levier supplémentaire (tri des déchets).
\end{exoresolu}

\begin{attention}[Pièges fréquents dans un bilan carbone scolaire]
\begin{itemize}
\item \textbf{Double comptage} : ne pas compter deux fois la même émission (ex. l'électricité de la cantine dans « énergie » \emph{et} dans « alimentation »).
\item \textbf{Oublier les unités} : un facteur d'émission en kg \ce{CO2}éq/km multiplié par des mètres donne un résultat 1000 fois faux.
\item \textbf{Confondre réduction absolue et relative} : « $-10$ kg » et « $-10~\%$ » ne disent pas la même chose.
\end{itemize}
\end{attention}

\begin{savaistu}[Le budget carbone mondial]
Les climatologues du GIEC estiment que le \textbf{budget carbone} restant — la quantité totale de \ce{CO2} que l'humanité peut encore émettre pour avoir une chance de limiter le réchauffement à $+1,5$~°C — est d'environ \textbf{400 Gt \ce{CO2}} à partir de 2020. Au rythme actuel des émissions mondiales (environ 40 Gt \ce{CO2}/an), ce budget serait épuisé en une dizaine d'années. Chaque tonne économisée, même à l'échelle d'une classe camerounaise, préserve une fraction de ce budget commun à toute l'humanité.
\end{savaistu}

\begin{aretenir}[Ce qu'il faut retenir de la section]
Mesurer (empreinte carbone = somme des activités $\times$ facteurs d'émission), hiérarchiser (agir d'abord sur le poste le plus émetteur), s'engager (charte chiffrée et datée), suivre (tableau de bord mensuel) et communiquer (présentation orale) : telle est la démarche complète qui transforme un citoyen informé en acteur du climat.
\end{aretenir}

\newpage

\coursec{Activité d'intégration}

\coursec{Lutte contre l'effet de serre}

\courssub{L'effet de serre naturel : principe et gaz impliqués}

\begin{definition}[Effet de serre naturel]
Phénomène physique naturel par lequel certains gaz de l'atmosphère, dits \cle{gaz à effet de serre (GES)}, absorbent une partie du rayonnement infrarouge émis par la surface terrestre et le réémettent dans toutes les directions, y compris vers le sol. Ce piégeage radiatif maintient la température moyenne du globe à environ $+15~^\circ\mathrm{C}$, contre $-18~^\circ\mathrm{C}$ en l'absence d'atmosphère.
\end{definition}

\begin{definition}[Gaz à effet de serre (GES) principaux]
Les GES identifiés par le GIEC et le Protocole de Kyoto sont : la vapeur d'eau ($\ce{H2O}$), le dioxyde de carbone ($\ce{CO2}$), le méthane ($\ce{CH4}$), le protoxyde d'azote ($\ce{N2O}$), les hydrofluorocarbures (HFC), les perfluorocarbures (PFC) et l'hexafluorure de soufre ($\ce{SF6}$). La vapeur d'eau est le plus abondant mais sa concentration est pilotée par la température (rétroaction) ; le $\ce{CO2}$ est le principal forçage anthropique de long terme.
\end{definition}

\begin{popfigure}
\begin{tikzpicture}[scale=0.85, every node/.style={font=\small}, >=Stealth]
% Soleil
\node[draw=poporange, fill=poporange!20, circle, minimum size=1.8cm, align=center] (sun) at (0,5){\textbf{Soleil}\\\SI{6000}{\kelvin}};
% Flèches rayonnement solaire
\draw[->, thick, poporange, decorate, decoration={snake, amplitude=1pt, segment length=4pt}] (sun) -- ++(0,-1.5) node[midway, right=2mm, align=center]{\textbf{Rayonnement solaire}\\(ondes courtes, UV/Visible)};
% Atmosphère
\draw[fill=popblue!10, draw=popblue, thick, rounded corners=5pt] (-4,3) rectangle (4,1.5) 
    node[midway, align=center, text width=5cm] {\textbf{Atmosphère}\\\cle{Gaz à effet de serre} : $\ce{CO2}$, $\ce{CH4}$, $\ce{H2O}$, $\ce{N2O}$...\\(transparent aux ondes courtes, absorbant aux ondes longues)};
% Sol
\draw[fill=popgreen!30, draw=popgreen, thick, rounded corners=5pt] (-4,1) rectangle (4,0) 
    node[midway] {\textbf{Surface terrestre (sol, océans, végétation)}};
% Rayonnement infrarouge émis par le sol
\draw[->, thick, popred] (-1.5,0) -- (-1.5,1.5) node[midway, left=2mm, align=center]{\textbf{Rayonnement infrarouge}\\(ondes longues, chaleur)};
\draw[->, thick, popred] (1.5,0) -- (1.5,1.5);
% Absorption / Réémission par les GES
\draw[->, thick, poppurple] (-1.5,2.2) -- (-1.5,2.7);
\draw[->, thick, poppurple] (1.5,2.2) -- (1.5,2.7);
\draw[->, thick, poppurple] (-1.5,2.2) -- (-3,2.2);
\draw[->, thick, poppurple] (1.5,2.2) -- (3,2.2);
\draw[->, thick, poppurple] (-1.5,2.2) -- (-1.5,1);
\draw[->, thick, poppurple] (1.5,2.2) -- (1.5,1);
% Échappement vers l'espace
\draw[->, thick, popblue] (-1.5,3) -- (-1.5,3.5) node[above] {Vers l'espace};
\draw[->, thick, popblue] (1.5,3) -- (1.5,3.5);
% Température
\node[draw=popdark, fill=popcream, rounded corners, align=center, below=0.5cm of sun] 
    {Équilibre radiatif : \textbf{Température moyenne $\approx +15~^\circ\mathrm{C}$}};
\end{tikzpicture}
\legende{Mécanisme de l'effet de serre naturel. 1. Le rayonnement solaire (ondes courtes) traverse l'atmosphère et chauffe le sol. 2. Le sol émet un rayonnement infrarouge (ondes longues). 3. Les GES absorbent ce rayonnement et le réémettent en partie vers le sol (flèches violettes), réchauffant l'atmosphère basse. 4. Une partie s'échappe vers l'espace (flèches bleues).}
\end{popfigure}

\begin{propriete}[Bilan radiatif terrestre]
L'équilibre thermique de la Terre impose : \textbf{Rayonnement solaire absorbé = Rayonnement infrarouge émis vers l'espace}. L'effet de serre naturel élève la température d'émission effective ($-18~^\circ\mathrm{C}$) à la température de surface ($+15~^\circ\mathrm{C}$). Tout accroissement de la concentration en GES perturbe cet équilibre (forçage radiatif positif) et impose un réchauffement jusqu'à un nouvel équilibre.
\end{propriete}

\courssub{Origines anthropiques du renforcement de l'effet de serre}

\begin{exemplebox}[Principales sources anthropiques d'émissions de GES (contexte mondial et camerounais)]
\begin{itemize}
\item \textbf{Combustion d'énergies fossiles} (charbon, pétrole, gaz) : production d'électricité (centrales thermiques), transports (essence, gasoil, kérosène), industrie. Au Cameroun : parc automobile vétuste, groupes électrogènes, centrale thermique de Kribi, Dibamba.
\item \textbf{Déforestation et changement d'affectation des sols} : coupe du bois d'œuvre, agriculture sur brûlis, extension des cultures (cacao, huile de palme) dans le bassin du Congo. Perte de puits de carbone et émissions par combustion/ décomposition.
\item \textbf{Agriculture et élevage} : fermentation entérique des ruminants ($\ce{CH4}$), gestion des déjections ($\ce{CH4}$, $\ce{N2O}$), engrais azotés ($\ce{N2O}$), riziculture inondée ($\ce{CH4}$). Au Cameroun : élevage bovin au Nord, cultures vivrières.
\item \textbf{Déchets} : décharges à ciel ouvert (méthanisation anaérobie), incinération à l'air libre, eaux usées non traitées. Problématique majeure à Yaoundé (Nkolfoulou), Douala (PK10).
\item \textbf{Procédés industriels} : cimenteries (décarbonatation du calcaire $\ce{CaCO3 -> CaO + CO2}$), aluminium, chimie.
\end{itemize}
\end{exemplebox}

\begin{aretenir}[Facteurs d'émission et CO$_2$ équivalent]
Pour comparer l'impact de différents GES, on utilise le \cle{Potentiel de Réchauffement Global (PRG)} sur 100 ans (valeurs GIEC AR6) : $\ce{CO2} = 1$ ; $\ce{CH4} \approx 27\text{--}30$ ; $\ce{N2O} \approx 273$. Les émissions sont ramenées en \textbf{kg $\ce{CO2}$ équivalent (kg $\ce{CO2}$éq)}. Un \cle{facteur d'émission} (FE) relie une donnée d'activité (km, kWh, kg de matière) à une émission en kg $\ce{CO2}$éq. Exemple : FE essence $\approx 2,3$ kg \ce{CO2}/L ou $0,21$ kg \ce{CO2}eq/km (valeur cycle de vie simplifiée).
\end{aretenir}

\courssub{Conséquences du réchauffement climatique}

\begin{propriete}[Conséquences observées et projetées (GIEC AR6)]
\begin{itemize}
\item \textbf{Thermiques :} Augmentation de la température moyenne mondiale ($+1,1~^\circ\mathrm{C}$ depuis 1850--1900), canicules plus fréquentes/intenses.
\item \textbf{Hydriques :} Fonte des glaces (Arctique, Antarctique, glaciers de montagne -- Mont Cameroun), élévation du niveau de la mer (menace sur le littoral camerounais : Limbé, Kribi, Douala), modification des régimes de pluies (sécheresses au Nord, inondations au Sud).
\item \textbf{Écosystèmes :} Déplacement des aires de répartition, blanchiment des coraux, pertes de biodiversité (forêts du bassin du Congo, mangroves).
\item \textbf{Socio-économiques :} Stress hydrique, baisse des rendements agricoles (maïs, mil, sorgho au Sahel), risques sanitaires (paludisme, choléra, maladies cardiovasculaires liées à la chaleur), migrations climatiques, coûts d'adaptation.
\end{itemize}
\end{propriete}

\begin{savaistu}[Le Cameroun face au changement climatique]
Le Cameroun, bien que faible émetteur (\textasciitilde 0,1 \% des émissions mondiales), est très vulnérable. Le \textbf{Plan National d'Adaptation au Changement Climatique (PNACC)} et la \textbf{Contribution Déterminée au Niveau National (CDN)} visent une réduction de 35 \% des émissions d'ici 2030 (conditionnelle). Projets concrets : reboisement \emph{Green Sahel} au Nord, gestion durable des forêts du bassin du Congo (REDD+), promotion des énergies renouvelables (solaire au Nord, hydroélectricité : barrages de Song Loulou, Édéa, Lagdo, Nachtigal, Memve'ele).
\end{savaistu}

\courssub{Moyens d'action : atténuation et adaptation}

\begin{definition}[Atténuation (Mitigation)]
Ensemble des actions visant à \textbf{réduire les émissions de GES} ou à \textbf{augmenter les puits de carbone} (reforestation, séquestration géologique). Exemples : efficacité énergétique, énergies renouvelables, mobilité douce, agriculture bas-carbone, économie circulaire.
\end{definition}

\begin{definition}[Adaptation]
Ensemble des actions visant à \textbf{réduire la vulnérabilité} des sociétés et écosystèmes aux impacts actuels et futurs du changement climatique. Exemples : variétés culturales résistantes à la sécheresse, digues de protection côtière, systèmes d'alerte précoce, urbanisation résiliente.
\end{definition}

\begin{methode}[Construire un plan d'action « Bas Carbone » à l'échelle d'un établissement]
\begin{enumerate}
\item \textbf{Diagnostic (Bilan carbone) :} Identifier les postes d'émission (Scopes 1, 2, 3 simplifiés : énergie, déplacements, achats, déchets). Collecter des données d'activité (consommations, km, tonnages).
\item \textbf{Quantification :} Appliquer les facteurs d'émission (Base Carbone\textsuperscript{\textregistered} ADEME, GIEC, données locales ENEO) : $\text{Émission} = \text{Activité} \times \text{FE}$.
\item \textbf{Hiérarchisation :} Classer les postes par contribution décroissante. Cibler les « gros postes » (règle de Pareto : 20 \% des actions pour 80 \% du gain).
\item \textbf{Scénario de réduction :} Pour chaque poste majeur, proposer des actions techniques (efficacité), organisationnelles (sobriété) ou de substitution (renouvelable). Estimer le gain (kg $\ce{CO2}$éq/an) et le coût.
\item \textbf{Suivi-évaluation :} Définir des indicateurs, un calendrier, des responsables. Communiquer (affichage, rapport).
\end{enumerate}
\end{methode}

\begin{attention}[Pièges fréquents dans le calcul d'empreinte carbone]
\begin{itemize}
\item Confondre \textbf{consommation finale} (ex: kWh facturés) et \textbf{consommation primaire} (intégrant les pertes de transport/transformation) -- utiliser le FE adapté au scope.
\item Oublier le \textbf{facteur temps} (jours ouvrés, semaines, heures de veille) ou le \textbf{nombre d'unités} (élèves, véhicules).
\item Utiliser un FE \textbf{inadapté au contexte local} (ex: FE électricité France $\approx 0,057$ kg \ce{CO2}/kWh vs Cameroun $\approx 0,05$ kg \ce{CO2}/kWh grâce à l'hydroélectricité).
\item Négliger le \textbf{Scope 3} (achats, déplacements domicile-travail, déchets) qui représente souvent $> 50 \%$ du bilan d'un service public.
\item Croire que \textbf{compenser (planter des arbres)} dispense de \textbf{réduire} : la séquestration est lente, incertaine (incendies, maladies) et les terres sont limitées. Priorité : \textbf{Éviter > Réduire > Compenser}.
\end{itemize}
\end{attention}

\courssub{Activité d'intégration : « Le lycée de Nkolbisson face au défi climatique »}

\courssub{Objectifs de l'activité}

À l'issue de cette activité d'intégration, tu seras capable de :

\begin{itemize}
\item mobiliser les savoirs de la leçon (mécanisme de l'effet de serre, gaz à effet de serre, origines anthropiques des émissions, conséquences du réchauffement) dans une situation concrète et complexe ;
\item mettre en œuvre les savoir-faire du programme : relier une activité humaine à une émission de gaz à effet de serre, calculer une émission à partir de données réelles et de facteurs d'émission, proposer et hiérarchiser des mesures de réduction ;
\item travailler en équipe, collecter des données, les organiser dans un tableau, effectuer des calculs simples et présenter des résultats sous forme d'affichage ;
\item développer un savoir-être : sens de la responsabilité collective, esprit critique face aux chiffres, engagement citoyen pour un développement durable de ton établissement.
\end{itemize}

\courssub{Situation}

\textbf{Le lycée de Nkolbisson face au défi climatique.}

Au lycée de Nkolbisson (Yaoundé), le club environnement a suivi la leçon sur la lutte contre l'effet de serre. Interpellés par les informations sur le renforcement anthropique de l'effet de serre, les élèves de la classe de Première C décident de répondre à la question suivante, posée par leur professeur de SVTEEHB :

\begin{quote}
\emph{« Quelle est l'empreinte carbone hebdomadaire de notre classe, et quelles actions concrètes pourraient la réduire de manière significative ? »}
\end{quote}

La classe compte \textbf{40 élèves}, répartis en 10 groupes de 4. Pendant une semaine de cours (5 jours), chaque groupe collecte des données sur quatre postes d'émission. Voici les résultats moyens obtenus pour l'ensemble de la classe sur la semaine :

\textbf{Poste 1 : déplacements domicile--lycée (aller-retour quotidien).}
\begin{itemize}
\item 15 élèves viennent en voiture (ou moto-taxi essence), distance moyenne aller-retour : \SI{8}{km} par jour ;
\item 18 élèves prennent le bus ou le taxi collectif, distance moyenne aller-retour : \SI{10}{km} par jour ;
\item 7 élèves viennent à pied ou à vélo (émission considérée comme nulle).
\end{itemize}

\textbf{Poste 2 : consommation électrique de la salle de classe.}
Le relevé du compteur de la salle (éclairage, ventilateurs, veilleuses des appareils) donne une consommation de \SI{60}{kWh} sur la semaine, dont on estime que \SI{12}{kWh} correspondent aux veilleuses et appareils laissés allumés inutilement la nuit et le week-end.

\textbf{Poste 3 : alimentation à la cantine.}
La cantine sert en moyenne, par semaine et pour l'ensemble des élèves de la classe qui y mangent : \SI{8}{kg} de bœuf, \SI{6}{kg} de poulet et \SI{10}{kg} de plats végétariens (haricots, folong, riz, manioc).

\textbf{Poste 4 : déchets.}
La classe produit \SI{25}{kg} de déchets par semaine, dont \SI{5}{kg} de plastiques et papiers effectivement triés et recyclés ; le reste (\SI{20}{kg}) part en décharge sauvage sans tri.

\textbf{Facteurs d'émission fournis par le professeur} (valeurs simplifiées, exprimées en kilogrammes de $\ce{CO2}$ équivalent, noté kg $\ce{CO2}$éq) :

\begin{center}
\begin{tabularx}{\linewidth}{|l|X|}\hline
\rowcolor{popblueL} \textbf{Source d'émission} & \textbf{Facteur d'émission} \\ \hline
Voiture / moto essence & 0,21 kg $\ce{CO2}$éq par km parcouru \\ \hline
Bus / taxi collectif & 0,09 kg $\ce{CO2}$éq par km parcouru \\ \hline
Électricité (réseau ENEO, mix hydroélectrique + thermique) & 0,05 kg $\ce{CO2}$éq par kWh \\ \hline
Bœuf & 2,5 kg $\ce{CO2}$éq par kg consommé \\ \hline
Poulet & 0,8 kg $\ce{CO2}$éq par kg consommé \\ \hline
Plat végétarien & 0,3 kg $\ce{CO2}$éq par kg consommé \\ \hline
Déchets non recyclés & 0,5 kg $\ce{CO2}$éq par kg \\ \hline
\end{tabularx}
\end{center}

\textbf{Problème à résoudre :} calculer le bilan carbone hebdomadaire de la classe, identifier les trois postes les plus émetteurs, puis construire un plan d'action « École bas carbone » chiffré, réaliste et affichable dans le couloir du lycée.

\courssub{Consignes}

\begin{enumerate}
\item Rappelle, en quelques phrases et à l'aide d'un schéma simple, le mécanisme de l'effet de serre naturel et explique pourquoi les émissions de $\ce{CO2}$ d'origine humaine le renforcent.
\item Pour chacun des quatre postes (déplacements, électricité, alimentation, déchets), calcule l'émission hebdomadaire de la classe en kg $\ce{CO2}$éq, en détaillant les calculs. Présente les résultats dans un tableau récapitulatif.
\item Calcule le bilan carbone total hebdomadaire de la classe, puis l'émission moyenne par élève et par semaine.
\item Classe les postes par ordre décroissant d'émission et identifie les trois postes majeurs. Ce classement correspond-il à ton intuition de départ ? Commente.
\item Propose \textbf{trois actions concrètes} de réduction, réalisables par la classe (une par poste majeur si possible). Pour chaque action, estime la réduction d'émission obtenue en kg $\ce{CO2}$éq par semaine, puis calcule la réduction totale en pourcentage du bilan initial.
\item Rédige l'affiche synthétique « École bas carbone » destinée au couloir : titre, chiffres clés, trois actions, slogan. Elle doit être compréhensible par un élève de Sixième.
\item (Pour aller plus loin) Extrapole le bilan de la classe à l'ensemble du lycée (30 classes) sur une année scolaire de 36 semaines. Compare avec l'émission absorbée par un arbre adulte, estimée à environ \SI{25}{kg} de $\ce{CO2}$ par an : combien d'arbres faudrait-il planter pour compenser ?
\end{enumerate}

\courssub{Ressources à mobiliser}

\begin{definition}[Bilan carbone / Empreinte carbone]
Quantité totale de gaz à effet de serre émise directement ou indirectement par une activité, une organisation, un produit ou un individu, exprimée en équivalent $\ce{CO2}$ (kg ou t $\ce{CO2}$éq). Elle couvre les scopes 1 (directes), 2 (énergie indirecte) et 3 (autres indirectes : déplacements, achats, déchets).
\end{definition}

\begin{methode}[Calculer une émission de GES à partir d'un facteur d'émission]
\begin{enumerate}
\item Identifier la \textbf{donnée d'activité} (quantité : km parcourus, kWh consommés, kg de matière, etc.).
\item Choisir le \textbf{facteur d'émission (FE)} adapté (unité : kg $\ce{CO2}$éq / unité d'activité), pertinent pour le contexte géographique et technologique.
\item Appliquer la formule : $\displaystyle \text{Émission (kg $\ce{CO2}$éq)} = \text{Quantité d'activité} \times \text{FE}$.
\item Vérifier l'homogénéité des unités et l'ordre de grandeur.
\end{enumerate}
\end{methode}

\begin{aretenir}[Boîte à outils de l'activité]
\begin{itemize}
\item \textbf{Mécanisme de l'effet de serre :} le rayonnement solaire (courtes longueurs d'onde) traverse l'atmosphère et chauffe le sol ; la Terre réémet un rayonnement infrarouge (grandes longueurs d'onde) que les gaz à effet de serre ($\ce{CO2}$, méthane $\ce{CH4}$, vapeur d'eau, protoxyde d'azote $\ce{N2O}$) absorbent en partie, ce qui maintient une température moyenne d'environ $+15~^\circ\mathrm{C}$ au lieu de $-18~^\circ\mathrm{C}$. L'augmentation de la concentration de ces gaz, liée aux activités humaines, renforce ce piégeage et provoque le réchauffement climatique.
\item \textbf{Formule du calcul d'émission :}
\[ \text{Émission (kg $\ce{CO2}$éq)} = \text{quantité d'activité} \times \text{facteur d'émission} \]
\item \textbf{Pourcentage de réduction :}
\[ \text{Réduction (\%)} = \frac{\text{émission évitée}}{\text{émission initiale}} \times 100 \]
\item \textbf{$\ce{CO2}$ équivalent :} unité qui ramène tous les gaz à effet de serre à l'équivalent $\ce{CO2}$ (par exemple, 1 kg de méthane équivaut, sur 100 ans, à environ 27--30 kg de $\ce{CO2}$ selon le dernier rapport GIEC).
\end{itemize}
\end{aretenir}

\begin{attention}[Points de vigilance pour l'activité]
\begin{itemize}
\item \textbf{Déplacements :} Ne pas oublier de multiplier par le nombre de jours (5) et par le nombre d'élèves concernés. Distinguer les modes de transport (FE différents).
\item \textbf{Électricité :} Le FE donné ($0,05$ kg $\ce{CO2}$éq/kWh) s'applique à la consommation \textbf{totale} relevée au compteur (\SI{60}{kWh}), pas seulement aux veilleuses. Les \SI{12}{kWh} de veilleuses servent pour l'action de réduction.
\item \textbf{Alimentation :} Les FE sont donnés par kg de produit \textbf{consommé} (tel que servi). Bien additionner les trois catégories.
\item \textbf{Déchets :} Seuls les déchets \textbf{non recyclés} (\SI{20}{kg}) émettent selon le FE donné. Les \SI{5}{kg} recyclés sont considérés comme évités (ou comptabilisés dans le cycle de vie amont, mais ici FE = 0 pour la fin de vie).
\item \textbf{Unités :} Tous les résultats en \textbf{kg $\ce{CO2}$éq / semaine}. Vérifier les virgules décimales.
\end{itemize}
\end{attention}

\courssub{Démarche et corrigé commenté}

\begin{exoresolu}[Bilan carbone de la classe de Première C -- Corrigé détaillé]
\textbf{Énoncé.} À partir des données de la situation, calculer le bilan carbone hebdomadaire de la classe, identifier les postes majeurs et chiffrer un plan d'action de réduction.

\tcblower

\textbf{Question 1 -- Rappel du mécanisme.} Le Soleil envoie vers la Terre un rayonnement riche en lumière visible (ondes courtes), qui traverse largement l'atmosphère et chauffe le sol. Le sol, ainsi réchauffé, émet à son tour un rayonnement infrarouge (ondes longues) vers l'espace. Les gaz à effet de serre ($\ce{CO2}$, $\ce{CH4}$, vapeur d'eau, $\ce{N2O}$) absorbent une partie de cet infrarouge et le réémettent dans toutes les directions, y compris vers le sol : c'est l'\cle{effet de serre}, phénomène naturel et indispensable à la vie. En brûlant des combustibles fossiles (essence, gasoil, charbon) et par certaines pratiques (élevage intensif, déforestation, décharges anaérobies), l'être humain rejette des quantités croissantes de $\ce{CO2}$ et de méthane : la « couverture » atmosphérique s'épaissit, le piégeage de la chaleur s'intensifie (forçage radiatif positif), et la température moyenne du globe augmente.

\begin{popfigure}
\begin{tikzpicture}[scale=0.7, every node/.style={font=\small}, >=Stealth]
% Sun
\node[draw=poporange, fill=poporange!20, circle, minimum size=1.5cm] (sun) at (0,4.5) {Soleil};
% Solar radiation
\draw[->, thick, poporange, decorate, decoration={snake, amplitude=1pt, segment length=3pt}] (sun) -- ++(0,-1.2) node[midway, right=2mm] {Ondes courtes};
% Atmosphere with enhanced GHG
\draw[fill=popblue!15, draw=popblue, thick, rounded corners=5pt] (-3.5,3) rectangle (3.5,1.8) 
    node[midway, align=center, text width=5cm] {\textbf{Atmosphère enrichie en GES}\\($\uparrow \ce{CO2}$, $\uparrow \ce{CH4}$)\\\textit{Absorption IR accrue}};
% Ground
\draw[fill=popgreen!30, draw=popgreen, thick, rounded corners=5pt] (-3.5,1.3) rectangle (3.5,0.3) 
    node[midway] {Surface terrestre};
% IR emission
\draw[->, thick, popred] (-1,0.3) -- (-1,1.8) node[midway, left=2mm] {Infrarouge};
\draw[->, thick, popred] (1,0.3) -- (1,1.8);
% Trapping
\draw[->, thick, poppurple] (-1,2.4) -- (-1,2.9);
\draw[->, thick, poppurple] (1,2.4) -- (1,2.9);
\draw[->, thick, poppurple] (-1,2.4) -- (-2.5,2.4);
\draw[->, thick, poppurple] (1,2.4) -- (2.5,2.4);
\draw[->, thick, poppurple] (-1,2.4) -- (-1,1.3);
\draw[->, thick, poppurple] (1,2.4) -- (1,1.3);
% Less escape
\draw[->, thick, popblue, dashed] (-1,3) -- (-1,3.3);
\draw[->, thick, popblue, dashed] (1,3) -- (1,3.3) node[right] {Moins vers l'espace};
% Warming label
\node[draw=popdark, fill=poppink!20, rounded corners, align=center, below=0.4cm of sun] 
    {\textbf{Renforcement de l'effet de serre} $\rightarrow$ \textbf{Réchauffement climatique}};
\end{tikzpicture}
\legende{Schématisation du renforcement anthropique : l'augmentation de la concentration en GES (zone bleue plus dense) augmente l'absorption/réémission du rayonnement infrarouge vers le sol (flèches violettes plus nombreuses), réduisant le flux sortant vers l'espace (flèches bleues pointillées).}
\end{popfigure}

\textbf{Question 2 -- Calcul des émissions par poste.}

\emph{Poste 1 : déplacements.}
\begin{align*}
\text{Voitures / motos : } & 15 \text{ élèves} \times 8 \text{ km/j} \times 5 \text{ j} \times 0{,}21 \text{ kg $\ce{CO2}$éq/km} = 126 \text{ kg $\ce{CO2}$éq} \\
\text{Bus / taxis collectifs : } & 18 \text{ élèves} \times 10 \text{ km/j} \times 5 \text{ j} \times 0{,}09 \text{ kg $\ce{CO2}$éq/km} = 81 \text{ kg $\ce{CO2}$éq} \\
\text{À pied / vélo : } & 0 \text{ kg $\ce{CO2}$éq}
\end{align*}
Total déplacements : $126 + 81 = \mathbf{207}$ kg $\ce{CO2}$éq.

\emph{Poste 2 : électricité.}
\[ 60 \text{ kWh} \times 0{,}05 \text{ kg $\ce{CO2}$éq/kWh} = \mathbf{3} \text{ kg $\ce{CO2}$éq} \]

\emph{Poste 3 : alimentation.}
\begin{align*}
\text{Bœuf : } & 8 \text{ kg} \times 2{,}5 \text{ kg $\ce{CO2}$éq/kg} = 20 \text{ kg $\ce{CO2}$éq} \\
\text{Poulet : } & 6 \text{ kg} \times 0{,}8 \text{ kg $\ce{CO2}$éq/kg} = 4{,}8 \text{ kg $\ce{CO2}$éq} \\
\text{Végétarien : } & 10 \text{ kg} \times 0{,}3 \text{ kg $\ce{CO2}$éq/kg} = 3 \text{ kg $\ce{CO2}$éq}
\end{align*}
Total alimentation : $20 + 4{,}8 + 3 = \mathbf{27{,}8}$ kg $\ce{CO2}$éq.

\emph{Poste 4 : déchets.}
\[ 20 \text{ kg non recyclés} \times 0{,}5 \text{ kg $\ce{CO2}$éq/kg} = \mathbf{10} \text{ kg $\ce{CO2}$éq} \]

\begin{center}
\begin{tabularx}{\linewidth}{|l|X|X|}\hline
\rowcolor{popblueL} \textbf{Poste} & \textbf{Émission hebdomadaire (kg $\ce{CO2}$éq)} & \textbf{Part du total} \\ \hline
Déplacements & 207 & 83,5 \% \\ \hline
Alimentation & 27,8 & 11,2 \% \\ \hline
Déchets & 10 & 4,0 \% \\ \hline
Électricité & 3 & 1,2 \% \\ \hline
\textbf{Total} & \textbf{247,8} & \textbf{100 \%} \\ \hline
\end{tabularx}
\end{center}

\textbf{Question 3 -- Bilan total et émission par élève.}
\[ \text{Total} = 207 + 3 + 27{,}8 + 10 = 247{,}8 \text{ kg $\ce{CO2}$éq/semaine} \]
\[ \text{Par élève : } \frac{247{,}8}{40} \approx 6{,}2 \text{ kg $\ce{CO2}$éq/semaine} \]

\textbf{Question 4 -- Classement et commentaire.} Les trois postes majeurs sont : (1) les \textbf{déplacements} (83,5 \%), (2) l'\textbf{alimentation} (11,2 \%), (3) les \textbf{déchets} (4,0 \%). L'électricité, souvent citée en premier dans les discussions grand public, est ici très minoritaire grâce à la forte part de l'hydroélectricité dans le mix camerounais (barrages de Song Loulou, Édéa, Lagdo, Nachtigal, Memve'ele) : un résultat qui invite à raisonner à partir des données locales et non des idées reçues importées.

\textbf{Question 5 -- Plan d'action « École bas carbone ».}

\emph{Action 1 (Poste Déplacements) : Journées « Mobilité douce » (2 jours/semaine).} Les 15 élèves motorisés prennent le bus/taxi collectif 2 jours sur 5 au lieu de la voiture/moto.
\[ \text{Gain} = 15 \times 8 \text{ km/j} \times 2 \text{ j} \times (0{,}21 - 0{,}09) \text{ kg $\ce{CO2}$éq/km} = 28{,}8 \text{ kg $\ce{CO2}$éq/semaine} \]

\emph{Action 2 (Poste Alimentation) : Un menu « Végétarien local » par semaine.} Remplacer la moitié de la ration de bœuf hebdomadaire (4 kg sur 8 kg) par des plats végétariens (haricots, folong, manioc).
\[ \text{Gain} = 4 \text{ kg} \times (2{,}5 - 0{,}3) \text{ kg $\ce{CO2}$éq/kg} = 8{,}8 \text{ kg $\ce{CO2}$éq/semaine} \]

\emph{Action 3 (Postes Électricité + Déchets) : « Zéro veille \& Tri total ».} Éteindre systématiquement les veilleuses (\SI{12}{kWh} économisés) et recycler la moitié des déchets résiduels (10 kg sur 20 kg).
\begin{align*}
\text{Gain électricité} &= 12 \text{ kWh} \times 0{,}05 = 0{,}6 \text{ kg $\ce{CO2}$éq} \\
\text{Gain déchets} &= 10 \text{ kg} \times 0{,}5 = 5 \text{ kg $\ce{CO2}$éq} \\
\text{Total Action 3} &= 5{,}6 \text{ kg $\ce{CO2}$éq/semaine}
\end{align*}

\emph{Réduction totale hebdomadaire :} $28{,}8 + 8{,}8 + 5{,}6 = 43{,}2$ kg $\ce{CO2}$éq/semaine.
\[ \text{Réduction (\%)} = \frac{43{,}2}{247{,}8} \times 100 \approx 17{,}4 \% \]
Le plan permet donc de réduire le bilan d'environ \textbf{17 \%} sans investissement matériel majeur, par des changements de comportement accessibles.

\textbf{Question 6 -- Affiche « École bas carbone » (contenu pour réalisation).}
\begin{center}
\begin{tabularx}{\linewidth}{|X|}\hline
\rowcolor{popgreenL} \textbf{\Large LYCÉE DE NKOLBISSON -- ÉCOLE BAS CARBONE} \\ \hline
\textbf{Notre empreinte :} $247,8$ kg $\ce{CO2}$éq/semaine (\SI{6,2}{kg/élève/semaine}) \\ 
\textbf{Top 3 des postes :} 
\begin{itemize}
\item \textbf{1. Déplacements} : \SI{83,5}{\%} (\SI{207}{kg})
\item \textbf{2. Alimentation} : \SI{11,2}{\%} (\SI{27,8}{kg})
\item \textbf{3. Déchets} : \SI{4,0}{\%} (\SI{10}{kg})
\end{itemize} \\
\textbf{Nos 3 actions pour -17 \% :} 
\begin{enumerate}
\item \textbf{Bus/Vélo/Marche 2j/sem} : $-$ \SI{28,8}{kg} ($\SI{11,6}{\%}$)
\item \textbf{1 Menu Végé local/sem} : $-$ \SI{8,8}{kg} ($\SI{3,6}{\%}$)
\item \textbf{Zéro veille + Tri total} : $-$ \SI{5,6}{kg} ($\SI{2,3}{\%}$)
\end{enumerate} \\
\textbf{Objectif :} \textbf{- 50 \% d'ici 2027 !} \\
\rowcolor{popgoldL} \textbf{\textit{« Moins de $\ce{CO2}$, plus d'avenir : notre lycée s'engage ! »}} \\ \hline
\end{tabularx}
\end{center}

\textbf{Question 7 -- Extrapolation à l'échelle du lycée et compensation.}
\begin{align*}
\text{Émission annuelle lycée (30 classes, 36 sem)} &= 247{,}8 \times 30 \times 36 \\
&= 267\,624 \text{ kg $\ce{CO2}$éq} \approx \mathbf{268 \text{ tonnes $\ce{CO2}$éq/an}} \\
\text{Nombre d'arbres pour compensation (25 kg $\ce{CO2}$/an/arbre)} &= \frac{267\,624}{25} \approx \mathbf{10\,705 \text{ arbres}}
\end{align*}
\textbf{Interprétation :} Planter \SI{10700}{arbres} représente une surface d'environ \SI{10}{ha} à \SI{20}{ha} (densité 500--1000 arbres/ha), soit un boisement significatif. Cela montre que la \textbf{compensation par plantation est un complément nécessaire mais insuffisant seule} : la priorité absolue reste la \textbf{réduction à la source} (atténuation), car la capacité d'absorption des écosystèmes est limitée et non permanente (risques incendies, maladies, coupe). L'adaptation (résilience du lycée : ombrage, eau, ventilation) doit aussi être préparée.

\end{exoresolu}

\courssub{Bilan de l'activité}

Cette activité d'intégration a permis de mobiliser l'ensemble des acquis de la leçon dans une démarche complète : collecte de données réelles, calculs d'émissions à l'aide de facteurs, interprétation critique des résultats et proposition d'un plan d'action chiffré. Trois enseignements majeurs en ressortent :

\begin{itemize}
\item \textbf{Hiérarchiser avant d'agir :} dans le contexte camerounais, ce sont les déplacements motorisés individuels qui dominent le bilan, loin devant l'électricité -- une conséquence directe du rôle central de l'hydroélectricité (barrages de Song Loulou, Édéa, Lagdo, Nachtigal, Memve'ele) dans la production nationale. Les actions efficaces ne sont pas toujours celles que l'intuition ou les messages médiatiques généraux désignent.
\item \textbf{Chiffrer pour convaincre :} un plan d'action crédible repose sur des estimations quantifiées (ici, environ 17 \% de réduction potentielle immédiate) ; la démarche scientifique transforme une bonne intention en projet argumenté, présentable à l'administration de l'établissement, aux parents d'élèves et aux partenaires (MINESEC, MINEPDED, communes).
\item \textbf{Relier atténuation et adaptation :} réduire les émissions (atténuation) reste la priorité absolue, car la compensation par la plantation d'arbres, bien que précieuse pour les sols, la biodiversité et le cadre de vie, ne suffit pas à elle seule à neutraliser nos émissions actuelles. L'école doit aussi s'adapter (végétalisation des cours, isolation, gestion de l'eau).
\end{itemize}

En réalisant l'affiche « École bas carbone », tu as enfin exercé un savoir-être essentiel : la responsabilité citoyenne. La lutte contre le renforcement de l'effet de serre commence par des gestes mesurés, partagés et répétés -- à l'échelle d'une classe, d'un lycée, puis d'un pays tout entier.

\begin{exercice}[Entraînement : Empreinte carbone d'un repas type]
Calculer l'empreinte carbone (en kg $\ce{CO2}$éq) de deux repas servis à la cantine :
\begin{enumerate}
\item \textbf{Repas A :} 150 g de bœuf, 200 g de riz, 100 g de légumes, 50 g de pain, 1 yaourt (lait).
\item \textbf{Repas B :} 100 g de poisson (maquereau), 200 g de manioc, 100 g de légumes (folong), 50 g de pain, 1 fruit (banane).
\end{enumerate}
Facteurs d'émission (simplifiés) : Bœuf 2,5 ; Poisson 0,6 ; Riz 0,4 ; Manioc 0,2 ; Légumes 0,3 ; Pain 0,8 ; Yaourt 1,2 ; Banane 0,1 (unités : kg $\ce{CO2}$éq/kg). Comparer et proposer une action de substitution.
\end{exercice}

\begin{corrige}[Exercice : Empreinte carbone d'un repas type]
\textbf{Repas A :}
\begin{align*}
\text{Bœuf} &: 0,150 \times 2,5 = 0,375 \\
\text{Riz} &: 0,200 \times 0,4 = 0,080 \\
\text{Légumes} &: 0,100 \times 0,3 = 0,030 \\
\text{Pain} &: 0,050 \times 0,8 = 0,040 \\
\text{Yaourt} &: 0,100 \times 1,2 = 0,120 \quad (\text{estimation 100 g}) \\
\text{Total A} &= \mathbf{0,645 \text{ kg $\ce{CO2}$éq}}
\end{align*}

\textbf{Repas B :}
\begin{align*}
\text{Poisson} &: 0,100 \times 0,6 = 0,060 \\
\text{Manioc} &: 0,200 \times 0,2 = 0,040 \\
\text{Légumes} &: 0,100 \times 0,3 = 0,030 \\
\text{Pain} &: 0,050 \times 0,8 = 0,040 \\
\text{Banane} &: 0,150 \times 0,1 = 0,015 \\
\text{Total B} &= \mathbf{0,185 \text{ kg $\ce{CO2}$éq}}
\end{align*}

\textbf{Comparaison :} Le repas B (protéine animale peu émettrice + féculent local bas carbone) émet \textbf{3,5 fois moins} de GES que le repas A (bœuf + riz importé). 
\textbf{Action :} Substituer une partie des repas « bœuf/riz » par des repas « poisson/manioc » ou « légumineuses/manioc » réduit drastiquement l'empreinte alimentation tout en valorisant les filières locales (pêche artisanale, tubercules).
\end{corrige}

\newpage

\coursec{Méthodes et exercices résolus}

\coursec{L'effet de serre naturel : principe et gaz impliqués}

\emph{Observation :} La température moyenne à la surface de la Terre est d'environ $15\ ^\circ\mathrm{C}$, alors que la Lune, à distance similaire du Soleil, voit sa surface varier entre $+120\ ^\circ\mathrm{C}$ (jour) et $-170\ ^\circ\mathrm{C}$ (nuit). Pourquoi la Terre bénéficie-t-elle d'un climat aussi tempéré et stable ? L'atmosphère joue le rôle d'une « couverture thermique » : c'est l'effet de serre naturel.

\begin{definition}[Effet de serre naturel]
L'\cle{effet de serre naturel} est le phénomène par lequel certains gaz de l'atmosphère, dits \cle{gaz à effet de serre (GES)}, absorbent une partie du rayonnement infrarouge émis par la surface terrestre chauffée par le Soleil, puis réémettent de la chaleur vers la surface, maintenant une température moyenne compatible avec l'eau liquide et la vie.
\end{definition}

\begin{propriete}[Les principaux gaz à effet de serre et leur origine naturelle]
\begin{center}
\begin{tabularx}{\linewidth}{|l|X|X|}\hline
\rowcolor{popblueL} \textbf{Gaz} & \textbf{Formule} & \textbf{Principales sources naturelles} \\ \hline
Dioxyde de carbone & \ce{CO2} & Respiration des êtres vivants, dégazage volcanique, décomposition organique, échanges océan-atmosphère \\ \hline
Méthane & \ce{CH4} & Zones humides (marais, tourbières), termites, océans, hydrates de méthane \\ \hline
Protoxyde d'azote & \ce{N2O} & Sols (nitrification/dénitrification), océans \\ \hline
Vapeur d'eau & \ce{H2O} & Évaporation des océans, lacs, évapotranspiration des végétaux (principal GES en quantité, mais durée de vie courte) \\ \hline
\end{tabularx}
\end{center}
\emph{Note :} La vapeur d'eau est le premier GES en contribution à l'effet de serre naturel ($\approx 60\ \%$), mais sa concentration est pilotée par la température (rétroaction positive) ; les gaz à longue durée de vie (\ce{CO2}, \ce{CH4}, \ce{N2O}) sont les \cle{forçages} qui initient le changement climatique.
\end{propriete}

\begin{popfigure}
\begin{tikzpicture}[scale=0.9, every node/.style={font=\small}]
% Sun
\node[circle, fill=popgold, minimum size=1.2cm] (sun) at (0,5) {};
\node[above] at (sun.north) {\textbf{Soleil}};
% Solar rays
\draw[poporange, thick, ->, decorate, decoration={snake, amplitude=1pt, segment length=4pt}] (-1.5,4.5) -- (-2,2.5) node[midway, left, align=center] {Rayonnement solaire \\(ondes courtes)};
\draw[poporange, thick, ->, decorate, decoration={snake, amplitude=1pt, segment length=4pt}] (1.5,4.5) -- (2,2.5);
% Atmosphere top
\draw[popblue, thick] (-3,3) -- (3,3) node[midway, above] {Sommet de l'atmosphère};
% Clouds / Reflection
\draw[popmuted, thick, ->] (-1,3) -- (-1,4) node[above] {$30\ \%$ réfléchis};
\draw[popmuted, thick, ->] (1,3) -- (1,4);
% Atmosphere layer
\draw[popblue!30, fill=popblue!5, thick] (-3,3) -- (3,3) -- (3,0.5) -- (-3,0.5) -- cycle;
\node[popblue, align=center] at (0,2.2) {Atmosphère \\ \textbf{Gaz à effet de serre :} \ce{CO2}, \ce{CH4}, \ce{N2O}, \ce{H2O}};
% Surface
\draw[poporange!80!black, fill=poporange!20, thick] (-3,0.5) -- (3,0.5) -- (3,-0.5) -- (-3,-0.5) -- cycle;
\node[poporange!80!black, align=center] at (0,0) {Surface terrestre \\ (sols, océans, végétation)};
% Absorption by surface
\draw[poporange, thick, ->] (-0.5,3) -- (-0.5,0.5) node[midway, right] {$70\ \%$ absorbés};
% Infrared emission from surface
\draw[poppink, thick, ->, decorate, decoration={snake, amplitude=1.5pt, segment length=5pt}] (-1,0.5) -- (-1,-1) node[below, align=center] {Rayonnement infrarouge \\(ondes longues)};
\draw[poppink, thick, ->, decorate, decoration={snake, amplitude=1.5pt, segment length=5pt}] (1,0.5) -- (1,-1);
% Trapping by GHG
\draw[poppink, thick, ->] (-1.5,1.5) -- (-1.5,0.5) node[midway, left] {Absorption par GES};
\draw[poppink, thick, ->] (1.5,1.5) -- (1.5,0.5);
% Re-emission down
\draw[poppink, thick, ->, dashed] (-1.5,1.5) .. controls (-2,1) and (-2,0) .. (-1.5,0.5) node[midway, left] {Réémission vers le sol};
\draw[poppink, thick, ->, dashed] (1.5,1.5) .. controls (2,1) and (2,0) .. (1.5,0.5);
% Escape to space
\draw[poppink, thick, ->] (-2.5,2) -- (-2.5,3.5) node[above] {Partie échappée vers l'espace};
\draw[poppink, thick, ->] (2.5,2) -- (2.5,3.5);
% Temperature
\node[draw, fill=popcream, rounded corners, align=center] at (4.5,0) {Température moyenne \\ $\approx 15\ ^\circ\mathrm{C}$ \\ (au lieu de $-18\ ^\circ\mathrm{C}$)};
\end{tikzpicture}
\legende{Schéma du mécanisme de l'effet de serre naturel. Le rayonnement solaire (ondes courtes, flèches orange) traverse l'atmosphère et chauffe la surface. La surface réémet un rayonnement infrarouge (ondes longues, flèches roses) partiellement piégé par les gaz à effet de serre, qui renvoient de la chaleur vers le sol.}
\end{popfigure}

\courssub{Méthode 1 : Expliquer le mécanisme de l'effet de serre à partir d'un schéma}

\begin{methode}[Décrire et expliquer un schéma de l'effet de serre]
\begin{enumerate}
\item \textbf{Identifier les trois zones du schéma} : le Soleil (source d'énergie), l'atmosphère (couche de gaz entourant la Terre) et la surface terrestre (sols, océans).
\item \textbf{Suivre le rayonnement solaire} : une partie est réfléchie vers l'espace (par les nuages, les glaces, les surfaces claires = \cle{albédo}), une partie traverse l'atmosphère et est \cle{absorbée} par la surface qui s'échauffe.
\item \textbf{Suivre le rayonnement infrarouge} : la Terre chauffée réémet de la chaleur sous forme de \cle{rayonnement infrarouge} vers l'espace.
\item \textbf{Localiser le piégeage} : une partie de ce rayonnement infrarouge est \cle{absorbée par les gaz à effet de serre} (GES : \ce{CO2}, méthane \ce{CH4}, vapeur d'eau, protoxyde d'azote \ce{N2O}), qui renvoient de la chaleur vers la surface : c'est l'\cle{effet de serre}.
\item \textbf{Conclure sur le rôle} : l'effet de serre \emph{naturel} maintient la température moyenne de la Terre autour de $15\ ^\circ\mathrm{C}$ au lieu d'environ $-18\ ^\circ\mathrm{C}$ ; il est donc indispensable à la vie. Le problème vient de son \emph{renforcement} par les activités humaines.
\end{enumerate}
\textbf{Réflexes de rédaction} : employer les verbes précis (absorber, réfléchir, réémettre, piéger) ; nommer les GES par leurs formules chimiques ; distinguer toujours « effet de serre naturel » (bénéfique) et « renforcement de l'effet de serre » (problématique).
\end{methode}

\begin{exoresolu}[Analyse d'un schéma de l'effet de serre]
Sur un schéma de l'effet de serre, on lit : sur $100$ unités d'énergie solaire arrivant au sommet de l'atmosphère, $30$ sont réfléchies vers l'espace et $70$ absorbées par la surface terrestre ; la surface réémet un rayonnement infrarouge dont une grande partie est absorbée par les gaz à effet de serre puis renvoyée vers le sol.
\begin{enumerate}
\item Définir l'effet de serre.
\item Citer deux gaz à effet de serre et donner leur formule chimique.
\item Expliquer pourquoi l'effet de serre naturel est indispensable à la vie.
\end{enumerate}
\tcblower
\textbf{1. Définition.} L'\cle{effet de serre} est le phénomène par lequel certains gaz de l'atmosphère (les gaz à effet de serre) absorbent une partie du rayonnement infrarouge émis par la surface terrestre chauffée par le Soleil, puis réémettent de la chaleur vers la surface, ce qui élève la température moyenne de la planète.

\textbf{2. Deux gaz à effet de serre.} Le dioxyde de carbone \ce{CO2} (rejeté par la combustion des énergies fossiles et la déforestation) et le méthane \ce{CH4} (rejeté par l'élevage des ruminants, les rizières et les décharges). On peut aussi citer la vapeur d'eau \ce{H2O} et le protoxyde d'azote \ce{N2O}.

\textbf{3. Rôle indispensable.} Sans effet de serre, la totalité du rayonnement infrarouge repartirait vers l'espace : la température moyenne de la Terre serait d'environ $-18\ ^\circ\mathrm{C}$, et l'eau y serait presque partout gelée. Grâce au piégeage d'une partie de la chaleur, la température moyenne est d'environ $15\ ^\circ\mathrm{C}$, ce qui permet l'existence de l'eau liquide et de la vie.

\textbf{Conclusion.} L'effet de serre est un mécanisme naturel et vital de régulation thermique ; ce n'est que son \emph{renforcement} par les émissions humaines de GES qui provoque le réchauffement climatique actuel.
\end{exoresolu}

\begin{savaistu}[L'effet de serre sur Vénus et Mars]
Vénus, avec une atmosphère à $96\ \%$ de \ce{CO2} et une pression de $92$ bars, subit un effet de serre « emballé » : température de surface $\approx 460\ ^\circ\mathrm{C}$. Mars, atmosphère ténue ($95\ \%$ \ce{CO2} mais $0,006$ bar), a un effet de serre très faible : température moyenne $-60\ ^\circ\mathrm{C}$. La Terre est dans la « zone habitable » grâce à une atmosphère « juste assez » épaisse et riche en GES pour maintenir l'eau liquide.
\end{savaistu}

\coursec{Origines anthropiques du renforcement de l'effet de serre}

\emph{Situation problème :} Depuis le début de l'ère industrielle (vers 1750), la concentration atmosphérique en \ce{CO2} est passée de $\approx 280$ ppm à plus de $420$ ppm en 2024 ($+50\ \%$). Le méthane a plus que doublé. Quelle est l'origine de cette hausse brutale ? Les activités humaines émettent des GES plus vite que les puits naturels (océans, forêts) ne peuvent les absorber.

\courssub{Méthode 2 : Relier une activité humaine à une émission de gaz à effet de serre}

\begin{methode}[Établir la chaîne activité humaine $\rightarrow$ émission de GES]
\begin{enumerate}
\item \textbf{Identifier l'activité} décrite dans l'énoncé (transport, agriculture, industrie, déforestation, gestion des déchets...).
\item \textbf{Repérer le processus physique, chimique ou biologique} mis en jeu : combustion (moteurs, centrales thermiques, feux de brousse), fermentation entérique des ruminants, décomposition anaérobie de la matière organique, déforestation (brûlis + perte de puits).
\item \textbf{Nommer le GES émis} avec sa formule : \ce{CO2} (combustions, déforestation), \ce{CH4} (élevage, rizières, décharges), \ce{N2O} (engrais azotés, combustion).
\item \textbf{Expliquer le mécanisme d'émission} en une phrase complète. Exemple : « La combustion de l'essence dans les moteurs réagit avec le dioxygène et libère du dioxyde de carbone : \ce{C_xH_y + O2 -> CO2 + H2O} (simplifié). »
\item \textbf{Relier à l'effet global} : cette émission s'ajoute aux GES naturels et \cle{renforce} l'effet de serre.
\end{enumerate}
\textbf{Piège classique} : ne pas confondre \emph{déforestation} (double effet : émission de \ce{CO2} lors des brûlis \emph{et} réduction de l'absorption de \ce{CO2} par la photosynthèse) avec une simple émission.
\end{methode}

\begin{exoresolu}[Activités humaines et émissions au Cameroun]
Voici quatre activités observées au Cameroun : (a) le trafic de motos-taxis et camions à Yaoundé ; (b) l'élevage de bovins dans l'Adamaoua ; (c) les feux de brousse et la culture sur brûlis ; (d) les décharges sauvages d'ordures en périphérie de Douala.
Pour chaque activité, indiquer le principal gaz à effet de serre émis et expliquer le mécanisme de l'émission.
\tcblower
\textbf{(a) Transport routier.} GES principal : dioxyde de carbone \ce{CO2}. Mécanisme : la combustion de l'essence ou du gasoil (hydrocarbures) avec le dioxygène de l'air dans les moteurs libère du \ce{CO2} et de la vapeur d'eau. Plus le trafic est dense et les véhicules anciens, plus les émissions par habitant sont élevées.

\textbf{(b) Élevage de bovins.} GES principal : méthane \ce{CH4}. Mécanisme : la \cle{fermentation entérique} des végétaux dans le rumen des bovins, réalisée par des archées méthanogènes, produit du méthane rejeté par éructations. Le méthane a un pouvoir de réchauffement global (PRG) $28$ à $34$ fois supérieur à celui du \ce{CO2} sur $100$ ans.

\textbf{(c) Feux de brousse et culture sur brûlis.} GES principal : \ce{CO2} (avec aussi du \ce{CH4} et du \ce{N2O} dus à la combustion incomplète). Mécanisme : la combustion de la biomasse végétale transforme le carbone stocké dans les plantes en \ce{CO2}. De plus, la destruction du couvert végétal réduit l'absorption de \ce{CO2} par la \cle{photosynthèse} : c'est un double effet aggravant (émission + perte de puits).

\textbf{(d) Décharges sauvages.} GES principal : méthane \ce{CH4}. Mécanisme : la décomposition de la matière organique (restes alimentaires, déchets verts) par des bactéries en l'absence de dioxygène (\cle{décomposition anaérobie}) dégage du méthane (biogaz).

\textbf{Conclusion.} Chaque secteur d'activité humaine possède ses sources propres de GES ; identifier le processus (combustion, fermentation, décomposition anaérobie, déforestation) permet de nommer correctement le gaz émis et de cibler ensuite les solutions.
\end{exoresolu}

\begin{aretenir}[Bilan anthropique global (données GIEC AR6)]
\begin{itemize}
\item \ce{CO2} fossile (énergie, industrie, transport) : $\approx 65\ \%$ du forçage radiatif anthropique total.
\item \ce{CO2} déforestation/changement d'affectation des sols : $\approx 11\ \%$.
\item \ce{CH4} (élevage, riz, déchets, fuites gaz/fossiles) : $\approx 18\ \%$.
\item \ce{N2O} (engrais, industrie) : $\approx 4\ \%$.
\item Les émissions anthropiques totales $\approx 59$ Gt\ce{CO2}eq/an (2019). Les puits naturels (océan + terres) absorbent environ la moitié ; le reste s'accumule dans l'atmosphère.
\end{itemize}
\end{aretenir}

\coursec{Conséquences du réchauffement climatique}

\emph{Observation :} Au Cameroun, l'Extrême-Nord subit une désertification accélérée et des sécheresses récurrentes ; le littoral (Kribi, Limbé) enregistre des érosions côtières et des inondations marines plus fréquentes ; dans l'Ouest, le paludisme gagne les hautes terres. Ces signaux locaux sont la signature du réchauffement global ($+1,1\ ^\circ\mathrm{C}$ depuis 1850--1900).

\courssub{Méthode 3 : Analyser les conséquences du réchauffement climatique dans un document}

\begin{methode}[Identifier et expliquer une conséquence du réchauffement]
\begin{enumerate}
\item \textbf{Classer la conséquence} dans la bonne catégorie : climatique (sécheresses, pluies extrêmes, cyclones), océanique (élévation du niveau des mers, acidification), écologique (recul de la biodiversité, déplacement des espèces), sanitaire (extension du paludisme vers les hautes terres, canicules) ou agricole (baisse des rendements, décalage des saisons).
\item \textbf{Expliquer la chaîne causale} : réchauffement $\rightarrow$ mécanisme physique ou biologique $\rightarrow$ conséquence observée. Exemple : réchauffement $\rightarrow$ fonte des glaciers et dilatation thermique de l'eau de mer $\rightarrow$ élévation du niveau des mers $\rightarrow$ inondation des zones côtières (Kribi, Limbé).
\item \textbf{Appuyer sur l'exemple local} quand l'énoncé le permet : avancée du désert dans l'Extrême-Nord, perturbation des saisons pluvieuses pour les cultures de cacao et de maïs, présence du moustique vecteur du paludisme en altitude (Mont Cameroun, région de l'Ouest).
\item \textbf{Rédiger en phrases complètes} en nommant précisément le mécanisme, sans se contenter d'une liste de mots.
\end{enumerate}
\end{methode}

\begin{popfigure}
\begin{tikzpicture}[scale=0.85, every node/.style={font=\small}]
% Sea level rise mechanism
% Water
\fill[popblue!30] (-4,-2) rectangle (4,0);
\draw[thick, popblue] (-4,0) -- (4,0) node[midway, above] {Niveau initial de la mer};
% Land
\fill[poporange!50!brown] (-4,-2) rectangle (-1,1);
\fill[poporange!50!brown] (1,-2) rectangle (4,1);
\node at (-2.5,1.5) {Terre continentale};
\node at (2.5,1.5) {Terre continentale};
% City
\draw[fill=popdark!80] (-1.5,0) rectangle (-0.5,1) node[midway, white, rotate=90] {Ville côtière};
\draw[fill=popdark!80] (0.5,0) rectangle (1.5,1) node[midway, white, rotate=90] {Ville côtière};
% Glacier
\draw[fill=popblue!80!white] (-3.5,1) rectangle (-2.5,2.5) node[midway, align=center] {Glacier/\\Calotte polaire};
\draw[fill=popblue!80!white] (2.5,1) rectangle (3.5,2.5) node[midway, align=center] {Glacier/\\Calotte polaire};
% Arrows for melting
\draw[popblue, thick, ->] (-3,1) -- (-3,0.2) node[midway, right] {Fonte};
\draw[popblue, thick, ->] (3,1) -- (3,0.2);
% Thermal expansion
\draw[poporange, thick, <->] (0,-1.5) -- (0,-0.5) node[midway, right] {Dilatation thermique};
% New sea level
\draw[thick, poppink, dashed] (-4,0.8) -- (4,0.8) node[midway, above] {Nouveau niveau (montée)};
% Flooding
\fill[poppink!30, opacity=0.5] (-1.5,0) rectangle (-0.5,0.8);
\fill[poppink!30, opacity=0.5] (0.5,0) rectangle (1.5,0.8);
\node[poppink, align=center] at (0,1.5) {Inondation côtière \\ (Kribi, Limbé, Douala)};
\end{tikzpicture}
\legende{Mécanismes de l'élévation du niveau des mers : fonte des glaces continentales (apport d'eau) et dilatation thermique de l'océan (volume accru). Conséquence directe : submersion des zones basses camerounaises.}
\end{popfigure}

\begin{exoresolu}[Conséquences du réchauffement au Cameroun]
Dans l'Extrême-Nord du Cameroun, les paysans observent depuis quelques décennies des saisons des pluies plus courtes et imprévisibles, l'avancée des zones arides et des épisodes de sécheresse répétés ; sur la côte, les habitants de certaines localités signalent des inondations marines plus fréquentes.
\begin{enumerate}
\item Relier chacune de ces observations au réchauffement climatique en expliquant le mécanisme.
\item Citer une conséquence sanitaire possible du réchauffement dans les hautes terres de l'Ouest camerounais.
\end{enumerate}
\tcblower
\textbf{1. Explication des observations.}

\emph{Sécheresses et avancée des zones arides (Extrême-Nord) :} le réchauffement climatique modifie la circulation atmosphérique et le régime des pluies ; l'augmentation des températures accroît l'\cle{évapotranspiration} des sols et des plantes, ce qui assèche les terres, réduit la durée de la saison des pluies et favorise la \cle{désertification} (avancée du Sahel vers le sud).

\emph{Inondations marines sur la côte :} le réchauffement provoque la \cle{fonte des glaciers} et des calottes polaires (Groenland, Antarctique), ainsi que la \cle{dilatation thermique} de l'eau des océans (l'eau chaude occupe un volume plus grand). Ces deux mécanismes élèvent le niveau moyen des mers (environ $3,3$ mm/an actuellement), ce qui augmente la fréquence des inondations des zones côtières basses.

\textbf{2. Conséquence sanitaire.} Avec des températures plus élevées en altitude, les moustiques du genre \emph{Anopheles}, vecteurs du \cle{paludisme}, peuvent survivre et se reproduire dans les hautes terres de l'Ouest autrefois trop fraîches (ex. Dschang, Bafoussam, pentes du Mont Cameroun) : le paludisme s'étend alors à des populations nouvelles, non immunisées.

\textbf{Conclusion.} Le réchauffement climatique agit par des mécanismes physiques (évaporation, fonte des glaces, dilatation de l'eau) et biologiques (extension des vecteurs de maladies) dont les conséquences touchent directement l'agriculture, les populations côtières et la santé publique au Cameroun.
\end{exoresolu}

\begin{savaistu}[Acidification des océans : l'autre problème du \ce{CO2}]
Environ $25\ \%$ du \ce{CO2} émis par l'homme est absorbé par les océans. Il y réagit avec l'eau : \ce{CO2 + H2O <=> H2CO3 <=> H+ + HCO3-}. L'augmentation des ions \ce{H+} abaisse le pH (acidification). Cela réduit la disponibilité en ions carbonate \ce{CO3^2-}, indispensables aux organismes calcifiants (coraux, mollusques, plancton) pour construire leur squelette (\ce{CaCO3}). Au Cameroun, cela menace les écosystèmes de mangroves et la pêche côtière.
\end{savaistu}

\coursec{Moyens d'action : atténuation et adaptation}

\emph{Problématique :} Face au réchauffement, deux stratégies complémentaires s'imposent : l'\cle{atténuation} (mitigation), qui s'attaque aux causes en réduisant les émissions nettes de GES, et l'\cle{adaptation}, qui réduit la vulnérabilité face aux effets déjà en cours.

\begin{definition}[Atténuation vs Adaptation]
\begin{itemize}
\item \textbf{Atténuation (Mitigation) :} Actions visant à \emph{réduire les sources} de GES (efficacité énergétique, énergies renouvelables, changement de mobilité) ou à \emph{augmenter les puits} de carbone (reboisement, restauration des sols, capture géologique). Agit sur la \emph{cause}.
\item \textbf{Adaptation :} Actions visant à \emph{réduire la vulnérabilité} aux impacts actuels ou futurs (digues, variétés résistantes à la sécheresse, systèmes d'alerte précoce, urbanisation résiliente). Agit sur les \emph{effets}.
\end{itemize}
\end{definition}

\courssub{Méthode 4 : Proposer des gestes et des politiques de réduction des émissions}

\begin{methode}[Construire une proposition d'action (atténuation ou adaptation)]
\begin{enumerate}
\item \textbf{Distinguer les deux types d'action} :
\begin{itemize}
\item \cle{Atténuation} : réduire les émissions de GES ou augmenter leur absorption (agir sur la \emph{cause}) ;
\item \cle{Adaptation} : s'organiser pour vivre avec les conséquences déjà présentes (agir sur les \emph{effets}).
\end{itemize}
\item \textbf{Partir de la source d'émission} identifiée (transport, énergie, déforestation, élevage, déchets) et proposer une action qui la réduit directement.
\item \textbf{Préciser l'échelle} : geste individuel (éteindre les appareils, privilégier la marche et le transport en commun, trier les déchets, planter des arbres), collectivité (reboisement communal, gestion des décharges avec captage du méthane), État (énergies renouvelables comme l'hydroélectricité et le solaire, transports publics, lois sur la déforestation), coopération internationale (Accord de Paris, Mécanisme pour un développement propre).
\item \textbf{Justifier chaque proposition} en la reliant au GES visé : « planter des arbres augmente l'absorption de \ce{CO2} par la photosynthèse ».
\item \textbf{Vérifier la cohérence} : une action d'adaptation (digues, variétés de semences résistantes à la sécheresse) ne réduit pas les émissions ; ne pas la présenter comme une atténuation.
\end{enumerate}
\end{methode}

\begin{exoresolu}[Plan de réduction des émissions pour une ville camerounaise]
La mairie d'une grande ville camerounaise souhaite réduire ses émissions de gaz à effet de serre. Les principales sources identifiées sont : le transport routier (vieux véhicules, embouteillages), la production d'électricité par des groupes électrogènes au gasoil, les décharges sauvages et la disparition des espaces verts.
Proposer quatre actions concrètes, en précisant pour chacune le GES visé, le mécanisme de réduction et s'il s'agit d'atténuation ou d'adaptation.
\tcblower
\textbf{Action 1 — Développer les transports en commun et renouveler le parc de véhicules.} GES visé : \ce{CO2}. Mécanisme : moins de véhicules individuels en circulation et des moteurs plus propres (normes Euro 4/5, électrique) réduisent la combustion d'essence et de gasoil, donc les émissions de \ce{CO2}. Type : \cle{atténuation}.

\textbf{Action 2 — Remplacer les groupes électrogènes par l'hydroélectricité et le solaire.} GES visé : \ce{CO2}. Mécanisme : les énergies renouvelables (barrages hydroélectriques comme ceux de la Sanaga -- Song Loulou, Edea, Nachtigal --, panneaux solaires photovoltaïques) produisent de l'électricité sans combustion de gasoil. Type : \cle{atténuation}.

\textbf{Action 3 — Aménager les décharges avec captage du méthane et tri-compostage.} GES visé : \ce{CH4}. Mécanisme : le compostage aérobie des déchets organiques et le captage du biogaz évitent la décomposition anaérobie qui dégage du méthane ; le biogaz capté peut même servir de combustible (économie circulaire). Type : \cle{atténuation}.

\textbf{Action 4 — Reboiser les espaces verts et planter des arbres en ville.} GES visé : \ce{CO2}. Mécanisme : par la \cle{photosynthèse}, les arbres absorbent le \ce{CO2} atmosphérique et stockent le carbone dans leur biomasse (bois, racines, sols) ; ils constituent un \cle{puits de carbone}. Type : \cle{atténuation} (augmentation de l'absorption).

\textbf{Conclusion.} Un plan efficace combine plusieurs actions d'atténuation ciblant chaque source d'émission ; il peut être complété par des mesures d'adaptation (digues, drainage des quartiers inondables, centres de santé climatisés) pour protéger la population des effets déjà sensibles du réchauffement.
\end{exoresolu}

\begin{savaistu}[Le Cameroun face au climat : engagements et atouts]
\begin{itemize}
\item \textbf{CDN (Contribution Déterminée au niveau National) :} Le Cameroun s'engage à réduire ses émissions de $35\ \%$ d'ici 2030 (scénario conditionnel) par rapport au scénario tendanciel, via la forêt (REDD+), l'énergie (hydroélectricité, solaire), l'agriculture climato-intelligente.
\item \textbf{Atout hydroélectrique :} Le potentiel hydroélectrique (Sanaga, Nachtigal, Grand Eweng) est le 2ème d'Afrique. Son exploitation réduit la dépendance aux centrales thermiques (gasoil/fuel lourd).
\item \textbf{Initiative « Sahel Vert » / Grande Muraille Verte :} Lutte contre la désertification dans l'Extrême-Nord par reboisement et gestion durable des terres (adaptation + atténuation).
\item \textbf{Forêt du Bassin du Congo :} 2ème poumon vert mondial. La lutte contre la déforestation illégale et la certification FSC du bois sont des leviers majeurs d'atténuation.
\end{itemize}
\end{savaistu}

\coursec{Synthèse, modélisation simple et engagement citoyen}

\emph{Synthèse :} L'effet de serre naturel est vital. Son renforcement par les émissions anthropiques (\ce{CO2}, \ce{CH4}, \ce{N2O}) cause un réchauffement global ($+1,1\ ^\circ\mathrm{C}$) aux conséquences multiples (climatiques, océaniques, écologiques, sanitaires, agricoles). La riposte combine atténuation (réduction émissions / augmentation puits) et adaptation (résilience). Chaque acteur (citoyen, collectivité, État, entreprise) a un rôle à jouer.

\courssub{Méthode 5 : Exploiter des données chiffrées sur les émissions et les températures}

\begin{methode}[Lire et interpréter un tableau ou une courbe de données climatiques]
\begin{enumerate}
\item \textbf{Lire les axes ou les colonnes} avec leurs unités (années, concentration de \ce{CO2} en ppm, température en $^\circ\mathrm{C}$, émissions en tonnes).
\item \textbf{Repérer la tendance} : croissance, stabilité, décroissance ; comparer la valeur initiale et la valeur finale.
\item \textbf{Calculer une variation} : variation absolue $\Delta = V_f - V_i$ ; variation relative (pourcentage) :
\[ t = \frac{V_f - V_i}{V_i} \times 100 \quad (\text{en } \%). \]
\item \textbf{Mettre en relation deux grandeurs} : par exemple l'augmentation de la concentration de \ce{CO2} et celle de la température moyenne, en précisant que la \emph{corrélation} observée s'explique par le mécanisme de l'effet de serre.
\item \textbf{Conclure par une phrase} reliant les données au phénomène étudié (renforcement de l'effet de serre).
\end{enumerate}
\textbf{Attention aux unités} : ppm = partie par million ; $1\ ^\circ\mathrm{C}$ d'augmentation moyenne globale est une valeur considérable à l'échelle du climat.
\end{methode}

\begin{exoresolu}[Évolution de la concentration en \ce{CO2} et de la température]
Le tableau ci-dessous donne la concentration atmosphérique moyenne en dioxyde de carbone et l'anomalie de température moyenne globale (écart par rapport à la moyenne 1951--1980).
\begin{center}
\begin{tabularx}{\linewidth}{|l|X|X|X|X|}
\hline
\rowcolor{popblueL} \textbf{Année} & \textbf{1960} & \textbf{1980} & \textbf{2000} & \textbf{2020} \\
\hline
Concentration en \ce{CO2} (ppm) & 317 & 339 & 369 & 414 \\
\hline
Anomalie de température ($^\circ\mathrm{C}$) & $-0{,}03$ & $+0{,}26$ & $+0{,}40$ & $+1{,}02$ \\
\hline
\end{tabularx}
\end{center}
\begin{enumerate}
\item Calculer l'augmentation de la concentration en \ce{CO2} entre 1960 et 2020, en valeur absolue puis en pourcentage.
\item Comparer l'évolution de la température à celle du \ce{CO2} et l'expliquer.
\end{enumerate}
\tcblower
\textbf{1. Augmentation de la concentration en \ce{CO2}.}

Variation absolue :
\[ \Delta = V_f - V_i = 414 - 317 = 97 \ \text{ppm}. \]

Variation relative :
\[ t = \frac{V_f - V_i}{V_i} \times 100 = \frac{97}{317} \times 100 \approx 30{,}6\ \%. \]

Entre 1960 et 2020, la concentration atmosphérique en \ce{CO2} a augmenté d'environ $97$ ppm, soit une hausse d'environ $31\ \%$.

\textbf{2. Comparaison et explication.} Sur la même période, l'anomalie de température passe de $-0{,}03\ ^\circ\mathrm{C}$ à $+1{,}02\ ^\circ\mathrm{C}$, soit un réchauffement d'environ $1{,}05\ ^\circ\mathrm{C}$. Les deux grandeurs augmentent donc simultanément : il y a \cle{corrélation} entre la concentration en \ce{CO2} et la température moyenne globale. Explication : le \ce{CO2} est un gaz à effet de serre ; l'augmentation de sa concentration, due aux activités humaines (combustion des énergies fossiles, déforestation), accroît l'absorption du rayonnement infrarouge émis par la Terre : c'est le \cle{renforcement de l'effet de serre}, responsable du réchauffement climatique.

\textbf{Conclusion.} Les données montrent une hausse de $31\ \%$ du \ce{CO2} atmosphérique accompagnée d'un réchauffement d'environ $1\ ^\circ\mathrm{C}$ en 60 ans, ce qui illustre quantitativement le lien entre émissions de GES et réchauffement climatique.
\end{exoresolu}

\courssub{Méthode 6 : Construire une modélisation simple (bilan carbone) et rédiger un argumentaire citoyen}

\begin{methode}[Réaliser un bilan carbone simplifié et argumenter]
\begin{enumerate}
\item \textbf{Lister les postes d'émission} de la situation (transport, alimentation, énergie, déchets) et relever les quantités données dans l'énoncé.
\item \textbf{Calculer les émissions de chaque poste} : émission = quantité consommée $\times$ facteur d'émission (donné). Additionner pour obtenir le total :
\[ E_{\text{total}} = \sum_i q_i \times f_i. \]
\item \textbf{Identifier le poste dominant} et proposer en priorité une action sur ce poste (efficacité maximale).
\item \textbf{Soustraire les absorptions} éventuelles (arbres plantés, sols restaurés) pour obtenir le bilan net.
\item \textbf{Rédiger l'argumentaire citoyen} : fait chiffré $\rightarrow$ explication scientifique $\rightarrow$ proposition concrète $\rightarrow$ appel à l'action. Ton respectueux, phrases courtes, vocabulaire scientifique correct.
\end{enumerate}
\end{methode}

\begin{exoresolu}[Bilan carbone simplifié d'un lycée]
Un lycée de Yaoundé relève sur un mois : $500$ litres de gasoil pour son groupe électrogène (facteur d'émission : $2{,}7\ \mathrm{kg}$ de \ce{CO2} par litre) ; $2\,000\ \mathrm{km}$ parcourus par les véhicules de service (facteur : $0{,}20\ \mathrm{kg}$ de \ce{CO2} par km). Par ailleurs, les élèves plantent des arbres qui absorberont environ $150\ \mathrm{kg}$ de \ce{CO2} par mois.
\begin{enumerate}
\item Calculer les émissions mensuelles totales de \ce{CO2} du lycée.
\item Calculer le bilan net après absorption par les arbres.
\item Rédiger en quelques phrases un argumentaire pour convaincre le proviseur d'installer des panneaux solaires.
\end{enumerate}
\tcblower
\textbf{1. Émissions mensuelles totales.}

Poste groupe électrogène :
\[ E_1 = 500 \times 2{,}7 = 1\,350\ \mathrm{kg\ de\ } \ce{CO2}. \]

Poste transport :
\[ E_2 = 2\,000 \times 0{,}20 = 400\ \mathrm{kg\ de\ } \ce{CO2}. \]

Total :
\[ E_{\text{total}} = E_1 + E_2 = 1\,350 + 400 = 1\,750\ \mathrm{kg\ de\ } \ce{CO2} \ \text{par mois}. \]

\textbf{2. Bilan net.} On soustrait l'absorption par les arbres :
\[ E_{\text{net}} = 1\,750 - 150 = 1\,600\ \mathrm{kg\ de\ } \ce{CO2} \ \text{par mois}. \]

\textbf{3. Argumentaire.} « Notre lycée rejette chaque mois $1\,750\ \mathrm{kg}$ de \ce{CO2}, dont $1\,350\ \mathrm{kg}$, soit plus des trois quarts ($77\ \%$), proviennent du seul groupe électrogène au gasoil. Ce dioxyde de carbone renforce l'effet de serre et contribue au réchauffement climatique dont nous subissons déjà les effets (sécheresses, inondations, santé). En installant des panneaux solaires, nous produirions une électricité propre, sans combustion, et nous supprimerions l'essentiel de ces émissions, tout en réalisant des économies de carburant. Engageons notre établissement dans la lutte contre le réchauffement climatique ! »

\textbf{Conclusion.} Le poste « groupe électrogène » domine largement le bilan ($77\ \%$ des émissions) : c'est sur lui qu'une action d'atténuation (solaire) sera la plus efficace ; la plantation d'arbres, utile, ne compense qu'environ $9\ \%$ des émissions.
\end{exoresolu}

\begin{attention}[Les 5 erreurs qui coûtent des points au Probatoire]
\begin{enumerate}
\item \textbf{Confondre effet de serre et renforcement de l'effet de serre.} L'effet de serre \emph{naturel} est indispensable à la vie (température moyenne de $15\ ^\circ\mathrm{C}$ au lieu de $-18\ ^\circ\mathrm{C}$) : écrire que « l'effet de serre est une catastrophe à supprimer » est une erreur conceptuelle grave. C'est son \emph{renforcement} par les émissions humaines qui pose problème.
\item \textbf{Confondre les gaz.} Le trou dans la couche d'ozone n'est \emph{pas} la cause de l'effet de serre ; les GES à citer sont \ce{CO2}, \ce{CH4}, \ce{N2O} et la vapeur d'eau, avec leurs formules correctes. Écrire « le CO » ou « l'oxygène » comme gaz à effet de serre est faux.
\item \textbf{Oublier le mécanisme dans l'explication.} Une liste de conséquences sans chaîne causale ne rapporte pas tous les points : il faut relier (réchauffement $\rightarrow$ fonte des glaces \emph{et} dilatation de l'eau $\rightarrow$ montée du niveau des mers). Nommer le mécanisme physique ou biologique est exigé.
\item \textbf{Erreurs de calcul sur les pourcentages.} La variation relative se calcule par rapport à la valeur \emph{initiale} : $t = \dfrac{V_f - V_i}{V_i} \times 100$, et non par rapport à $V_f$. Vérifier aussi la cohérence des unités (kg de \ce{CO2}, ppm, $^\circ\mathrm{C}$) et donner le résultat avec son unité.
\item \textbf{Mélanger atténuation et adaptation.} Construire une digue ou cultiver des variétés résistantes à la sécheresse, ce sont des mesures d'\emph{adaptation} : elles ne réduisent aucune émission. Seules les actions qui diminuent les émissions ou augmentent l'absorption des GES (énergies renouvelables, reboisement, transports propres) sont des mesures d'\emph{atténuation}.
\end{enumerate}
\end{attention}

\courssub{Entraînement au Probatoire : Exercices types}

\begin{exercice}[Exercice 1 : Analyse de document — Émissions sectorielles]
Le graphique ci-dessous représente la répartition des émissions mondiales de gaz à effet de serre par secteur d'activité en 2019 (total : $59$ Gt\ce{CO2}eq).
\begin{center}
\begin{tikzpicture}
\begin{axis}[
    ybar stacked,
    bar width=15pt,
    width=\linewidth,
    height=7cm,
    xlabel={Secteurs},
    ylabel={Émissions (Gt\ce{CO2}eq)},
    symbolic x coords={Énergie, Industrie, AFOLU, Transport, Bâtiment, Déchets},
    xtick=data,
    x tick label style={rotate=45, anchor=east},
    legend style={at={(0.5,-0.25)}, anchor=north, legend columns=-1},
    ymin=0, ymax=25
]
\addplot[fill=popblue] coordinates {(Énergie,15) (Industrie,8) (AFOLU,2) (Transport,2) (Bâtiment,1) (Déchets,0.5)};
\addplot[fill=poporange] coordinates {(Énergie,5) (Industrie,6) (AFOLU,1) (Transport,6) (Bâtiment,2) (Déchets,0.5)};
\addplot[fill=popgreen] coordinates {(Énergie,0) (Industrie,0) (AFOLU,10) (Transport,0) (Bâtiment,0) (Déchets,0)};
\addplot[fill=poppurple] coordinates {(Énergie,0) (Industrie,0) (AFOLU,5) (Transport,0) (Bâtiment,0) (Déchets,3)};
\legend{\ce{CO2} fossile, \ce{CO2} procédés, \ce{CH4}, \ce{N2O}}
\end{axis}
\end{tikzpicture}
\end{center}
\textbf{AFOLU} = Agriculture, Foresterie et Autres Utilisations des Terres.
\begin{enumerate}
\item Identifier le secteur émetteur dominant et le principal gaz associé.
\item Expliquer pourquoi le secteur AFOLU émet à la fois du \ce{CO2}, du \ce{CH4} et du \ce{N2O} (citer les processus pour chaque gaz).
\item Calculer le pourcentage des émissions totales dues au \ce{CH4} (somme sur tous les secteurs).
\end{enumerate}
\end{exercice}

\begin{corrige}[Exercice 1]
\textbf{1.} Le secteur dominant est \textbf{Énergie} ($\approx 20$ Gt\ce{CO2}eq, soit $\approx 34\ \%$). Le principal gaz est le \ce{CO2} d'origine fossile (combustion charbon, pétrole, gaz).
\textbf{2.} AFOLU : \ce{CO2} (déforestation, brûlis, dégradation des sols) ; \ce{CH4} (fermentation entérique ruminants, rizières inondées) ; \ce{N2O} (utilisation d'engrais azotés, déjections animales, brûlis).
\textbf{3.} \ce{CH4} total = $2$ (Énergie) + $1$ (Industrie) + $10$ (AFOLU) + $0$ + $0$ + $3$ (Déchets) = $16$ Gt\ce{CO2}eq. Pourcentage = $\frac{16}{59} \times 100 \approx 27,1\ \%$.
\end{corrige}

\begin{exercice}[Exercice 2 : Argumentation — Projet d'aménagement]
La mairie de Maroua (Extrême-Nord) prévoit de créer une ceinture verte de $500$ ha d'arbres (espèces sahéliennes type \emph{Acacia senegal}) autour de la ville pour lutter contre l'avancée du sable et capter du carbone. Le facteur de séquestration moyen est estimé à $5$ t\ce{CO2}/ha/an.
\begin{enumerate}
\item Calculer la quantité de \ce{CO2} séquestrée annuellement par ce projet.
\item Classer cette action en atténuation ou adaptation, et justifier.
\item Proposer une mesure complémentaire d'adaptation pour la population face aux canicules.
\end{enumerate}
\end{exercice}

\begin{corrige}[Exercice 2]
\textbf{1.} Séquestration = $500\ \mathrm{ha} \times 5\ \mathrm{t\ce{CO2}/ha/an} = 2\,500\ \mathrm{t\ce{CO2}/an}$ (soit $2\,500\,000\ \mathrm{kg}$).
\textbf{2.} C'est une action d'\textbf{atténuation} (augmentation d'un puits de carbone par photosynthèse) \textbf{ET} d'\textbf{adaptation} (frein à l'érosion éolienne, ombrage, microclimat local, ressources non-ligneuses). Double bénéfice.
\textbf{3.} Mesure d'adaptation : création de « centres de fraîcheur » (bâtiments bioclimatiques, points d'eau, végétalisation dense) accessibles aux personnes vulnérables lors des pics de chaleur ($>40\ ^\circ\mathrm{C}$) ; ou système d'alerte précoce canicule couplé à une distribution d'eau.
\end{corrige}

\begin{exercice}[Exercice 3 : Modélisation — Bilan carbone d'un ménage]
Un ménage de Douala consomme par an : $1\,200$ kWh d'électricité (mix camerounais : facteur $0{,}45\ \mathrm{kg\ce{CO2}/kWh}$), $800$ litres d'essence pour la voiture (facteur $2{,}3\ \mathrm{kg\ce{CO2}/L}$), $200$ kg de viande bovine (facteur $27\ \mathrm{kg\ce{CO2}eq/kg}$). Il composte ses déchets organiques (évite $150\ \mathrm{kg\ce{CO2}eq}$ d'émission de décharge).
\begin{enumerate}
\item Calculer l'empreinte carbone brute annuelle du ménage (en t\ce{CO2}eq).
\item Calculer l'empreinte nette après compostage.
\item Identifier le poste prioritaire pour réduire l'empreinte et proposer une action concrète.
\end{enumerate}
\end{exercice}

\begin{corrige}[Exercice 3]
\textbf{1.} Électricité : $1\,200 \times 0{,}45 = 540\ \mathrm{kg}$. Voiture : $800 \times 2{,}3 = 1\,840\ \mathrm{kg}$. Viande : $200 \times 27 = 5\,400\ \mathrm{kg}$. Total brut = $540 + 1\,840 + 5\,400 = 7\,780\ \mathrm{kg} = 7{,}78\ \mathrm{t\ce{CO2}eq}$.
\textbf{2.} Net = $7\,780 - 150 = 7\,630\ \mathrm{kg} = 7{,}63\ \mathrm{t\ce{CO2}eq}$.
\textbf{3.} Poste dominant : \textbf{Viande bovine} ($69\ \%$ des émissions). Action prioritaire : réduire la consommation de viande rouge (ex. $1$ jour sans viande/semaine, substitution volaille/légumineuses) ; impact massif sur \ce{CH4} et \ce{N2O} de l'élevage. Action secondaire : covoiturage/transport en commun pour la voiture.
\end{corrige}

\newpage

\coursec{Exercices}

\courssub{Je vérifie mes connaissances}

\begin{exercice}[QCM : l'effet de serre]
Pour chaque question, choisir la (ou les) bonne(s) réponse(s).
\begin{enumerate}
\item L'effet de serre naturel est :
\begin{enumerate}
\item un phénomène apparu avec l'industrialisation ;
\item un phénomène naturel indispensable à la vie sur Terre ;
\item un phénomène qui chauffe directement le sol par les rayons ultraviolets.
\end{enumerate}
\item Parmi les gaz suivants, lesquels sont des gaz à effet de serre ?
\begin{enumerate}
\item le dioxygène \ce{O2} ;
\item le dioxyde de carbone \ce{CO2} ;
\item le méthane \ce{CH4} ;
\item la vapeur d'eau \ce{H2O}.
\end{enumerate}
\item Sans effet de serre naturel, la température moyenne à la surface de la Terre serait d'environ :
\begin{enumerate}
\item $+15\ ^\circ\mathrm{C}$ ;
\item $-18\ ^\circ\mathrm{C}$ ;
\item $+33\ ^\circ\mathrm{C}$.
\end{enumerate}
\item Le rayonnement renvoyé par le sol terrestre vers l'atmosphère est principalement :
\begin{enumerate}
\item un rayonnement ultraviolet ;
\item un rayonnement visible ;
\item un rayonnement infrarouge.
\end{enumerate}
\end{enumerate}
\end{exercice}

\begin{exercice}[Vrai ou faux, à justifier]
Dire si chaque affirmation est vraie ou fausse, puis justifier en une ou deux phrases.
\begin{enumerate}
\item « L'effet de serre est un phénomène totalement artificiel créé par les activités humaines. »
\item « Le méthane est un gaz à effet de serre dont le pouvoir de réchauffement, à masse égale, est supérieur à celui du \ce{CO2}. »
\item « La déforestation contribue au renforcement de l'effet de serre uniquement parce que les feux de brousse dégagent du \ce{CO2}. »
\item « Le barrage hydroélectrique de Song Loulou émet autant de \ce{CO2} par kilowattheure qu'une centrale thermique au gaz. »
\end{enumerate}
\end{exercice}

\begin{exercice}[Définitions]
Définir en deux ou trois lignes chacun des termes suivants :
\begin{enumerate}
\item effet de serre ;
\item gaz à effet de serre (GES) ;
\item atténuation (du changement climatique) ;
\item adaptation (au changement climatique).
\end{enumerate}
\end{exercice}

\begin{exercice}[Calcul direct d'émissions]
Une moto-taxi de Douala consomme en moyenne 4 litres d'essence par jour. On donne le facteur d'émission de l'essence : 2,3 kg de \ce{CO2} par litre brûlé.
\begin{enumerate}
\item Calculer la masse de \ce{CO2} émise par cette moto en une journée de travail.
\item Calculer la masse de \ce{CO2} émise en un mois de 26 jours travaillés.
\item Un jeune arbre absorbe en moyenne 10 kg de \ce{CO2} par an. Combien d'arbres faudrait-il pour compenser les émissions annuelles (312 jours travaillés) de cette moto-taxi ?
\end{enumerate}
\end{exercice}

\courssub{Je m'entraîne}

\begin{exercice}[Relier activité et émission]
Compléter le tableau suivant en indiquant pour chaque activité humaine le principal gaz à effet de serre émis et le secteur concerné (énergie, agriculture, transport, industrie, déchets) :
\begin{enumerate}
\item riziculture inondée dans la plaine de Yagoua ;
\item décharge sauvage d'ordures ménagères à Yaoundé ;
\item circulation des véhicules à essence et gasoil ;
\item combustion du charbon de bois pour la cuisson ;
\item élevage de bovins dans l'Adamaoua.
\end{enumerate}
\end{exercice}

\begin{exercice}[Schéma du bilan radiatif]
Reproduire et compléter un schéma simplifié du bilan radiatif de la Terre en y plaçant les flèches et les légendes suivantes : rayonnement solaire incident, rayonnement réfléchi par l'atmosphère et le sol, rayonnement infrarouge émis par le sol, rayonnement infrarouge absorbé puis réémis par les gaz à effet de serre, rayonnement infrarouge s'échappant vers l'espace. Expliquer ensuite en trois lignes pourquoi une augmentation de la concentration en GES élève la température moyenne du globe.
\end{exercice}

\begin{exercice}[Bilan carbone d'un déplacement]
Un élève se rend chaque jour au lycée. Deux options : moto-taxi (facteur d'émission : 0,080 kg de \ce{CO2} par km et par passager) ou marche à pied (0 kg). La distance domicile-lycée est de 3 km.
\begin{enumerate}
\item Calculer les émissions annuelles de \ce{CO2} liées aux trajets aller-retour en moto-taxi, pour 180 jours de classe.
\item Comparer avec l'option marche à pied et citer deux autres avantages de ce choix.
\item Proposer un mode de transport intermédiaire permettant de diviser les émissions par au moins quatre, en justifiant.
\end{enumerate}
\end{exercice}

\begin{exercice}[Lecture de données climatiques]
Le tableau ci-dessous donne la concentration moyenne en \ce{CO2} atmosphérique mesurée à différentes dates.
\begin{center}
\begin{tabularx}{\linewidth}{|l|X|X|X|X|}
\hline
\rowcolor{popblueL} Année & 1960 & 1980 & 2000 & 2020 \\
\hline Concentration en \ce{CO2} (ppm) & 317 & 339 & 369 & 414 \\
\hline
\end{tabularx}
\end{center}
\begin{enumerate}
\item Calculer l'augmentation de la concentration entre 1960 et 2020, en ppm puis en pourcentage.
\item Calculer l'augmentation moyenne par décennie entre 1960 et 1980, puis entre 2000 et 2020. Que constate-t-on ?
\item Citer deux activités humaines responsables de cette évolution.
\end{enumerate}
\end{exercice}

\courssub{Je me prépare au Probatoire}

\begin{exercice}[Le bilan radiatif de la Terre]
Un document présente le bilan radiatif moyen de la planète : pour 100 unités d'énergie solaire arrivant au sommet de l'atmosphère, environ 30 unités sont réfléchies vers l'espace (nuages, surfaces claires), 70 unités sont absorbées par le sol et l'atmosphère. Le sol, chauffé, émet un rayonnement infrarouge dont une grande partie est absorbée par les gaz à effet de serre (vapeur d'eau, \ce{CO2}, \ce{CH4}, \ce{N2O}), qui réémettent à leur tour de la chaleur vers la surface.
\begin{enumerate}
\item Expliquer en cinq lignes maximum le rôle des gaz à effet de serre dans le maintien de la température moyenne de la Terre voisine de $+15\ ^\circ\mathrm{C}$.
\item Schématiser ce bilan radiatif en faisant figurer les quatre flux d'énergie cités dans le texte.
\item Indiquer deux activités humaines qui augmentent la concentration des gaz à effet de serre, en précisant pour chacune le gaz principalement émis.
\item Expliquer pourquoi on parle d'un « renforcement » et non d'une « création » de l'effet de serre par l'Homme.
\end{enumerate}
\end{exercice}

\begin{exercice}[Facture d'électricité et émissions d'un ménage]
Un ménage de Yaoundé consomme 300 kWh d'électricité par mois. Le mix électrique camerounais est composé d'environ 70 \% d'hydroélectricité (émissions négligeables) et de 30 \% de production thermique au gaz, dont le facteur d'émission est de 0,45 kg de \ce{CO2} par kWh.
\begin{enumerate}
\item Calculer la quantité d'électricité mensuelle d'origine thermique consommée par ce ménage.
\item En déduire les émissions mensuelles de \ce{CO2} de ce ménage, puis ses émissions annuelles.
\item Le ménage souhaite réduire ses émissions de 20 \%. Calculer la masse de \ce{CO2} à économiser chaque mois, puis la quantité d'électricité thermique correspondante en kWh.
\item Proposer deux gestes concrets permettant d'atteindre cette réduction, en précisant pour chacun le mécanisme d'économie d'énergie.
\item Expliquer pourquoi le développement des barrages hydroélectriques (Memve'ele, Nachtigal) contribue à limiter les émissions nationales de \ce{CO2}.
\end{enumerate}
\end{exercice}

\begin{exercice}[Émissions par secteur et plaidoyer pour l'agroforesterie]
Le tableau ci-dessous donne la répartition estimée des émissions de gaz à effet de serre du Cameroun (en pourcentage du total national) en 1990 et en 2020.
\begin{center}
\begin{tabularx}{\linewidth}{|l|X|X|}
\hline
\rowcolor{popblueL} Secteur & Part en 1990 (\%) & Part en 2020 (\%) \\
\hline Déforestation et changement d'usage des terres & 55 & 45 \\
\hline Énergie et transport & 15 & 28 \\
\hline Agriculture et élevage & 20 & 17 \\
\hline Déchets et autres & 10 & 10 \\
\hline
\end{tabularx}
\end{center}
\begin{enumerate}
\item Identifier le secteur dont la part relative a le plus augmenté entre 1990 et 2020, et calculer cette augmentation en points de pourcentage.
\item Suggérer une politique sectorielle adaptée pour limiter les émissions de ce secteur, en donnant deux mesures concrètes.
\item La déforestation reste le premier secteur émetteur. Expliquer le double rôle de la forêt (source et puits de carbone).
\item Rédiger une argumentation d'environ 150 mots à l'intention du maire d'une commune rurale pour l'inciter à intégrer l'agroforesterie (association d'arbres et de cultures) dans le plan communal de développement. L'argumentaire s'appuiera sur au moins deux arguments : le stockage du carbone par les arbres et la résilience des cultures face aux aléas climatiques (sécheresses, pluies irrégulières).
\end{enumerate}
\end{exercice}

\newpage

\coursec{Corrigés des exercices}

\coursec{Synthèse, modélisation simple et engagement citoyen}
\courssub{Banque d'exercices corrigés type Probatoire}

\begin{corrige}[Exercice 1 — QCM : l'effet de serre]
\begin{enumerate}
\item \textbf{Bonne réponse : b.} L'effet de serre est un phénomène \cle{naturel}, dû à la présence de gaz à effet de serre dans l'atmosphère depuis des temps géologiques. Sans lui, la température moyenne du globe serait d'environ $-18\ ^\circ\mathrm{C}$ et la vie telle que nous la connaissons serait impossible. Il n'est donc pas apparu avec l'industrialisation : c'est son \emph{renforcement} qui est lié aux activités humaines. Le sol n'est pas chauffé directement par les ultraviolets mais surtout par le rayonnement visible du Soleil.
\item \textbf{Bonnes réponses : b, c et d.} Le \ce{CO2}, le méthane \ce{CH4} et la vapeur d'eau \ce{H2O} absorbent le rayonnement infrarouge émis par le sol : ce sont des gaz à effet de serre. Le dioxygène \ce{O2}, molécule diatomique symétrique, n'absorbe pas l'infrarouge : ce n'est pas un GES.
\item \textbf{Bonne réponse : b.} Sans effet de serre naturel, la température moyenne serait d'environ $-18\ ^\circ\mathrm{C}$, soit $33\ ^\circ\mathrm{C}$ de moins que la valeur actuelle ($+15\ ^\circ\mathrm{C}$).
\item \textbf{Bonne réponse : c.} Tout corps chauffé émet un rayonnement dont la longueur d'onde dépend de sa température. Le sol terrestre, à environ $+15\ ^\circ\mathrm{C}$, émet dans l'\cle{infrarouge} (rayonnement de grande longueur d'onde), contrairement au Soleil, très chaud, qui émet surtout dans le visible.
\end{enumerate}
\end{corrige}

\begin{corrige}[Exercice 2 — Vrai ou faux, à justifier]
\begin{enumerate}
\item \textbf{Faux.} L'effet de serre est un phénomène naturel indispensable : il maintient la température moyenne du globe à $+15\ ^\circ\mathrm{C}$ au lieu de $-18\ ^\circ\mathrm{C}$. Les activités humaines ne créent pas l'effet de serre, elles le \emph{renforcent} en augmentant la concentration des GES.
\item \textbf{Vrai.} À masse égale, le méthane possède un pouvoir de réchauffement global environ 25 à 28 fois supérieur à celui du \ce{CO2} sur un horizon de 100 ans, même s'il est émis en quantités moindres et persiste moins longtemps dans l'atmosphère.
\item \textbf{Faux.} La déforestation renforce l'effet de serre par un \emph{double} mécanisme : la combustion ou la décomposition du bois libère du \ce{CO2}, \textbf{et} la disparition des arbres supprime un puits de carbone (la photosynthèse qui absorbait le \ce{CO2} atmosphérique n'a plus lieu). La réduire aux feux de brousse est donc incomplet.
\item \textbf{Faux.} L'hydroélectricité n'implique aucune combustion : ses émissions directes de \ce{CO2} par kilowattheure sont négligeables, alors qu'une centrale thermique au gaz émet environ \SI{0,45}{kg} de \ce{CO2} par kWh. C'est pourquoi le développement des barrages (Song Loulou, Memve'ele, Nachtigal) limite les émissions nationales.
\end{enumerate}
\end{corrige}

\begin{corrige}[Exercice 3 — Définitions]
\begin{enumerate}
\item \textbf{Effet de serre} : phénomène naturel par lequel certains gaz de l'atmosphère absorbent le rayonnement infrarouge émis par la surface terrestre chauffée par le Soleil, puis réémettent de la chaleur vers le sol, maintenant la température moyenne du globe à environ $+15\ ^\circ\mathrm{C}$.
\item \textbf{Gaz à effet de serre (GES)} : gaz atmosphérique capable d'absorber le rayonnement infrarouge émis par la Terre et de le réémettre dans toutes les directions, contribuant ainsi au réchauffement de la basse atmosphère. Exemples : vapeur d'eau, \ce{CO2}, méthane \ce{CH4}, protoxyde d'azote \ce{N2O}.
\item \textbf{Atténuation} : ensemble des actions visant à réduire les causes du changement climatique, c'est-à-dire à diminuer les émissions de GES (énergies renouvelables, économies d'énergie) ou à augmenter les puits de carbone (reboisement).
\item \textbf{Adaptation} : ensemble des mesures visant à réduire la vulnérabilité des sociétés et des écosystèmes face aux effets déjà observés ou attendus du changement climatique (variétés de cultures résistantes à la sécheresse, aménagement des zones côtières, systèmes d'alerte précoce).
\end{enumerate}
\end{corrige}

\begin{corrige}[Exercice 4 — Calcul direct d'émissions]
\begin{enumerate}
\item Émissions journalières :
\[ m = 4\ \mathrm{L} \times 2{,}3\ \mathrm{kg/L} = 9{,}2\ \mathrm{kg\ de\ \ce{CO2}\ par\ jour.} \]
\item Émissions mensuelles (26 jours travaillés) :
\[ m_{\text{mois}} = 9{,}2 \times 26 = 239{,}2\ \mathrm{kg\ de\ \ce{CO2}.} \]
\item Émissions annuelles (312 jours travaillés) :
\[ m_{\text{an}} = 9{,}2 \times 312 = 2\,870{,}4\ \mathrm{kg\ de\ \ce{CO2}.} \]
Nombre d'arbres nécessaires :
\[ N = \frac{2\,870{,}4}{10} = 287{,}04. \]
Comme un arbre entier est nécessaire pour absorber le reliquat, il faut arrondir à l'entier supérieur : il faudrait environ \textbf{288 jeunes arbres} pour compenser les émissions annuelles d'une seule moto-taxi.
\end{enumerate}
\textbf{Point de vigilance :} ne pas arrondir 287,04 à 287, car 287 arbres n'absorberaient que \SI{2870}{kg}, insuffisant ; en compensation carbone, on arrondit toujours à l'entier supérieur.
\end{corrige}

\begin{corrige}[Exercice 5 — Relier activité et émission]
\begin{center}
\begin{tabularx}{\linewidth}{|X|l|l|}
\hline
\rowcolor{popblueL} Activité humaine & GES principal & Secteur \\
\hline Riziculture inondée (Yagoua) & Méthane \ce{CH4} (fermentation anaérobie dans l'eau des rizières) & Agriculture \\
\hline Décharge sauvage d'ordures (Yaoundé) & Méthane \ce{CH4} (décomposition des déchets organiques sans oxygène) & Déchets \\
\hline Véhicules à essence et gasoil & Dioxyde de carbone \ce{CO2} (combustion des carburants fossiles) & Transport \\
\hline Combustion du charbon de bois & Dioxyde de carbone \ce{CO2} (et \ce{CH4} si combustion incomplète) & Énergie \\
\hline Élevage de bovins (Adamaoua) & Méthane \ce{CH4} (fermentation entérique lors de la digestion) & Agriculture / élevage \\
\hline
\end{tabularx}
\end{center}
\textbf{Point de vigilance :} ne pas attribuer systématiquement le \ce{CO2} à toutes les activités : le méthane domine dans la riziculture, les décharges et l'élevage.
\end{corrige}

\begin{corrige}[Exercice 6 — Schéma du bilan radiatif]
Le schéma attendu fait figurer :
\begin{itemize}
\item une flèche jaune descendante du Soleil vers la Terre : \emph{rayonnement solaire incident} (visible) ;
\item une flèche jaune ascendante renvoyée par les nuages et les surfaces claires : \emph{rayonnement réfléchi} ;
\item une flèche rouge ascendante partant du sol : \emph{rayonnement infrarouge émis par le sol} ;
\item une flèche rouge redescendant de l'atmosphère vers le sol : \emph{infrarouge absorbé puis réémis par les GES} ;
\item une flèche rouge quittant l'atmosphère : \emph{infrarouge s'échappant vers l'espace}.
\end{itemize}
\begin{popfigure}
\begin{center}
\begin{tikzpicture}[scale=1]
\draw[fill=popblueL,draw=none] (0,2.6) rectangle (9,4.2);
\node[popdark] at (4.5,3.9) {\small Atmosphère (gaz à effet de serre)};
\draw[fill=popgreenL,draw=popdark] (0,0) rectangle (9,0.6);
\node[popdark] at (4.5,0.3) {\small Surface terrestre (sol et océans)};
\draw[fill=popgold] (0.6,5.6) circle (0.45);
\node[popdark] at (0.6,6.3) {\small Soleil};
\draw[-{Stealth},very thick,popgold] (1.1,5.2) -- (2.6,0.7) node[midway,left,popdark]{\small rayonnement solaire incident};
\draw[-{Stealth},very thick,popgold] (3.4,0.7) -- (4.6,4.4) node[midway,right,popdark]{\small rayonnement réfléchi};
\draw[-{Stealth},very thick,poporange] (5.4,0.7) -- (5.4,2.5) node[midway,right,popdark]{\small infrarouge émis par le sol};
\draw[-{Stealth},very thick,poporange] (6.2,2.5) -- (6.2,0.8) node[midway,right,popdark]{\small infrarouge réémis vers le sol par les GES};
\draw[-{Stealth},very thick,poporange] (7.6,0.7) -- (8.6,4.6) node[midway,right,popdark]{\small infrarouge s'échappant vers l'espace};
\end{tikzpicture}
\end{center}
\legende{Schéma simplifié du bilan radiatif de la Terre : le sol absorbe le rayonnement solaire, émet de l'infrarouge dont une partie est piégée et réémise vers la surface par les gaz à effet de serre.}
\end{popfigure}
\textbf{Explication :} si la concentration en GES augmente, l'atmosphère absorbe une fraction plus grande du rayonnement infrarouge émis par le sol et en réémet davantage vers la surface. La part d'énergie qui s'échappe vers l'espace diminue : le bilan radiatif devient excédentaire et la température moyenne du globe s'élève jusqu'à un nouvel équilibre.
\end{corrige}

\begin{corrige}[Exercice 7 — Bilan carbone d'un déplacement]
\begin{enumerate}
\item Distance quotidienne (aller-retour) : $3 \times 2 = 6$ km. Émissions annuelles :
\[ m = 6\ \mathrm{km} \times 0{,}080\ \mathrm{kg/km} \times 180\ \mathrm{jours} = 86{,}4\ \mathrm{kg\ de\ \ce{CO2}\ par\ an.} \]
\item La marche à pied émet $0$ kg de \ce{CO2} : l'économie est donc de 86,4 kg par an. Autres avantages : amélioration de la santé (activité physique régulière, prévention des maladies cardiovasculaires) et économie d'argent (aucun frais de transport).
\item Le \textbf{vélo} émet 0 kg de \ce{CO2} à l'usage. On peut aussi proposer le \textbf{transport collectif} (bus) : un bus émettant par exemple 0,8 kg de \ce{CO2} par km avec 20 passagers ramène l'émission à $0{,}8/20 = 0{,}04$ kg par km et par passager, soit une division par deux ; avec un bus bien rempli (40 passagers), on atteint 0,02 kg/km, c'est-à-dire une division par quatre du facteur 0,080. Le covoiturage à quatre dans un véhicule dont le facteur est de 0,16 kg/km donne de même $0{,}16/4 = 0{,}04$ kg/km/passager.
\end{enumerate}
\textbf{Point de vigilance :} bien compter l'aller \emph{et} le retour (6 km et non 3 km par jour) ; c'est l'erreur la plus fréquente.
\end{corrige}

\begin{corrige}[Exercice 8 — Lecture de données climatiques]
\begin{enumerate}
\item Augmentation entre 1960 et 2020 :
\[ \Delta C = 414 - 317 = 97\ \mathrm{ppm}. \]
En pourcentage :
\[ \frac{97}{317} \times 100 \approx 30{,}6\ \%. \]
La concentration en \ce{CO2} a augmenté d'environ 31 \% en 60 ans.
\item Entre 1960 et 1980 (deux décennies) :
\[ \frac{339 - 317}{2} = 11\ \mathrm{ppm\ par\ décennie.} \]
Entre 2000 et 2020 :
\[ \frac{414 - 369}{2} = 22{,}5\ \mathrm{ppm\ par\ décennie.} \]
\textbf{Constat :} la vitesse d'augmentation a doublé entre les deux périodes : l'accroissement de la concentration en \ce{CO2} s'\emph{accélère}.
\item Deux activités responsables : la combustion des énergies fossiles (charbon, pétrole, gaz) pour l'industrie, le transport et l'électricité ; la déforestation, qui libère le carbone stocké dans la biomasse et réduit les puits de carbone.
\end{enumerate}
\end{corrige}

\begin{corrige}[Exercice 9 — Le bilan radiatif de la Terre (type Probatoire)]
\begin{enumerate}
\item Les gaz à effet de serre (vapeur d'eau, \ce{CO2}, \ce{CH4}, \ce{N2O}) laissent passer l'essentiel du rayonnement solaire visible, qui chauffe le sol. Le sol chauffé émet un rayonnement infrarouge que ces gaz absorbent en grande partie, puis réémettent dans toutes les directions, notamment vers la surface. Ce « piégeage » partiel de la chaleur maintient la température moyenne du globe à environ $+15\ ^\circ\mathrm{C}$ au lieu de $-18\ ^\circ\mathrm{C}$.
\item Schéma attendu : quatre flux fléchés et légendés — (1) 100 unités de rayonnement solaire incident au sommet de l'atmosphère ; (2) 30 unités réfléchies vers l'espace ; (3) 70 unités absorbées par le sol et l'atmosphère ; (4) rayonnement infrarouge émis par le sol, absorbé puis réémis par les GES, une partie s'échappant vers l'espace. (Reprendre le schéma du corrigé de l'exercice 6 en y portant les valeurs 100 / 30 / 70.)
\item Exemples : la combustion des carburants fossiles dans les transports et les centrales thermiques émet principalement du \ce{CO2} ; l'élevage de bovins et la riziculture inondée émettent principalement du méthane \ce{CH4} (on accepte aussi : déforestation $\to$ \ce{CO2} ; engrais azotés $\to$ \ce{N2O}).
\item L'effet de serre existe naturellement depuis que la Terre possède une atmosphère contenant des GES : il est indispensable à la vie. Les activités humaines ne le créent donc pas ; elles \emph{augmentent la concentration} des GES, ce qui accroît la fraction du rayonnement infrarouge piégée : on parle de \textbf{renforcement} de l'effet de serre, responsable du réchauffement climatique actuel.
\end{enumerate}
\textbf{Barème indicatif (sur 10) :} explication du rôle des GES : 3 pts ; schéma avec les quatre flux correctement fléchés et légendés : 3 pts ; deux activités avec le bon gaz : 2 pts ; justification du terme « renforcement » : 2 pts.

\textbf{Points de vigilance :} citer explicitement le rayonnement infrarouge émis par le sol (et non « la chaleur » en général) ; nommer au moins deux GES ; sur le schéma, chaque flèche doit porter une légende ; pour la question 4, la justification attendue est l'existence naturelle et indispensable de l'effet de serre avant toute influence humaine.
\end{corrige}

\begin{corrige}[Exercice 10 — Facture d'électricité et émissions d'un ménage (type Probatoire)]
\begin{enumerate}
\item Part thermique de la consommation mensuelle :
\[ E_{\text{th}} = 300 \times \frac{30}{100} = 90\ \mathrm{kWh\ par\ mois.} \]
\item Émissions mensuelles :
\[ m_{\text{mois}} = 90 \times 0{,}45 = 40{,}5\ \mathrm{kg\ de\ \ce{CO2}.} \]
Émissions annuelles :
\[ m_{\text{an}} = 40{,}5 \times 12 = 486\ \mathrm{kg\ de\ \ce{CO2}.} \]
\item Réduction visée de 20 \% :
\[ m_{\text{éco}} = 40{,}5 \times \frac{20}{100} = 8{,}1\ \mathrm{kg\ de\ \ce{CO2}\ à\ économiser\ chaque\ mois.} \]
Quantité d'électricité thermique correspondante :
\[ E_{\text{éco}} = \frac{8{,}1}{0{,}45} = 18\ \mathrm{kWh\ par\ mois.} \]
(Autre raisonnement équivalent : réduire la consommation totale de 20 \%, soit $300 \times 0{,}20 = 60$ kWh, dont 30 \% thermique : $60 \times 0{,}30 = 18$ kWh.)
\item Deux gestes concrets :
\begin{itemize}
\item Remplacer les ampoules à incandescence par des lampes LED : à éclairement égal, une LED consomme environ 5 à 8 fois moins d'énergie électrique, ce qui réduit directement les kWh facturés ;
\item Éteindre les appareils en veille et débrancher les chargeurs : les veilles consomment un courant résiduel permanent ; les supprimer diminue la consommation de base du ménage.
\end{itemize}
(On accepte aussi : utiliser un chauffe-eau solaire, cuire avec un couvercle, régler le climatiseur à $24\ ^\circ\mathrm{C}$, etc., à condition d'expliquer le mécanisme d'économie.)
\item Les centrales hydroélectriques produisent de l'électricité sans combustion : leurs émissions directes de \ce{CO2} sont négligeables. Chaque kilowattheure produit par Memve'ele ou Nachtigal se substitue à un kilowattheure thermique qui aurait émis environ 0,45 kg de \ce{CO2}. En augmentant la part de l'hydroélectricité dans le mix national, on réduit donc le facteur d'émission moyen de l'électricité camerounaise et les émissions nationales.
\end{enumerate}
\textbf{Barème indicatif (sur 10) :} question 1 : 1 pt ; question 2 : 2 pts ; question 3 : 2 pts ; question 4 : 2 pts (1 pt par geste avec mécanisme) ; question 5 : 3 pts.

\textbf{Points de vigilance :} appliquer le pourcentage de 30 \% \emph{avant} le facteur d'émission (ne pas multiplier 300 kWh par 0,45) ; conserver les unités à chaque étape ; pour la question 4, un geste sans mécanisme d'économie expliqué ne vaut que la moitié des points ; pour la question 5, la justification attendue repose sur la \emph{substitution} d'une source émettrice par une source non émettrice.
\end{corrige}

\begin{corrige}[Exercice 11 — Émissions par secteur et plaidoyer pour l'agroforesterie (type Probatoire)]
\begin{enumerate}
\item Évolution des parts entre 1990 et 2020 :
\begin{itemize}
\item Déforestation : $45 - 55 = -10$ points (baisse) ;
\item Énergie et transport : $28 - 15 = +13$ points ;
\item Agriculture et élevage : $17 - 20 = -3$ points ;
\item Déchets et autres : $10 - 10 = 0$ point.
\end{itemize}
Le secteur dont la part relative a le plus augmenté est \textbf{l'énergie et le transport}, avec une hausse de \textbf{13 points de pourcentage}.
\item Politique sectorielle adaptée (deux mesures concrètes) :
\begin{itemize}
\item développer le transport collectif urbain (bus à haut niveau de service, réseau structuré à Douala et Yaoundé) et le covoiturage, pour réduire le nombre de véhicules individuels en circulation ;
\item promouvoir les énergies renouvelables et l'efficacité énergétique : extension de l'hydroélectricité et du solaire photovoltaïque, normes d'importation limitant l'âge et la consommation des véhicules d'occasion, campagnes d'éco-conduite.
\end{itemize}
\item Double rôle de la forêt : la forêt est un \textbf{puits de carbone} lorsque, par photosynthèse, les arbres absorbent le \ce{CO2} atmosphérique et stockent le carbone dans leur biomasse (troncs, racines) et dans les sols. Elle devient une \textbf{source de carbone} lorsqu'elle est défrichée ou brûlée : la combustion et la décomposition de la biomasse relâchent le carbone stocké sous forme de \ce{CO2}, et le puits disparaît en même temps. C'est ce double effet qui explique que la déforestation reste le premier secteur émetteur au Cameroun.
\item \textbf{Exemple d'argumentation (environ 150 mots) :}

« Monsieur le Maire, notre commune subit déjà les effets du changement climatique : pluies irrégulières, sécheresses prolongées, rendements en baisse. L'agroforesterie, qui associe arbres et cultures sur les mêmes parcelles, offre une réponse concrète et peu coûteuse. D'abord, les arbres plantés absorbent le dioxyde de carbone de l'atmosphère et stockent le carbone dans leur bois et dans le sol : chaque hectare d'agroforesterie devient un puits de carbone qui contribue à la lutte contre le réchauffement et peut ouvrir l'accès de la commune aux financements climatiques. Ensuite, les arbres protègent les cultures : leur ombrage limite l'évaporation pendant les saisons sèches, leurs racines retiennent l'eau et le sol, leurs feuilles enrichissent la terre en matière organique. Les rendements deviennent ainsi plus stables face aux aléas climatiques, ce qui sécurise les revenus des producteurs. Intégrer l'agroforesterie au plan communal de développement, c'est investir à la fois dans le climat, la fertilité des sols et la résilience de nos populations. »
\end{enumerate}
\textbf{Barème indicatif (sur 10) :} question 1 : 2 pts (identification 1 pt, calcul 1 pt) ; question 2 : 2 pts ; question 3 : 2 pts ; question 4 : 4 pts (deux arguments développés : 2 pts ; structure argumentative et registre adapté au destinataire : 1 pt ; respect de la longueur et clarté : 1 pt).

\textbf{Points de vigilance :} distinguer « points de pourcentage » (différence simple : 13 points) et « pourcentage relatif » ; pour la question 3, exiger les deux termes \emph{puits} et \emph{source} avec le mécanisme (photosynthèse / combustion-décomposition) ; pour le plaidoyer, vérifier que les deux arguments imposés (stockage du carbone, résilience des cultures) sont explicitement présents et que le texte s'adresse bien au maire avec un registre de persuasion.
\end{corrige}

\newpage

\coursec{Fiche bilan}

\coursec{Lutte contre l'effet de serre}

\courssub{Carte mentale de la leçon}
\begin{popfigure}
\begin{tikzpicture}[
    mindmap,
    grow cyclic,
    every node/.style={concept, circular drop shadow, execute at begin node=\hskip0pt},
    root concept/.append style={
        concept color=popdark, font=\bfseries\large, text width=3.5cm, align=center
    },
    level 1 concept/.append style={
        level distance=5.2cm, sibling angle=360/7,
        font=\bfseries\normalsize, text width=2.8cm, align=center
    },
    level 2 concept/.append style={
        level distance=3.2cm, font=\small, text width=2.5cm, align=center
    }
]
\node[root concept, align=center]{Effet de serre \&\\ R\'echauffement climatique}
    child[concept color=popblue] { node[align=center]{1. Effet de serre\\ naturel}
        child { node {Principe : rayonnement SW $\to$ LW} }
        child { node {GES : \ce{CO2}, \ce{CH4}, \ce{H2O}, \ce{N2O}} }
        child { node {R\^ole : maintien T° moy. +15°C} }
    }
    child[concept color=poporange] { node[align=center]{2. Renforcement\\ anthropique}
        child { node {\'Energie : combustion fossiles} }
        child { node {Transport : routier, a\'erien} }
        child { node {Agriculture : \'elevage, riziculture, engrais} }
        child { node {D\'eforestation : perte puits C} }
        child { node {Industrie : ciment, chimie, froid} }
    }
    child[concept color=popgreen] { node[align=center]{3. Cons\'equences\\ du r\'echauffement}
        child { node {Hausse T° globale (+1,1°C depuis 1850)} }
        child { node {Mont\'ee niveau mers (dilatation + fonte)} }
        child { node {\'Ev\'enements extr\^emes (s\'echeresses, inondations)} }
        child { node {Biodiversit\'e : migration, extinction} }
        child { node {Sant\'e : vecteurs, canicules} }
        child { node {Agriculture : rendements, sols} }
    }
    child[concept color=poppurple] { node[align=center]{4. Att\'enuation\\ (r\'eduire causes)}
        child { node {\'Energies renouvelables (solaire, hydro, \'eolien)} }
        child { node {Efficacit\'e \'energ\'etique (b\^atiment, industrie)} }
        child { node {Mobilit\'e douce, transport en commun, VE} }
        child { node {Agriculture durable, agroforesterie} }
        child { node {S\'equestration C : reboisement, sols} }
        child { node {\'Economie circulaire, sobri\'et\'e} }
    }
    child[concept color=poppink] { node[align=center]{5. Adaptation\\ (g\'erer effets)}
        child { node {Gestion eau : barrages, irrigation goutte-\`a-goutte} }
        child { node {Vari\'et\'es r\'esistantes (s\'echeresse, chaleur)} }
        child { node {Protection c\^otes : digues, mangroves} }
        child { node {Syst\`emes d'alerte pr\'ecoce m\'et\'eo} }
        child { node {Sant\'e : surveillance maladies vectorielles} }
        child { node {Urbanisme r\'esilient (infiltration, \^ilots frais)} }
    }
    child[concept color=popgold] { node[align=center]{6. Mod\'elisation\\ \& indicateurs}
        child { node {Bilan radiatif : for\c{c}age (W/m$^2$)} }
        child { node {Sc\'enarios GIEC (SSP1-1.9 \`a SSP5-8.5)} }
        child { node {Empreinte carbone (t eq \ce{CO2}/an/hab)} }
        child { node {Budget carbone restant} }
    }
    child[concept color=popgreen!70!popblue] { node[align=center]{7. Engagement\\ citoyen}
        child { node {Gestes quotidiens (tri, T° chauffage, v\'elo)} }
        child { node {Consommer local, de saison, moins de viande} }
        child { node {Isoler son logement, appareils classe A} }
        child { node {Sensibiliser, voter, soutenir ONG} }
        child { node {Compenser \'emissions (projets certifi\'es)} }
    };
\end{tikzpicture}
\legende{Carte mentale synth\'etique de la leçon NCT12 : 7 branches structurant les savoirs, savoir-faire et savoir-\^etre du module \'Education \`a l'Environnement et au D\'eveloppement Durable.}
\end{popfigure}

\courssub{Bilan structuré}

\begin{definition}[D\'efinitions cl\'es \`a conna\^itre par c\oe ur]
\begin{itemize}
    \item \cle{Effet de serre naturel} : Ph\'enom\`ene physique par lequel l'atmosph\`ere (gr\^ace aux GES) absorbe une partie du rayonnement infrarouge \'emis par la Terre et le r\'e\'emet vers le sol, maintenant une temp\'erature moyenne de +15°C (au lieu de -18°C sans atmosph\`ere).
    \item \cle{Gaz \`a effet de serre (GES)} : Gaz traces (\ce{CO2}, \ce{CH4}, \ce{N2O}, \ce{H2O}, HFC, PFC, \ce{SF6}) transparents au rayonnement solaire (ondes courtes) mais absorbants dans l'infrarouge (ondes longues).
    \item \cle{Renforcement de l'effet de serre} : Augmentation de la concentration atmosph\'erique en GES d'origine anthropique, pi\'egeant davantage de chaleur et causant le r\'echauffement climatique.
    \item \cle{R\'echauffement climatique} : Augmentation durable de la temp\'erature moyenne de l'atmosph\`ere et des oc\'eans (environ +1,1°C depuis l'ère préindustrielle).
    \item \cle{Att\'enuation} : Actions visant à r\'eduire les \'emissions de GES ou à augmenter les puits de carbone (agir sur la cause).
    \item \cle{Adaptation} : Ajustements dans les syst\`emes naturels ou humains face aux effets actuels ou attendus du changement climatique (agir sur la cons\'equence).
    \item \cle{Empreinte carbone} : Quantit\'e totale de GES \'emise directement ou indirectement par une activit\'e, un produit, une personne ou un territoire, exprim\'ee en \'equivalent \ce{CO2} (t eq \ce{CO2}).
\end{itemize}
\end{definition}

\begin{propriete}[Formules, grandeurs et rep\`eres num\'eriques (Probatoire)]
\begin{center}
\begin{tabularx}{\linewidth}{|l|X|}\hline
\rowcolor{popblueL}
\textbf{Notion} & \textbf{\'El\'ements \`a ma\^itriser} \\ \hline
Bilan radiatif simplifi\'e & Rayonnement solaire entrant $\approx 340\ \text{W/m}^2$ ; Alb\'edo Terre $\approx 0,3$ ; Rayonnement IR \'emis par la surface $\approx 390\ \text{W/m}^2$ ; Effet de serre $\approx 150\ \text{W/m}^2$ pi\'eg\'es. \\ \hline
For\c{c}age radiatif (GIEC AR6, 2019 vs 1750) & \ce{CO2} : +2,16 W/m$^2$ ; \ce{CH4} : +0,54 W/m$^2$ ; \ce{N2O} : +0,21 W/m$^2$ ; Total anthropique net : +2,72 W/m$^2$. \\ \hline
Concentrations actuelles (2023) & \ce{CO2} $\approx$ 420 ppm ; \ce{CH4} $\approx$ 1920 ppb ; \ce{N2O} $\approx$ 335 ppb. \\ \hline
Pouvoir de r\'echauffement global (PRG-100, AR6) & \ce{CO2} = 1 (r\'ef\'erence) ; \ce{CH4} $\approx$ 27,9--29,8 ; \ce{N2O} $\approx$ 273 ; \ce{HFC-134a} $\approx$ 1300 ; \ce{SF6} $\approx$ 25200. \\ \hline
Objectif Accord de Paris (2015) & Limiter le r\'echauffement \textbf{bien en-dessous de 2°C} et \textbf{de pr\'ef\'erence \`a 1,5°C} vs \'ere préindustrielle. \\ \hline
Budget carbone restant (pour 1,5°C, 50\% chance) & $\approx$ 500 Gt \ce{CO2} d\`es 2020 ($\sim$ 12--13 ans aux \'emissions actuelles $\approx$ 40 Gt/an). \\ \hline
Secteurs \'emetteurs mondiaux (GIEC AR6) & \'Energie (34\%) ; Industrie (24\%) ; AFAT (Agriculture, For\^et, Autre Usage des Terres) (22\%) ; Transport (15\%) ; B\^atiments (6\%). \\ \hline
\end{tabularx}
\end{center}
\end{propriete}

\begin{methode}[M\'ethodes express (Savoir-faire Probatoire)]
\begin{enumerate}
    \item \textbf{Expliquer le m\'ecanisme (sch\'ema) :} Dessiner Terre + Atmosph\`ere. Fl\`eches : Soleil $\to$ SW (traverse) $\to$ Sol (absorbe, chauffe) $\to$ IR/LW (\'emet) $\to$ GES (absorbent/r\'e\'emettent vers bas) $\to$ R\'echauffement surface. \emph{L\'egender : "Rayonnement solaire (ondes courtes)", "Rayonnement infrarouge (ondes longues)", "Absorption par GES", "R\'e\'emission vers le sol"}.
    \item \textbf{Relier activit\'e $\to$ GES :} Identifier le secteur $\to$ Proc\'ed\'e physique/chimique $\to$ Mol\'ecule \'emise. Ex: \textit{Cimenterie} $\to$ Calcinage \ce{CaCO3 -> CaO + CO2 ^} ; \textit{Riziculture inond\'ee} $\to$ Ana\'erobiose $\to$ M\'ethanogen\`ese $\to$ \ce{CH4 ^}.
    \item \textbf{Proposer mesures (Att\'enuation/Adaptation) :} Classer par levier : \textbf{\'Eviter} (sobri\'et\'e), \textbf{R\'eduire} (efficacit\'e, renouvelables), \textbf{Compenser/S\'equestrer} (reboisement, CSC). Pour l'adaptation : \textbf{Prot\'eger}, \textbf{Accommoder}, \textbf{Replier}.
    \item \textbf{Analyser un document (courbe de Keeling, carte T°, graphique \'emissions) :} 1. Nature/Source. 2. Variables/Unit\'es. 3. Tendances (hausse, saisonnalit\'e). 4. Interpr\'etation (lien cause-effet, r\^ole puits). 5. Conclusion argument\'ee.
\end{enumerate}
\end{methode}

\courssub{Checklist avant le Probatoire}
\begin{aretenir}[Checklist avant le Probatoire]
$\square$ Je sais \textbf{d\'efinir} l'effet de serre naturel et citer au moins 3 GES majeurs avec leur formule brute (\ce{CO2}, \ce{CH4}, \ce{N2O}, \ce{H2O}).
$\square$ Je sais \textbf{sch\'ematiser et l\'egender} le m\'ecanisme complet : rayonnement SW entrant, absorption sol, \'emission LW, absorption/r\'e\'emission par GES, r\'echauffement surface.
$\square$ Je connais les \textbf{principales sources anthropiques} par secteur (\'Energie, Transport, Agriculture/\'Elevage, D\'eforestation/Changement d'affectation des sols, Industrie, D\'echets) et le GES \'emis pour chacune.
$\square$ Je sais \textbf{diff\'erencier} att\'enuation (agir sur les causes) et adaptation (agir sur les cons\'equences) et donner 2 exemples concrets \textbf{au Cameroun} pour chacune (ex: Barrage de Lom Pangar / Reboisement Mont Cameroun vs Vari\'et\'es ma\"is tol\'erant s\'echeresse / Digues Wouri).
$\square$ Je sais \textbf{citer 4 cons\'equences majeures} du r\'echauffement (T° globale, niveau mers, \'ev\'enements extr\^emes, biodiversit\'e/sant\'e/agriculture) et en expliquer le m\'ecanisme physique.
$\square$ Je ma\^itrises les \textbf{ordres de grandeur} cl\'es : T° actuelle (+1,1°C), Objectif Paris (1,5--2°C), \ce{CO2} atm (420 ppm), Budget carbone restant ($\sim$500 Gt), PRG du \ce{CH4} ($\sim$28--30).
$\square$ Je sais \textbf{calculer/estimer une empreinte carbone} simple (ex: trajet voiture, conso \'electricit\'e, repas) en utilisant des facteurs d'\'emission standards.
$\square$ Je connais le \textbf{r\^ole des puits de carbone} (Oc\'ean $\approx$ 25--30\%, V\'eg\'etation/Sols $\approx$ 25--30\%, Atmosph\`ere $\approx$ 45--50\% des \'emissions annuelles) et le risque de saturation/inversion.
$\square$ Je sais \textbf{analyser un document scientifique} (courbe de Keeling, sc\'enarios GIEC SSP, carte anomalies thermiques) : d\'ecrire, expliquer, conclure.
$\square$ Je peux \textbf{proposer 3 gestes citoyens quotidiens} \`a fort impact (Mobilit\'e douce/Transport en commun, Isolation/\'Economie \'energie, Alimentation locale/v\'eg\'etale/R\'eduction gaspillage) et justifier leur effet sur les \'emissions.
$\square$ Je comprends le \textbf{principe de responsabilit\'e commune mais diff\'erenci\'ee} (RCPD) et l'enjeu de justice climatique (Pays du Nord historiques vs Pays du Sud vuln\'erables comme le Cameroun).
$\square$ Je sais \textbf{argumenter} sur la n\'ecessit\'e d'une action \textbf{syst\'emique} (politiques publiques, prix du carbone, transition juste) et non seulement individuelle.
\end{aretenir}

\findecours

\end{document}
